\renewcommand{\textbf}[1]{#1}
\chapter{Onde está o laboratório e os seus cientistas?}
\label{capitulo1}

\epigrafe{%
A Pesquisa é a zona para a qual são arrastados humanos e não humanos, onde ao longo das eras foi feito o mais extraordinário dos experimentos coletivos para distinguir, em tempo real, o ``cosmos'' da ``desordem'' sem que ninguém, cientista ou ``estudioso da ciência'', pudesse saber de antemão qual seria a resposta provisória. Talvez, afinal das contas, os estudos de ciência sejam Anticiência.}%
{Bruno Latour, \textit{A ciência em ação: como seguir cientistas e engenheiros sociedade afora}, 2011.}

Em dezembro de 2022, estive no México para participar do congresso da \textit{Society for Social Studies of Science} (4S), onde apresentei pela primeira vez os resultados preliminares desta pesquisa. Além da experiência de conhecer o país e participar de um evento internacional na área de estudos sociais da ciência e tecnologia, viajei com um grupo diverso de pesquisadores e amigos que contribuíram para aprofundar a reflexão sobre o que significa fazer ciências sociais e investigar como a ciência é produzida em outros áreas do conhecimento. Clarissa, uma das companheiras da viagem e do doutorado na Unicamp, investigava um tema muito diferente sobre o trabalho de antropólogas que menstruam em campo. Ainda assim, compartilhávamos uma preocupação: a ideia de ``não ferir os sentimentos estabelecidos'', uma reflexão de Isabelle Stengers,\footnote{Stengers propõe em \textit{Another Science is Possible} uma ``ciência lenta'' (\textit{slow science}) como resistência às pressões de produtividade que degradam a qualidade do pensamento científico \parencite{stengers2018another}. A defesa da lentidão pressupõe condições de autonomia que nem sempre são possíveis. Esta tese documenta como a temporalidade acadêmica pode se tornar condição de vulnerabilidade quando articulada às lógicas temporais corporativas. Enquanto universidades públicas operam em horizontes de décadas, corporações decidem em ciclos de curto prazo segundo lógicas de retorno sobre investimentos. A capa da edição inglesa traz um caracol, imagem que ressoa com o lema zapatista ``lento, pero avanzo'' e que esta tese adota na dedicatória.} ``como estratégia para criar alianças com interlocutores cientistas, engenheiros e áreas afins''. Longe de se configurar como atitude conformista, trata-se justamente do oposto: ``uma postura calcada apenas no confronto e na denúncia das mazelas da tecnociência é infrutífera, pois produz um resultado alienante e perverso: oculta para nós, pessoas que pesquisam desde as ciências humanas, nossa responsabilidade enquanto coabitantes desse mesmo espaço onde estão os cientistas que denunciamos'' \parencite[p. 314]{costa2024manchando}. Isso não significa que as denúncias sejam irrelevantes, mas ``reconhecer nossa responsabilidade como coabitantes do sistema de produção de conhecimento científico é também aprender a construir e sustentar diálogos, por mais difíceis que sejam'' \parencite[p. 314]{costa2024manchando}.\footnote{Costa desenvolve o conceito de \textit{armadilha} como estratégia metodológica e política para que corpos marcados produzam conhecimento crítico dentro de estruturas acadêmicas sem serem devorados por elas. A armadilha opera em duas etapas, primeiro, encanta através da beleza estética; depois, revela que as imagens fractais são feitas de sangue menstrual, vazando questões feministas de forma que não sejam imediatamente rejeitadas. A etnografia também pode operar como armadilha. Seguir os atores sem julgamento prévio, ganhar confiança, descrever o que é dito e feito para depois tornar visíveis tensões e contradições que os próprios atores não nomeiam \parencite{costa2024manchando}.}

Quase três anos depois do congresso no México, em novembro de 2025, a IBM fechou seu laboratório de pesquisa no Brasil, o único da América Latina, operando há 15 anos com mais de 100 pesquisadores. Em dezembro de 2025, encerrou-se o convênio C4AI-IBM. Desde o início desta pesquisa eu me perguntava até quando aquilo duraria. Na época, havia expectativa por parte do C4AI de que se estendesse por dez anos. Acabou em cinco. A confiança com que meus interlocutores falavam sobre a participação da IBM e a solidez dos investimentos exemplifica de um modo brutal a ``economia especulativa de promessas'' que Stengers critica \parencite{stengers2018another}. Eu, no entanto, repetia para mim mesma: é uma empresa, quero dizer, funciona através de lógicas muito distintas das da universidade. 

O encerramento chegou quando eu entrava na fase final de escrita desta tese e obrigou-me a rever algumas questões que atravessam esta pesquisa. No entanto, acabou por ser também um presente etnográfico sobre o registro de uma parceria público-privada controversa, tensa, contraditória e, a seu modo, bem-sucedida na pesquisa em IA em uma universidade pública brasileira. Essa situação me aproximou do estudo de \textcite{Latour1993Aramis} sobre o Aramis, um sistema de transporte que nunca chegou a existir. Latour conduziu sua pesquisa entre dezembro de 1987 e janeiro de 1989 com o intuito de investigar por que o projeto havia sido cancelado em fevereiro de 1987.\footnote{Minha posição temporal distingue-se: iniciei a pesquisa em março de 2022, quando a parceria parecia consolidada; o encerramento em dezembro de 2025 não estava previsto. Registrei processos enquanto ocorriam, capturando tensões que permaneciam latentes porque a rede ainda resistia.} Tanto o Aramis quanto o C4AI são casos de parcerias público-privadas onde atores com interesses diversos mobilizam promessas de longo prazo que se desfazem quando cálculos estratégicos corporativos mudam. Em fevereiro de 1987, a Matra distribuiu uma nota ordenando ``pare tudo''; em agosto de 2025, a IBM comunicou internamente que o convênio não seria renovado. No Aramis, a retirada resultou no cancelamento completo após quase vinte anos de investimentos públicos.\footnote{O que restou do Aramis foi um protótipo no saguão da empresa, uma casca vazia sem trilhos, sem eletricidade, sem passageiros: ``um objeto de museu'' antes mesmo de ter sido um objeto no mundo. Latour pergunta quando uma tecnologia se torna um objeto. ``Nunca'', ele responde, ``exceto quando porções extraídas da instituição são colocadas em exposição dentro de museus técnicos'' \parencite[p. 9]{Latour1993Aramis}. É só no ferro-velho que um objeto técnico finalmente se torna objeto. Antes disso, é sempre uma instituição, uma rede de humanos e não humanos, uma transação em curso. Law radicaliza esse argumento ao propor vocabulário para pensar objetos que nunca alcançam fechamento, investigando ``as vidas incertas e complexas dos objetos em um mundo onde não há fechamento'' \parencite[p. 59]{Law2004After}. Objetos são múltiplos mesmo quando parecem singulares, e essa multiplicidade persiste sem se resolver com estabilização.} No C4AI, a saída da IBM não acabou com o Centro, mas forçou sua reconfiguração como a perda de financiamento corporativo, incertezas sobre projetos em andamento e o esvaziamento das expectativas construídas ao longo de cinco anos. A promessa de dez anos durou cinco. O Centro continua a existir, mas a parceria IBM-C4AI tornou-se objeto do passado; uma relação que deixou rastros em artigos científicos, infraestrutura computacional e promessas não cumpridas. Esta tese começou como uma etnografia de um arranjo sociotécnico aparentemente estável e terminou descrevendo justamente a sua instabilidade e mudanças em um campo volátil e de difícil acesso, e o que isso pode vir a dizer sobre a pesquisa pública em IA no Brasil.

A fotografia da figura~\ref{fig:dgx1-openai} circulou pela internet depois que Musk a publicou no Twitter em 2016, e tornou-se uma das imagens de referência da história recente da inteligência artificial: Jensen Huang (CEO da Nvidia) entrega a Elon Musk (co-fundador da OpenAI em 2015) um servidor de US\$~129.000 com oito GPUs, e Musk assina o chassi da supermáquina. Ao fundo, uma frase na parede, parcialmente legível, anuncia que o destino do mundo depende do uso que se faz do conhecimento. Menos de uma década depois, a arquitetura de treinamento inaugurada por aquela máquina produziria sistemas capazes de gerar texto, imagem, código e voz com fluência suficiente para redesenhar mercados de trabalho, práticas científicas, sistemas educacionais e disputas geopolíticas. Treinar esses sistemas exige acesso a dezenas, centenas ou milhares de GPUs; a infraestrutura que a cena de 2016 inaugura tornou-se condição de entrada para qualquer laboratório que pretenda produzir modelos competitivos, e a busca por mais poder computacional e mais dados funcionou, por anos, como lei não escrita do desenvolvimento de modelos de linguagem. Em janeiro de 2026, o lançamento do \textit{DeepSeek}, modelo chinês que atingiu desempenho comparável aos sistemas americanos de ponta com apenas uma fração dos recursos computacionais e dos dados, abalou esse paradigma e reabriu a questão sobre quais são, afinal, as condições necessárias para fazer pesquisa em IA. O Spira, projeto desenvolvido no laboratório de processamento de linguagem natural do IME-USP antes mesmo de o C4AI existir como instituição, chegou a essa condição por outra via através de uma doação pontual de GPU da própria Nvidia, registrada nos agradecimentos do primeiro artigo do projeto \parencite{Casanova2021DeepLearning}, sem contrato nem vínculo institucional. O C4AI, por sua vez, operou ao longo de cinco anos com um cluster de GPUs financiado pela USP-Fapes-IBM; em dezembro de 2025, dias após o fim da parceria com a IBM, a USP inaugurou o JAIRU, o maior \textit{cluster} (conjunto de GPUs que trabalham em paralelo) de IA da América Latina, com 96 GPUs Nvidia B200 e um investimento de aproximadamente R\$~40 milhões. As possibilidades que se abriram com aquela máquina são inseparáveis dos problemas que trouxeram, a concentração de poder computacional em poucas corporações, a dependência infraestrutural de países periféricos, a extração de trabalho não remunerado de milhões de usuários cujos dados alimentam os modelos, a opacidade dos sistemas e dificuldade de responsabilização quando erram, entre outros problemas. 

%%%%%%%%%%%%%%%%%%%%%%%%%%
% TRANSIÇÃO PARA TIMELINE
%%%%%%%%%%%%%%%%%%%%%%%%%%

\begin{figure}[!htb]
\centering
\includegraphics[width=1.0\textwidth]{figuras/cap.1/Nvidia-entrega-openai.jpg}
\caption[Entrega do primeiro supercomputador DGX-1 à OpenAI (2016)]{Entrega do primeiro supercomputador DGX-1 à OpenAI (2016). Jensen Huang, CEO da Nvidia, registra o momento em que Elon Musk, co-fundador da OpenAI, assina o chassi do sistema no escritório da OpenAI em São Francisco. O DGX-1, avaliado em US\$~129.000, reunia oito GPUs Tesla P100 com capacidade de até 170 teraflops (FP16), tornando-se a primeira infraestrutura de computação dedicada exclusivamente ao treinamento de redes neurais em escala. Ao fundo, inscrição na parede com trecho parcialmente legível: \textit{Man has a large capacity for [...]. Than we think it is few [...]. Value the faculty of knowing [...]. The sill to do it. Knowing. The critical issue is not [...]. We know ... I believe that [...]. The fate of the worrld.} A fotografia foi compartilhada pelo próprio Musk no Twitter, plataforma que ele viria a adquirir em 2022, renomeando-a como X, e circulou amplamente pela internet. Fonte: Twitter/@elonmusk, 2016.}
\label{fig:dgx1-openai}
\end{figure}

Esta etnografia foi feita durante esses acontecimentos. Iniciei o trabalho de campo no C4AI em março de 2022, seis meses antes do lançamento do ChatGPT, quando a inteligência artificial generativa ainda era objeto de especialistas e o debate público sobre seus efeitos mal havia começado. Entre março de 2022 e dezembro de 2025, assisti, dentro e a partir do Centro, à aceleração que converteu modelos de linguagem em infraestrutura cotidiana, ao crescimento das publicações sobre processamento de linguagem natural pelo C4AI, ao surgimento do primeiro modelo de linguagem brasileiro, à aprovação do primeiro marco regulatório global e ao encerramento da parceria que havia dado origem ao campo que eu estudava. O C4AI, fundado quatro anos depois do DGX-1 com recursos públicos brasileiros e co-financiado pela IBM, é um nó nessa rede e herdeiro das possibilidades, limites e todo o tipo de problema que a fotografia de Jensen Huang e Elon Musk evoca. Acompanhar esse Centro por quase quatro anos é também acompanhar uma pergunta que a fotografia enseja: qual é o papel da universidade pública na produção de conhecimento em inteligência artificial quando quem está nas fronteiras da pesquisa em IA são essas mesmas corporações? Esta tese documenta uma tentativa de resposta situada à essa questão: o que o C4AI fez, como se organizou e o que ficou quando a IBM deixou de co-financiar o centro. 

\begin{figure}[!htb]
  \centering
  \includegraphics[width=0.7\textwidth,keepaspectratio]{figuras/cap.1/urdidura-capitulo1-2026-05-12-231748.pdf}
  \caption[Urdidura do capítulo~\ref{capitulo1}]{\textbf{Urdidura do capítulo~\ref{capitulo1}: condições materiais, sanitárias, regulatórias, bibliométricas e etnográficas em que se tramou esta pesquisa.} A urdidura nomeia, em diálogo com \textcite{Law2004After}, o \textit{hinterland} que sustenta o que aparece como presença no texto da tese: cinco fios longitudinais que correm em paralelo e se atravessam ao longo de uma década, em que a faixa etnográfica se torna texto enquanto as outras quatro permanecem como condição de possibilidade. (i)~A faixa \textit{tecnológica} marca a aceleração da infraestrutura computacional para IA entre 2016 e 2026, do DGX-1 entregue à OpenAI à arquitetura Blackwell da NVIDIA, e inclui o lançamento de modelos cujas trajetórias atravessaram minha pesquisa: GPT-3, ChatGPT, GPT-4, LLaMA, Sabiá e Sabiá-3 da Maritaca AI, \texttt{claude.ai}3.5 Sonnet da Anthropic, com o Nobel de Física de 2024 para Hopfield e Hinton inscrevendo o reconhecimento institucional do campo. (ii)~A faixa \textit{sanitária} é a única barra crítica contínua: a pandemia de Covid-19 entre março de 2020 e maio de 2023 é condição de possibilidade do SPIRA (que coletou vozes em covideiro) e do C4AI (fundado em outubro de 2020, em plena pandemia), e atravessa toda a fase inicial do trabalho de campo. (iii)~A faixa \textit{regulatória} inscreve a aprovação ética do CEP-HCFMUSP para o SPIRA em abril de 2020 (que torna possível a coleta de voz no covideiro), as políticas científicas brasileiras de fomento à IA (EBIA 2021, PBIA 2024) e os arranjos regulatórios internacionais que reposicionaram o setor (CHIPS Act, AI Act), até a Portaria CNPq 2.664 de 2026, que minha tese antecipou por escolha metodológica. (iv)~A faixa \textit{bibliométrica} mostra a aceleração das publicações em IA nas ciências humanas brasileiras (SciELO, CAPES) entre 2020 e 2026, com o ápice de 313 publicações do C4AI em 2023 e o declínio bibliométrico associado à saída da IBM em 2025-2026. (v)~A faixa \textit{etnográfica} concentra os eventos que se tornam texto na tese: o acordo-quadro FAPESP-IBM de 2016 que precede a fundação do Centro, meu ingresso no doutorado em 2019, a fundação do C4AI em outubro de 2020, minha aproximação ao campo em outubro de 2021, as 196 semanas de trabalho de campo (março de 2022 a dezembro de 2025), o letramento técnico acumulado em paralelo (MBA, especialização, cursos), o rascunho integral da tese redigido entre junho e dezembro de 2024 sem assistência de inteligência artificial, a colaboração com \texttt{claude.ai} e a \textit{skill} \texttt{juliane-tese} entre janeiro de 2025 e maio de 2026, as três entrevistas com Marcelo Finger, Fábio Cozman e Cláudio Pinhanez, e os dois eventos críticos que abriram a caixa-preta institucional durante a finalização da pesquisa: o fechamento da IBM Research Brasil em novembro de 2025 e o encerramento da parceria C4AI-IBM em dezembro de 2025. Inscrição produzida em colaboração com \texttt{claude.ai}(Anthropic) e hospedada no \texttt{MermaidChart}\footnote{Versão interativa do diagrama disponível em: \url{https://mermaid.ai/d/b27ca78c-b037-42fe-a3eb-939a092e1455}. Acesso em: \today.}, em que cada fio longo do Gantt opera como linha de \textit{hinterland} que o texto principal precisou deixar fora para sustentar seu fluxo narrativo, e cuja explicitação aqui é exigência da posição reflexiva que este capítulo assume. A urdidura é, nesse sentido, par solidário da Figura~\ref{fig:trama-cap1}: onde aquela inscreve a estrutura argumentativa em cima do tear, esta inscreve os fios longos que correm por baixo dela.}
  \label{fig:urdidura-cap1}
\end{figure}

Estar no campo durante esse período implicou um desafio que vai além do habitual esforço de atualização bibliográfica. Nas seções que se seguem, detalho os fundamentos teórico-metodológicos que orientaram essa trajetória, as operações de triangulação que permitiram articular dados heterogêneos, os cortes analíticos estratégicos que transformaram redes infinitas em objetos etnográficos analisáveis, e as quatro lições que o encontro entre teoria e campo produziu como contribuições metodológicas desta tese. A etnografia, como o caracol, avançou devagar, secretando seu rastro. O que esse rastro deixou está escrito (e inscrito) nesta tese. Mas, esse rastro também deixa meus vestígios enquanto pesquisadora aprendendo a caminhar sobre superfícies estranhas, e deixa, sobretudo, os vestígios dos próprios pesquisadores do C4AI construindo redes, mobilizando inscrições e recrutando aliados, negociando com corporações, fazendo, por fim, \textit{tecnociência}.\footnote{Em \textit{Ciência em Ação}, \textcite{Latour2011CienciaEmAcao} mobiliza o termo para nomear um campo em que ciência e tecnologia se produzem como um só processo, sem partição prévia entre fatos e artefatos, descobertas e invenções, conhecimento e instrumentação. O argumento opera contra a separação que isola o trabalho conceitual dos cientistas do trabalho material dos engenheiros, pois, para Latour,  as controvérsias científicas se resolvem pela capacidade de mobilizar redes longas e heterogêneas, nas quais, teoremas, máquinas, documentos, financiamentos e alianças institucionais atuam como aliados que tornam o custo da discordância proibitivo para o oponente. A tecnociência aparece, nesse sentido, como prática coletiva de fechamento de caixas-pretas, ou seja, de estabilização de associações cuja construção desaparece da vista quando o fato passa a circular como inquestionável. \textcite{stengers2018another} mobiliza o mesmo termo em outra direção. Para Stengers, a tecnociência é a forma assumida pelas práticas científicas quando recrutadas pela lógica industrial e pelos imperativos de competitividade, sob o regime do que chama de \textit{fast science}, em que a velocidade da produção e a captura por interesses corporativos reorganizam as condições de possibilidade do trabalho científico, e em que a formação do pesquisador opera como um tipo de anestesia que o torna insensível à recalcitrância dos mundos atravessados pela pesquisa. Stengers renomeia a economia do conhecimento como \textit{economia especulativa de promessas}, regime em que o financiamento se sustenta pela produção de possibilidades orientadas a investidores futuros e em que a avaliação crítica entre pares se enfraquece porque ninguém quer cortar o galho da promessa que sustenta o campo. Latour e Stengers descrevem dimensões diferentes do mesmo fenômeno: o vocabulário latouriano permite mapear como redes tecnocientíficas se estendem e se estabilizam; o diagnóstico de Stengers expõe como essas redes, sob condições contemporâneas, operam regimes específicos de captura, aceleração e despolitização. Esta tese mobiliza os dois registros. O aspecto interessante na proposta de Stengers é que ela consegue articular, ao mesmo tempo, a crítica (não no sentido de simples denúncia) e a construção de alianças, lembrando que sua formação é em química e filosofia da ciência. A trajetória da parceria C4AI-IBM, com a promessa inicial de dez anos reduzida a cinco quando os cálculos corporativos se reorganizaram, é caso empírico em que as duas leituras se sobrepõem: a rede longa que Latour permite descrever é também a economia de promessas que Stengers diagnostica. Estudá-las em ação, antes do fechamento que tentam estabilizar e da retirada que pode desfazê-las, é o gesto que orienta a primeira regra metodológica reproduzida adiante e a orientação desta tese.} É desse rastro que trato agora. Para isso, começo por mobilizar a noção de \textit{actante}.\footnote{O termo é tomado da semiótica de Algirdas Julien Greimas, que o utilizava para descrever qualquer entidade dotada de papel narrativo em um relato, sem prejulgamento sobre sua natureza humana ou não humana, individual ou coletiva, figurativa ou abstrata. Latour transporta essa generalidade para os estudos de ciência e tecnologia precisamente para escapar do antropocentrismo embutido na palavra \enquote{ator}, que tende a pressupor intenção, consciência e personificação. Em \textit{Ciência em Ação}, oferece a formulação mínima: \enquote{chamarei de actante qualquer pessoa e qualquer coisa que seja representada}, observando que não há diferença prática, para fins de análise, entre representar trabalhadores num sindicato e representar hormônios num experimento, pois uns e outros são mudos e dependem de porta-vozes que articulem suas posições \parencite[p. 127, 138-140]{Latour2011CienciaEmAcao}. Em \textit{A Esperança de Pandora}, a definição é refinada em chave pragmática: um actante é definido pela lista de provas a que é submetido e pelas performances que apresenta nessas provas; sua \enquote{competência} ou essência é deduzida só depois, retrospectivamente, a partir do que conseguiu fazer ou impedir de acontecer \parencite[Glossário]{Latour2017Esperanca}. Um microrganismo, por exemplo, antes de receber nome e estabilizar-se como entidade reconhecível, é apenas a soma das vitórias e derrotas que registra ao transformar açúcar em álcool, resistir à fervura, sobreviver a antissépticos. Em \textit{Políticas da Natureza}, a ênfase desloca-se para a relacionalidade: actante é \enquote{qualquer entidade que modifica outra entidade em uma prova} \parencite[Glossário]{latour2004politics}. A agência, nesse registro, nunca é propriedade isolada de uma entidade; é \textit{fait-faire}, fazer-fazer, distribuído entre os mediadores da rede. No artigo sobre o dispositivo de fechamento automático de portas, Latour acrescenta uma distinção que importa para a descrição etnográfica: actantes podem ser figurativos, quando aparecem personificados (um funcionário, um zelador), ou não figurativos, quando operam de modo técnico ou abstrato (uma mola, um quebra-molas, um algoritmo) \parencite{johnson1988mixing}. O que define o actante é, nesse sentido, o que ele faz na cadeia de associações, não o tipo de entidade que parece ser. Esses quatro registros são complementares e operam juntos nesta tese: actante é qualquer entidade que (i) é representada por um porta-voz, (ii) é definida pragmaticamente pelo que faz nas provas a que é submetida, (iii) modifica outras entidades nessas provas, e (iv) pode aparecer sob figura humana ou sob forma técnica sem que isso altere seu estatuto analítico. Convém marcar uma cautela. Tratar humanos e não humanos como actantes não significa afirmar sua equivalência ontológica nem dissolver as diferenças entre seus modos de agir; significa, mais recusar a partição prévia que decidiria, antes da pesquisa, quem age e quem é apenas cenário ou contexto. Os actantes que importam analiticamente nesta tese são, em sua maioria, entidades que transformam, traduzem e tensionam aquilo que atravessam, em contraste com os intermediários, que apenas transportam algo ou alguma coisa sem alteração. Contratos, modelos computacionais, GPUs, vozes gravadas em enfermarias, o vírus SARS-CoV-2 e linguagens formais entram nas redes que descrevo porque deslocam o curso das associações e das ações em que participam. E tornam possíveis fazer coisas que antes era impossível ou que ocorriam de maneira diferente.} Esse conceito sustenta a noção de simetria teórico-metodológica que, em geral, orienta a teoria ator-rede. Considerando a interpretação de um de seus principais formuladores, Bruno Latour, por exemplo, argumenta que não devemos aceitar acriticamente o que os cientistas dizem sobre a Natureza, tampouco o que os cientistas sociais dizem sobre a Sociedade. A sua quarta regra metodológica estabelece que ``uma vez que a resolução de uma controvérsia é a causa da estabilidade da sociedade, não podemos usar a sociedade para explicar como e por que uma controvérsia foi resolvida. Devemos considerar simetricamente os esforços para alistar e controlar recursos humanos e não-humanos'' \parencite[p. 225]{Latour2011CienciaEmAcao}. Isso desloca Natureza e Sociedade da posição de causas explicativas do desfecho das controvérsias para a posição de consequências estabilizadas ao final do processo. Para esta tese, essa regra tem implicação metodológica direta: não posso explicar a pesquisa em inteligência artificial no C4AI recorrendo a fatores sociais externos como se fossem mais estáveis ou mais reais do que os algoritmos, os dados ou os modelos. Em outras palavras, devo observar empiricamente quem ou o que efetivamente \enquote{age} em determinada situação, não esquecendo também de como e para quê no sentido de Haraway e Stengers. 

O desafio desse tipo de posicionamento analítico consiste em se apoiar no arcabouço teórico-metodológico das ciências sociais e, simultaneamente, voltar esses mesmos conceitos, centrais às investigações sociológicas e antropológicas, para a própria análise. Isso exige uma suspensão recorrente das nossas próprias crenças enquanto cientistas sociais, suspensão que precisa ser refeita a cada nova passagem do texto e que se torna condição para o aprendizado, para um aprender antropologicamente no sentido proposto por \textcite{Ingold2019Antropologia}. É nesse ponto que a análise alcança sua maior densidade, porque os actantes mobilizados nesta tese vão de políticas científicas e contratos institucionais a GPUs, algoritmos, vozes humanas gravadas em enfermarias, inscrições tecnocientíficas dos artigos (imagens microscópicas, diagramas, gráficos, equações) e o próprio vírus SARS-CoV-2. Os capítulo s 3 e 4 documentam como esses actantes se associam em cadeias heterogêneas que dão forma a redes sociotécnicas. Isso ocorre tanto em reuniões de \textit{design thinking}, como meus interlocutores chamam as conversas que decidem os rumos das pesquisas e suas condições, quanto nas enfermarias hospitalares onde pacientes com Covid-19 tiveram suas falas gravadas e convertidas em dados computacionais para treinamento e inferência do Spira. A pandemia atravessa esta tese: o C4AI foi inaugurado em outubro de 2020, em plena crise sanitária, e o trabalho de campo começou em março de 2022, ainda sob seus efeitos. Seguir os atores levou-me ao vírus SARS-CoV-2, que se mostrou um actante que reconfigurou as condições de possibilidade tanto das pesquisas em IA que observo quanto da pesquisa que realizo. Esse posicionamento teórico-metodológico é, no sentido que \textcite{Barad2007Meeting} dá ao termo, ético-onto-epistemológico.\footnote{Como adiantado no capítulo ~\ref{capitulo0}, Karen Barad cunha a expressão \enquote{ethico-onto-epistem-ology} para nomear a inseparabilidade entre as questões do que existe, do que se pode conhecer e da responsabilidade que se assume ao conhecer e fazer existir: \enquote{conhecer, pensar, medir, teorizar e observar são práticas materiais de intra-ação no mundo e como parte dele} \parencite[p.~90]{Barad2007Meeting}. A escrita com hífens do termo indica que as três dimensões são inseparáveis em qualquer prática de pesquisa. Para o argumento desta tese, isso significa que tratar simetricamente vírus, GPUs, algoritmos e políticas científicas é, ao mesmo tempo, decisão sobre o que existe, decisão sobre como conhecer e decisão sobre quem responde pelas associações que se estabilizam. O capítulo ~\ref{capitulo4} mobiliza essa orientação no plano quando o \textit{pipeline} (representação visual de como foi feito um sistema de IA) do Spira produz, pelo próprio aparato, o objeto que mede, calcula, treina e infere.} Esse esforço, contudo, não está livre de contradições. Em muitos momentos, contradigo a mim mesma e aos princípios que digo seguir. Pensar dessa maneira exige virar uma chave interpretativa que nem sempre conseguimos acionar. Anos de formação nas ciências sociais nos ensinam a buscar o social nos fenômenos, a desconfiar das explicações técnicas, a revelar os interesses ocultos. Desaprender esse costume acadêmico é um trabalho contínuo e necessariamente incompleto. Manter essa disciplina (não poderia ser um tipo de recursividade?) ao longo desta tese é uma prova de força (e não estou falando da mesma prova de força no sentido latouriano, apesar de também ter a ver com isso). Por isso, quando o leitor encontrar passagens em que recorro apressadamente a fatores sociais para explicar o que deveria ser descrito em sua heterogeneidade, saiba que estou mais ou menos ciente desses deslizes. 

A formulação ético-onto-epistemológica de Barad encontra o seu antecedente operacional, por assim dizer, nas sete regras metodológicas e nos seis princípios que Latour propõe ao final de \textit{Ciência em Ação}. Onde Barad nomeia a inseparabilidade entre o que existe, o que se pode conhecer e a responsabilidade que se assume ao conhecer, Latour elaborou um guia prático para sustentar essa inseparabilidade ao longo de uma pesquisa: o que considerar como objeto, em que momento entrar no campo, como tratar simetricamente humanos e não humanos, o que fazer diante de divisores entre interior e exterior. Reproduzo-os aqui na íntegra porque cada um deles intervém no texto desta tese como decisão tomada a cada novo capítulo, a cada nova entrevista, a cada decisão sobre qual associação seguir e qual deixar fora da análise.

\begin{citacaoabnt}
\textit{Regra 1.} Estudamos a ciência em ação, e não a ciência ou tecnologia pronta; para isso, ou chegamos antes que fatos e máquinas se tenham transformado em caixas-pretas, ou acompanhamos as controvérsias que as reabrem.

\textit{Regra 2.} Para determinar a objetividade ou subjetividade de uma afirmação, a eficiência ou a perfeição de um mecanismo, não devemos procurar por suas qualidades intrínsecas, mas por todas as transformações que ele sofre depois, nas mãos dos outros.

\textit{Regra 3.} Como a solução de uma controvérsia é a causa da representação da Natureza, e não sua consequência, nunca podemos utilizar essa consequência, a Natureza, para explicar como e por que uma controvérsia foi resolvida.

\textit{Regra 4.} Como a resolução de uma controvérsia é a causa da estabilidade da sociedade, não podemos usar a sociedade para explicar como e por que uma controvérsia foi dirimida. Devemos considerar simetricamente os esforços para alistar recursos humanos e não-humanos.

\textit{Regra 5.} Com relação àquilo de que é feita a tecnociência, devemos permanecer tão indecisos quanto os vários atores que seguimos; sempre que se constrói um divisor entre interior e exterior, devemos estudar os dois lados simultaneamente e fazer uma lista (não importa se longa e heterogênea) daqueles que realmente trabalham.

\textit{Regra 6.} Diante da acusação de irracionalidade, não olhamos para que regra da lógica foi infringida nem que estrutura social poderia explicar a distorção, mas sim para o ângulo e a direção do deslocamento do observador, bem como para a extensão da rede que assim está sendo construída.

\textit{Regra 7.} Antes de atribuir qualquer qualidade especial à mente ou ao método das pessoas, examinemos os muitos modos como as inscrições são coligidas, combinadas, interligadas e devolvidas. Só se alguma coisa ficar sem explicação depois do estudo da rede é que deveremos começar a falar em fatores cognitivos \parencite[p. 420-422]{Latour2011CienciaEmAcao}.
\end{citacaoabnt}

Se as regras orientam a prática etnográfica, os seis princípios extraem dela consequências teóricas mais amplas sobre a natureza dos fatos, das máquinas e das redes que os sustentam:

\begin{citacaoabnt}
\textit{Primeiro princípio.} O destino de fatos e máquinas está nas mãos dos consumidores finais; suas qualidades, portanto, são consequência, e não causa, de uma ação coletiva.

\textit{Segundo princípio.} Os cientistas e engenheiros falam em nome de novos aliados que conformaram e alistaram; representantes entre outros representantes, com esses recursos inesperados, fazem o fiel da balança de forças pender em seu favor.

\textit{Terceiro princípio.} Nunca somos postos diante da ciência, da tecnologia e da sociedade, mas sim diante de uma gama de associações mais fracas e mais fortes; portanto, entender o que são fatos e máquinas é o mesmo que entender o que as pessoas são.

\textit{Quarto princípio.} Quanto mais esotérico o conteúdo da ciência e da tecnologia, mais elas se expandem externamente; portanto, 'ciência e tecnologia' é apenas um subconjunto da tecnociência.

\textit{Quinto princípio.} A acusação de irracionalidade é sempre feita por alguém que está construindo uma rede em relação a outra pessoa que atravessa seu caminho; portanto, não há Grande Divisor entre mentes, mas apenas redes maiores ou menores.

\textit{Sexto princípio.} A história da tecnociência é, em grande parte, a história dos recursos espalhados ao longo das redes para acelerar a mobilidade, a fidedignidade, a combinação e a coesão dos traçados que possibilitam a ação a distância \parencite[p. 422-424]{Latour2011CienciaEmAcao}.
\end{citacaoabnt}

Conforme desenvolvia esta pesquisa, orientada por esses princípios e regras, percebi que, ao fazer associações com os pesquisadores do C4AI, tanto eu quanto esta tese passamos a fazer parte das redes do Centro, e tanto eu quanto a tese nos tornamos actantes nessas redes; assim, esta tese tornou-se também uma inscrição tecnocientífica.\footnote{A percepção de que esta tese também fazia redes me foi sugerida por Fernanda Rosa, em conversa após eu apresentar as primeiras impressões desta pesquisa no encontro anual da \textit{Society for Social Studies of Science} (4S), em Cholula, em dezembro de 2022. Fernanda Rosa investiga infraestruturas digitais no Sul Global e aparece também na revisão do capítulo ~\ref{capitulo2}. A formulação dela deslocou a posição em que eu me colocava até então: o que eu vinha descrevendo como acompanhamento dos cientistas e engenheiros do C4AI também era construção de associações que iam ganhando forma na medida em que as registrava. Levei tempo para perceber o alcance metodológico daquela formulação, e foi essa percepção que veio a desenhar o eixo desta tese.} Latour define instrumento como qualquer estrutura que possibilite uma exposição visual de qualquer tipo em um texto científico \parencite[p.~102]{Latour2011CienciaEmAcao}, e inscrição como aquilo que esse instrumento produz: o traçado, o gráfico, a tabela, o registro.\footnote{O conceito de inscrição é desenvolvido em três frentes. Em \textit{Ciência em Ação} \parencite[caps.~1--2]{Latour2011CienciaEmAcao}, Latour analisa o artigo científico como construção sobre camadas de inscrições e descreve o laboratório como fábrica de inscrições; em \textit{A Esperança de Pandora} \parencite[cap.~2]{Latour2017Esperanca}, acompanha uma expedição pedológica à Amazônia e mostra como o solo se transforma sucessivamente em amostras, em cores de um código de Munsell, em números e em um diagrama de relatório, cadeia que ele chama de referência circulante; em \enquote{Visualization and Cognition} \parencite{Latour1986Visualization}, define as inscrições como \textit{móveis imutáveis}, condição da dominação à distância. Para que o conhecimento se acumule em centros de cálculo, as inscrições precisam ser móveis, transportáveis dos lugares distantes para o centro; imutáveis, sem distorção no transporte; e combináveis, rearranjáveis em uma única superfície plana. A construção de fatos opera por uma cascata de simplificações: uma substância complexa se reduz a um traço, que se reduz a um número, que se torna média em uma tabela, até chegar a uma afirmação dura em um artigo. No texto científico, a inscrição traz o referente para dentro da narrativa: o autor aponta para a figura e mostra a prova, em vez de pedir que o leitor acredite em sua palavra.} Esta etnografia mobiliza uma série dessas inscrições por meio de imagens, diagramas, gráficos e do site que apresenta e performa esta tese. E quando digo que a tese em si é uma inscrição tecnocientífica, é porque ela possibilita certos tipos de ações. De alguma maneira, traduzi aquele pedaço do mundo envolvendo o C4AI para o formato de papel, ou melhor, para um arquivo digital que começou como anotações em \texttt{.txt} no bloco de notas do celular, passou por um \texttt{.docx} no Word e no Google Docs, depois para LaTeX no Overleaf, até chegar a um PDF e, talvez, a um volume encadernado ou a um livro impresso.\footnote{Diferentemente de um texto impresso em papel (que também depende de uma complexa cadeia de técnicas e tecnologias, afinal, até a escrita à mão exige algum suporte físico e um conjunto de habilidades aprendidas ao longo do tempo), um arquivo digital só pode ser compreendido por um dispositivo, como um computador, para que possamos ver as letras na tela. Isso pode parecer um detalhe simples, mas envolve diversas mediações técnicas: o computador lê os diferentes formatos de arquivo convertendo os dados em sequências de bits (0s e 1s, desligado/ligado, falso/verdadeiro) e transformando-os em informações compreensíveis, em um processo que passa pelo sistema operacional, pelas extensões de arquivo e pelos aplicativos específicos instalados na máquina. Sem um software de aplicação que funcione como tradutor dos dados brutos contidos nos arquivos e os transforme em algo compreensível para o usuário, os dados de arquivos como o desta tese em PDF permanecem apenas como dados crus, isto é, mera eletricidade; e, sem a ação do software, essa eletricidade não pode ser convertida em uma mensagem ou informação inteligível. Por exemplo, o código binário 01000001, para um processador de texto, corresponde à letra \enquote{A}. Não é impressionante o que a computação tornou possível? Quando consideramos os modelos de linguagem, tudo isso fica ainda mais interessante, porque esses modelos funcionam em uma escala de indeterminação, ao contrário de qualquer outro sistema computacional tradicional, e aliás contra as suas próprias bases fundacionais que opera segundo lógicas de determinação e continuidade. Arrisco sugerir que é justamente essa indeterminação e descontinuidade que torna a inteligência artificial tão instigante para a antropologia. Por mais que se tente controlar tais aspectos por meio dos famosos \textit{prompts} (como ficaram conhecidas as instruções que são feitas à IA generativa), eles não serão totalmente eliminados. O \textit{prompt} expressa, em certo sentido, a ilusão frustrada de que seria possível controlar integralmente esses sistemas. Isso não significa que esses sistemas sejam incontroláveis, mas talvez os \textit{prompts} não sejam a melhor forma de fazê-lo. Eu mesma realizei experimentações e considerei o resultado tão bizarro em tentar impedir que o sistema falhasse, cometesse erros ou reproduzisse suas marcas estilísticas características; pareceu-me que a IA foi assim \enquote{desumanizada} de uma vez por todas.} Esse recorte dizia respeito a algo muito maior do que eu, no qual se desenrolavam acontecimentos que só podiam ser notados a partir de uma certa escala de análise, incluindo não humanos invisíveis ao olhar humano sem a mediação da tecnociência, como o coronavírus. Ao mobilizar inscrições de diversos tipos, a tese transforma eventos dispersos no tempo e no espaço em traçados estáveis, combináveis e transportáveis. À medida que essa rede se expande para além do Centro, associando actantes cada vez mais inesperados e heterogêneos, a tese adquire o status do que Latour denomina inscrição também literária. É essa inscrição que permite à tecnociência agir à distância e converter fenômenos complexos, remotos ou invisíveis em um arquivo digital que pode ser lido, interpretado, analisado, isto é, dominado, medido, controlado, utilizado tanto para apreciação quanto para persuadir outros acerca da realidade de um fato tecnocientífico cujas cadeias de associações ganham força conforme a extensão da rede e o número de elementos que ela consegue associar \parencite{Latour2011CienciaEmAcao}, ou fazer alianças \parencite{stengers2010cosmopolitics1}.

% =====================================================================
% capítulo  1 — Esboço de seção/interlúdio
% Sugestão de inserção: imediatamente após o parágrafo que termina com
% "[...] fazer alianças \parencite{stengers2010cosmopolitics1}." (tese
% como inscrição tecnocientífica), e antes do parágrafo que abre com
% "De certo modo, em sintonia com as críticas de Haraway e Stengers [...]".
%
% Observação técnica: para o build do LaTeX encontrar a figura, copie
%   infranodus/infranodus_cap1_focus.png
% para
%   figuras/cap.1/infranodus_cap1_focus.png
% (mantendo a convenção do capítulo). O caminho abaixo já assume essa
% localização. Se preferir manter a figura em `infranodus/`, basta ajustar
% o argumento de \includegraphics.
% =====================================================================
\section{Etnografias: fazendo o método a partir da experiência}
\label{sec:etnografia-heterodoxa}

\epigrafe{%
Há muitos métodos para o estudo da construção de fatos científicos e artefatos técnicos. A primeira regra metodológica pela qual nos decidimos na Introdução é a mais simples de todas. Não tentaremos analisar os produtos finais; em vez disso, seguiremos os passos de cientistas e engenheiros nos momentos e nos lugares nos quais planejam, desfazem, modificam. Vamos dos produtos finais à produção, de objetos instáveis e frios a objetos instáveis e mais quentes. Estar ali \textbf{antes} que a caixa se fechasse e ficasse preta. Com esse método simples precisamos apenas seguir os melhores de todos os guias, os próprios cientistas, em sua tentativa de fechar uma caixa-preta e abrir outra.}%
{Bruno Latour, \textit{A ciência em ação: como seguir cientistas e engenheiros sociedade afora}, 2011.}

Esta pesquisa fundamenta-se em uma acepção específica de etnografia elaborada por Suely Kofes.\footnote{A acepção de etnografia mobilizada nesta tese é a que Suely Kofes desenvolve em \textit{Dilemas na maçonaria contemporânea: um experimento antropológico} \parencite{Kofes2015Maconaria}. A formulação aparece em duas seções complementares do livro: na parte dois, intitulada \enquote{Etnografia heterodoxa}, e na seção final, \enquote{Pesquisa antropológica como método e/ou como experiência} \parencite[pp.~291--303]{Kofes2015Maconaria}. Nessas seções, Kofes articula o entrelaçamento entre biografia, etnografia e ciências sociais como alternativa às categorizações dicotômicas supostamente universais. Essa acepção dialoga com obra anterior da autora, \textit{Uma trajetória, em narrativas} \parencite{kofes2001trajetoria}, na qual Kofes contrasta sua proposta com outras definições consolidadas de etnografia: a etnografia como totalidade na formulação de Bronislaw Malinowski em \textit{Argonautas do Pacífico Ocidental}, como trabalho de campo na formulação de Martyn Hammersley e Paul Atkinson em \textit{Ethnography: Principles in Practice}, e como escritura na formulação de James Clifford e George Marcus em \textit{Writing Culture: The Poetics and Politics of Ethnography}. Kofes observa que nenhuma dessas definições dá conta isoladamente de explicar o que faz uma etnógrafa.} A proposta de Kofes é pensar a pesquisa antropológica como método e como experiência, formulação que ela desenvolve no capítulo final do livro de 2015 ao fazer dialogar Strathern, Favret-Saada e Turner. Kofes parte de uma desconfiança sobre a verdade do empirismo como método e da consciência dos limites do método: o que está em jogo é a relação entre pesquisa de campo, escrita e teoria, e a posição da etnógrafa nessa relação \parencite[pp.~298, 300]{Kofes2015Maconaria}. Mobilizando Marilyn Strathern, Kofes situa a etnografia como prática reveladora ou \enquote{revelatória}, que convida o leitor a considerar o que não é revelado na evidência do que o é \parencite[p.~297]{Kofes2015Maconaria}. Com Jeanne Favret-Saada, sustenta que a etnógrafa precisa entrar no processo da palavra e se tornar uma falante como outra, posicionar-se no e como discurso, e ao escrever os resultados retornar à situação de enunciação, fazendo desse vai e vem entre a tomada inicial e sua retomada teórica o objeto mesmo da reflexão \parencite[p.~299]{Kofes2015Maconaria}.\footnote{A formulação que Suely Kofes recupera de Jeanne Favret-Saada vem de \textit{Les mots, la mort, les sorts: la sorcellerie dans le Bocage} \parencite[p.~32]{FavretSaada1977}, etnografia da feitiçaria no Bocage francês na qual a autora reconstrói a posição da etnógrafa como interna ao processo de enunciação que descreve. O argumento de Favret-Saada se desdobra em ensaio posterior, \enquote{Être affecté}, publicado originalmente na revista \textit{Gradhiva} e mais tarde recolhido em \textit{Désorceler} \parencite{FavretSaada1990SerAfetado}, no qual a autora formula a categoria do \enquote{ser afetado} como modalidade epistemológica específica da etnografia, distinta da observação participante e da empatia identificatória. A tradução brasileira desse ensaio, feita por Paula Siqueira Lopes e Tânia Stolze Lima, está disponível em \textit{Cadernos de Campo} \parencite{FavretSaada2005SerAfetado}, edição que torna acessível ao leitor brasileiro a formulação metodológica que Favret-Saada sintetiza retrospectivamente sobre a pesquisa do Bocage.} Essa posição encontra eco no que Strathern descreve como recriação imaginativa dos efeitos do próprio campo, gesto que torna a escrita antropológica parte do nexo entre pesquisa e teoria \parencite[p.~299]{Kofes2015Maconaria}. É a partir dessa sequência argumentativa que Kofes recupera Victor Turner: a tensão entre método e experiência se torna desnecessária quando dissolvida pela dramatização narrativa \parencite[p.~301]{Kofes2015Maconaria}. A experiência, ao se expressar narrativamente, conecta eventos e afecções e lhes confere significações e valores; essa expressão narrativa da experiência é o que Turner chama de \textit{structured unit of experience}, unidade estrutural da experiência. A narrativa, ao ligar o que se concebe como cesura, cria uma estrutura de relação, a estrutura de experiência. É por isso que a narrativa da etnógrafa é parte do método: o método se ensaia conforme se experimenta, e só pode ser narrado retrospectivamente porque se constitui junto com o objeto, no encontro de campo.\footnote{A formulação de Suely Kofes ressoa com o conceito de \textit{method assemblage} proposto por John Law em \textit{After Method: Mess in Social Science Research} \parencite{Law2004After}. Ambas as formulações recusam a normatividade prévia dos métodos convencionais, embora por vias distintas: Law aponta para a multiplicidade das realidades que o método produz e para a inseparabilidade entre o método e os mundos que ele torna pensáveis; Kofes sustenta a recusa pela via da experiência, argumentando que o método se ensaia através da experiência, inclusive da experiência de campo, e que a escrita antropológica é parte constitutiva desse ensaio \parencite[pp.~298--299]{Kofes2015Maconaria}.}

Essa concepção de etnografia como pesquisa antropológica que articula método e experiência dialoga com a teoria ator-rede de maneiras que interessam a esta tese. Primeiro, porque recusa categorias dicotômicas supostamente universais, como a separação entre humano e não humano. Segundo, porque trata o campo como processo de aprendizado: assim como Latour propõe que sigamos os atores sem decidir de antemão o que é técnico e o que é social, Kofes propõe que o método se ensaia conforme nos deixamos afetar pelo encontro com nossos interlocutores. Terceiro, porque ambas as abordagens organizam-se em torno das inscrições, dos documentos, das narrativas, dos artefatos e dos traçados que constituem simultaneamente nossos dados e nossos modos de produzir conhecimento.

A etnobiografia, desdobramento metodológico que Kofes formula em \textit{Uma trajetória, em narrativas} \parencite{kofes2001trajetoria}, também pode ser uma forma de investigar instituições, como faço nesta tese. Kofes adota o conceito de pessoa, em vez do de indivíduo, a partir de Marilyn Strathern, para quem a pessoa abriga o potencial para relações e é, ela própria, constituída por uma matriz de relações sociais, de sociabilidade, conforme exposto em \textit{Property, Substance and Effect} \parencite{strathern2000pse}.\footnote{A concepção de pessoa que Suely Kofes mobiliza em diálogo com Marilyn Strathern parte de uma formulação relacional da pessoa, desenvolvida pela autora a partir de sua etnografia na Melanésia. Em \textit{The Gender of the Gift: Problems with Women and Problems with Society in Melanesia} \parencite{strathern1988gender}, Strathern formula a pessoa como microcosmo de relações sociais, constituída pelas ações de outros, alimentação, cuidados, trocas, e capaz de destacar partes de si para criar novas relações. O argumento aprofunda-se em \textit{Property, Substance and Effect: Anthropological Essays on Persons and Things} \parencite{strathern2000pse} e em \textit{Partial Connections} \parencite{strathern2004partial}. Strathern designa essa pessoa como \textit{dividual}, entidade compósita e partível: na economia da dádiva melanésia, as pessoas trocam perspectivas, e o que circula nas trocas são porções de pessoa incorporadas em substâncias e objetos. A formulação contrasta com a noção euro-americana de indivíduo como unidade atômica e indivisível, que se completa pela posse de si mesmo. Strathern propõe pensar a pessoa como circuito de conexões parciais, em que cada parte estende a outra sem encerrar-se em uma identidade unitária final. Quando Kofes recorre a essa formulação para sustentar sua etnobiografia, transporta para o registro etnográfico a superação da pessoa como entidade preexistente à sociedade onde vive: a pessoa, e por extensão a etnógrafa e o coletivo etnografado, é o efeito momentâneo de uma matriz de relações que a constitui. A tradução brasileira dos ensaios de Strathern sobre o efeito etnográfico está disponível em coletânea \parencite{strathern2014efeito}, edição que torna acessível ao leitor brasileiro o conjunto de elaborações sobre o que Strathern chama de efeito etnográfico.} Se, como Kofes, é possível fazer a etnobiografia de uma Loja maçônica, então é igualmente possível fazer a etnobiografia de um centro de pesquisa em inteligência artificial. Isso significa reconstruir a trajetória e a vida de um coletivo, abrangendo suas práticas e estruturas, suas fundações e refundações, suas alianças e rupturas, suas transformações ao longo do tempo. Assim, o C4AI é também um actante cuja biografia se entrelaça com a minha trajetória de pesquisadora, e nesse entrelaçamento ressoa a formulação de Roy Wagner que Kofes recupera no mesmo capítulo: todo entendimento de uma outra cultura é um experimento com a nossa própria \parencite[p.~296]{Kofes2015Maconaria}. Esse gesto abre a possibilidade de pensar como cientistas e engenheiros fazem inteligência artificial a partir das mediações que a etnografia e a teoria ator-rede tornam visíveis nesta pesquisa.

Essa formulação ressoa com a atenção de Latour aos porta-vozes e às cadeias de tradução. Assim como o cientista no laboratório fala em nome de entidades que não podem falar por si mesmas, a antropóloga fala em nome de interlocutores cujas vozes são traduzidas, editadas, recontextualizadas no texto final. A responsabilidade dessa tradução é constitutiva do trabalho antropológico tanto quanto do trabalho científico. Concretamente, baseio-me em técnicas qualitativas diversas: observação de eventos acadêmicos e corporativos, entrevistas semiestruturadas com líderes de grupos de pesquisa, análise documental de acordos institucionais e acompanhamento de atividades de divulgação científica, de forma presencial ou remota. Essa diversidade de fontes caracteriza o que Kofes chama de pesquisa artesanal: trabalho que articula observação rigorosa e registro documental com a sensibilidade da experiência antropológica. Os dados desta tese provêm de conjunto de observações, escuta, conversas, anotações, entrevistas e convivência com pessoas ligadas ao Centro ao longo de quase quatro anos (março de 2022 a dezembro de 2025). Embora a parceria C4AI-IBM tenha sido inaugurada em outubro de 2020, antes do início do trabalho de campo, a articulação entre observação direta a partir de 2022, entrevistas retrospectivas e análise de relatórios institucionais cobrindo o período 2020--2025 permitiu reconstituir o ciclo completo da parceria, da fundação à dissolução. As entrevistas formais com Fábio Cozman (setembro de 2025) e Cláudio Pinhanez (outubro de 2025) foram realizadas após a comunicação interna da IBM sobre não-renovação da parceria (agosto de 2025), mas antes do anúncio público. Na época das entrevistas, nenhum dos interlocutores mencionou o fim iminente da parceria. Somente em dezembro de 2025, quando a IBM comunicou publicamente o encerramento, soube que ambos já conheciam essa decisão no momento das entrevistas. Essa descoberta retrospectiva conferiu dimensão temporal singular às narrativas: o que havia registrado como relatos sobre um projeto vivo revelou-se, em retrospecto, como testemunhos de quem já sabia de sua morte iminente mas não podia ou não quis dizê-lo.
 
A imersão prolongada exigiu também aprendizado de linguagens e códigos específicos. Assim como Kofes precisou decifrar a grafia e as abreviaturas maçônicas, precisei aprender a linguagem dos pesquisadores em inteligência artificial, abrangendo o vocabulário técnico, os modos de argumentar, as hierarquias implícitas entre abordagens e os critérios tácitos de avaliação. Esse processo de alfabetização foi condição de possibilidade para o aprofundamento da inteligência artificial. Mas a ausência inicial de conhecimento especializado teve um impacto duplo. Por um lado, constituiu uma vantagem metodológica, ao tratar a atividade científica como elemento novo, evitei a tendência de adotar antecipadamente a distinção entre social e técnico. Por outro lado, trouxe vulnerabilidade quando o foco mudou para a necessidade de aprender com os cientistas. Foi assim que compreendi, por exemplo, que o laboratório do cientista da computação é o computador, ideia aparentemente óbvia mas que possui implicações que merecem atenção. O computador, portátil e móvel, é mais do que simples objeto ou mediador: é o próprio laboratório, um actante, ponto na rede e simultaneamente também uma rede. Mas essa mobilidade tem limites que só se tornam visíveis quando a infraestrutura falha ou exige intervenção. Latour observa que objetos técnicos, enquanto funcionam, são intermediários silenciosos e mudos; é a crise, a quebra, o mau funcionamento, que revela sua existência, suas partes e sua agência \parencite[cap.~6]{Latour2017Esperanca}.\footnote{Em 4 de março de 2026, Luciano Antonio Digiampietri, pesquisador do Grupo de Pesquisa em Inteligência Artificial da EACH-USP (GrIA), fundado em 2006 e que divide com o C4AI a infraestrutura computacional, publicou no Instagram uma fotografia da placa de entrada do Centro com a legenda: \enquote{Hoje foi dia de presencial no laboratório. (Tive que apertar um botão, no servidor.)} (Digiampietri, 2026, comunicação pessoal). O servidor em questão executa experimentos de IA, hospeda modelos de linguagem e realiza \textit{fine-tuning} para tarefas como detecção de \textit{bots} e identificação de notícias falsas. Luciano acessa o servidor remotamente de sua sala na EACH, a cem metros de distância, da reitoria da USP, ou de sua casa, e seus orientandos também o acessam remotamente. Naquele dia, o servidor desligou e não religou automaticamente, e foi necessário que ele fosse até a máquina para pressionar o botão de reinicialização. Em seguida, voltou a acessá-lo remotamente, da sala a cem metros de distância. A postagem condensa o paradoxo: o laboratório distribuído, que dispensa presença contínua, convocou o corpo do pesquisador para um gesto mínimo e indispensável. Mostra ainda que o C4AI não construiu sua infraestrutura do zero, instalou-se sobre redes já constituídas, como o GrIA, que existia desde 2006. Luciano explicou o contexto à pesquisadora via mensagem privada no Instagram em março de 2026 e autorizou expressamente o uso da postagem e da explicação como dado etnográfico; o termo de consentimento assinado encontra-se em posse da autora.} O laboratório de IA às vezes exige que alguém se desloque até onde as máquinas estão quando elas desligam ou param de funcionar por algum motivo para \enquote{apertar um botão}.

 % =====================================================================
% Cap. 1 — Fechamento da seção: articulação Kofes-Latour em parágrafo único
%
% Substitui o conjunto de 7 parágrafos que ia de "Apoio-me em diferentes
% abordagens..." até "...essa decisão de corte é parte do trabalho da
% etnógrafa".
%
% Os 7 parágrafos foram condensados em UM parágrafo longo e denso. Foram
% dispensados: a apresentação ECT/ESCT, a menção nominal à sócio-lógica
% (operação absorvida na frase sobre "redes mais fortes ou mais fracas"),
% o anúncio dos capítulo s 3 e 4 e a referência à linhagem específica.
% Foram preservados: a articulação Kofes-Latour, o conceito de hinterland
% (com a nota expandida), os dispositivos de inscrição como modulações do
% mesmo problema, e a transição para "Cortando redes".
% =====================================================================

Tanto a teoria ator-rede de Bruno Latour e a etnografia heterodoxa de Suely Kofes funcionam juntos nesta tese porque enfrentam, a partir de temas e questões distintas, o mesmo problema. Latour o formula no registro epistemológico, onde a coerência sobre o que os atores dizem ou fazem deve ser realizada dentro do próprio movimento de associação em que as redes se fazem e se tornam mais fortes ou mais fracas, onde os fatos sempre estarão relacionados ao tamanho de uma rede e da quantidade de suas associações \parencite[p.~331]{Latour2011CienciaEmAcao}. Kofes formula a questão a partir da escrita etnográfica, pela qual o método ensaia-se conforme a etnógrafa se deixa afetar pelo que encontra, e esse ensaio só pode ser narrado retrospectivamente. Em Latour, esse gesto se realiza como seguir os actantes e mapear as redes que eles compõem; em Kofes, como deixar-se afetar pelo encontro de campo e compor a narrativa em que a experiência se estrutura. A própria adoção da teoria ator-rede nesta tese emergiu e foi personalizada de acordo com o processo de escrita, e isso é, na formulação de Kofes, o método ensaiando-se através da experiência. O gesto tem precedente etnográfico em \textit{Vida de Laboratório}, onde \textcite{Latour1997VidaLaboratorio} mostraram que o \textit{hinterland} de instrumentos de inscrição, protocolos experimentais e negociações constitui a própria ciência \parencite[p.~38]{Law2004After}.\footnote{O conceito de \textit{hinterland} é desenvolvido por Law em \textit{After Method} \parencite[pp.~27--42, 159]{Law2004After} e definido no glossário como \enquote{a bundle of indefinitely extending and more or less routinised and costly literary and material relations that include statements about reality and the realities themselves; a hinterland includes inscription devices, and enacts a topography of reality possibilities, impossibilities, and probabilities. A concrete metaphor for absence and presence} \parencite[p.~159]{Law2004After}. Mantenho o termo em itálico, sem traduzir, ao longo da tese, dado que \textit{After Method} não tem tradução brasileira disponível e que nenhuma palavra do português recobre integralmente o campo semântico que Law mobiliza, articulando a imagem geográfica do território interior que abastece um porto sem aparecer nas trocas desse porto, a topografia de possibilidades de realidade encenada por dispositivos de inscrição, e a metáfora concreta para a relação entre presença e ausência. Discuto a aproximação com a palavra \enquote{contexto}, recuperada pela etimologia de \textit{contexere} (tecer junto), na subseção~\ref{subsec:cortando-redes}. Em termos metodológicos, o \textit{hinterland} de uma descrição é o conjunto de práticas, pressupostos, infraestruturas e recursos sedimentados que sustentam o que aparece sem serem tematizados. O conceito é distinto, mas inseparável da noção de \textit{corte}, desenvolvida na subseção~\ref{subsec:cortando-redes} a partir de \textcite{strathern2011cortando}: o corte é a operação ativa que delimita o que entra e o que sai da análise; o \textit{hinterland} é a condição material e discursiva sedimentada que torna essa operação possível antes mesmo de qualquer gesto analítico, e que continua operando silenciosamente para que o objeto da análise possa existir. No contexto desta tese, o conceito retorna nos capítulos~\ref{capitulo3} e~\ref{capitulo4} para designar a rede longa de contratos, políticas de fomento e genealogia institucional que deixa de ser objeto direto de análise quando a atenção se volta para as práticas situadas de Marcelo, mas que não desaparece, pois continua sendo condição de possibilidade do que aparece.} Aquilo que Law nomeia como problema, fazer pesquisa quando o próprio fazer modifica aquilo que se investiga, opera nesta tese através de todos os dispositivos de inscrição mobilizados: os diagramas em Mermaid que performam conexões entre actantes, a curadoria bibliográfica que produz linhagens teóricas, a narrativa etnográfica que constrói o C4AI que documenta, e a composição com \texttt{claude.ai}que torna essa agência impossível de naturalizar. Cada dispositivo de inscrição \parencite{Latour2011CienciaEmAcao} decide o que torna visível e o que invisibiliza, e por isso a rede que a descrição abre precisa ser interrompida em algum ponto, isto é, precisa ser cortada.
%%%%%

\subsection{Cortando redes: sobre os limites da descrição etnográfica}
\label{subsec:cortando-redes}

\epigrafe{%
Ninguém vive em todos os lugares; todo mundo vive em algum lugar. Nada está conectado a tudo; tudo está conectado a alguma coisa.}%
{Donna Haraway, \textit{Ficar com o problema}, 2023 \parencite[p.~60]{Haraway2023Ficar}.}

A pesquisa com o C4AI colocou-me diante do dilema de como descrever uma rede que aparenta se expandir sem fim. O Centro está ligado a cerca de 250 pesquisadores, cada qual com trajetórias próprias; à USP, com toda a sua história institucional; à FAPESP, imbricada em políticas científicas estaduais e nacionais; à IBM, uma corporação global com mais de um século de existência; a projetos e sistemas específicos como Spira e Corpus Carolina; a infraestruturas computacionais espalhadas pelo planeta; ao mercado global de IA; a colaborações internacionais; e ao contexto da pandemia de Covid-19. Ao seguir os atores, eu era constantemente levada a novos lugares (e a novas questões) e, em algum ponto, eu tinha que parar de segui-los e de fazer associações, isto é, \textit{cortar} a rede.

\textcite{latour1996clarifications} explicita as propriedades topológicas que tornam o conceito de rede simultaneamente potente e problemático para a descrição etnográfica. Em \textit{On actor-network theory: a few clarifications}, ele caracteriza a rede através de três deslocamentos em relação às metáforas espaciais convencionais. O primeiro abandona a tirania da distância: elementos próximos podem estar infinitamente distantes se as conexões são analisadas, e elementos aparentemente remotos podem estar próximos quando seus elos são trazidos à descrição \parencite[p.~4]{latour1996clarifications}. O segundo supera a oposição micro-macro: ``uma rede nunca é maior que outra, ela é simplesmente mais longa ou mais intensamente conectada'' \parencite[p.~5]{latour1996clarifications}. O terceiro recusa a distinção interno-externo: ``uma rede é toda fronteira, sem interior nem exterior. A única questão que se pode colocar é se há ou não conexão estabelecida entre dois elementos'' \parencite[p.~6]{latour1996clarifications}. Esses três movimentos liberam a análise das categorias herdadas de níveis, esferas, territórios e estruturas. A consequência metodológica dessa formulação é que a rede não traz consigo limite ontológico interno. ``Literalmente não há nada além de redes, não há nada entre elas'' \parencite[p.~4]{latour1996clarifications}: cada conexão abre para outras conexões, cada actante conecta-se a outros actantes, e o trabalho da etnógrafa consiste em seguir esses elos. ``Cada laço, por mais forte que seja, é ele próprio tecido de fios mais fracos'' \parencite[p.~3]{latour1996clarifications}. O conceito que estipula como descrever as associações não estipula onde a descrição deve parar. O dilema enfrentado em campo era a forma empírica desse impasse conceitual: por onde terminar de seguir o C4AI?

\textcite{strathern2011cortando} formula esse dilema ao mostrar que redes não se tornam objetos de conhecimento sem antes serem cortadas. Em \textit{Cortando a rede}, ela retoma o conceito latouriano para identificar uma de suas propriedades, a autoilimitação: o conceito de rede funciona como metáfora das ``ilimitadas extensões e entrelaçamentos dos fenômenos'' \parencite[p.~7]{strathern2011cortando}. Frente a essa propriedade, ``a análise, como a interpretação, precisa ter um ponto; precisa ser encenada como um ponto de chegada'' \parencite[p.~7]{strathern2011cortando}. Uma rede é tão longa quanto seus elementos possam ser enumerados, e essa enumeração ``pressupõe uma totalização; ou seja, uma enumeração residindo em um objeto identificável (a soma). Ao parar, a rede pode ser `cortada' em um ponto'' \parencite[p.~7]{strathern2011cortando}. Strathern articula rede e híbrido como duas figuras da mesma operação: ``se tomarmos certos tipos de redes como híbridos socialmente expandidos então podemos tomar híbridos como redes condensadas. O trabalho de condensação trabalha como uma totalização ou parada'' \parencite[pp.~7--8]{strathern2011cortando}. Cada corte produz simultaneamente uma rede estendida, com elos enumeráveis e seguíveis, e um híbrido condensado, que retém os elementos heterogêneos da rede como objeto manipulável de conhecimento ou de transação. O corte não é gesto externo imposto a uma rede que estaria já dada, pois, as próprias práticas observadas realizam cortes, condensam relações em objetos, totalizam redes em pontos de chegada. O dote nupcial entre os \textit{'are'are} das Ilhas Salomão converte uma rede de prestações em objeto consumível na sequência mortuária; o sistema de patentes converte uma rede de quarenta colaboradores científicos em uma única invenção atribuível a seis inventores \parencite[pp.~8--9]{strathern2011cortando}.\footnote{Strathern desenvolve essa argumentação em diálogo etnográfico com sociedades cujos sistemas de troca tornam visíveis operações de corte que os euroamericanos tendem a naturalizar. Entre os \textit{'are'are} das Ilhas Salomão, a pessoa viva é composta de corpo, sopro e imagem, e na sequência funerária esses componentes são convertidos, respectivamente, em taro, porcos e dinheiro de conchas, este último figurando como a própria forma da imagem ancestral \parencite[pp.~10--12]{strathern2011cortando}. Os \textit{'are'are} designam essa sequência finalizadora como ``parada'' ou ``freio'': o dinheiro de conchas funciona simultaneamente como circulador e como interruptor de fluxos, condensando trocas anteriores e bloqueando reivindicações futuras.} O que Strathern quer dizer é que a análise reconhece e transporta cortes que as práticas já operam sem a intervenção do antropólogo. No argumento de Strathern, o conceito de \textit{propriedade}, por exemplo, funciona como mecanismo euroamericano por excelência de cortar redes. Ela é ``poderosa por causa de seus efeitos duplos, sendo simultaneamente uma questão de pertencer e de propriedade'' \parencite[p.~16]{strathern2011cortando}: separa quem participa de quem não participa, e converte o trabalho cooperativo de uma rede longa em objeto sobre o qual se exerce direito. ``Quando a tecnologia pode aumentar as redes, a condição de propriedade pode garantidamente cortá-las em um dado tamanho'' \parencite[pp.~16--17]{strathern2011cortando}.

Pensar o corte da rede a partir do conceito de propriedade permitiu-me identificar, no caso desta tese, onde os cortes já operavam na prática antes que eu mesma tivesse de produzi-los analiticamente. O exemplo mais dramático disso foi a rescisão da parceria entre C4AI e IBM em 2025, examinada no capítulo~\ref{capitulo3}. O que descrevo com Latour como desestabilização ou reconfiguração da rede foi, com Strathern, um corte que mostrou o poder do direito e das decisões contratuais, como sobre a propriedade intelectual, patentes, código aberto e assim por diante.

Mobilizo \textcite{Haraway2023Ficar} para reforçar o argumento de Strathern com a figura do \enquote{pensamento tentacular} que dá epígrafe a esta seção: nada se conecta a tudo, tudo se conecta a algo. Cada elo, cada nó da rede amarra e se liga a uma conexão particular. A rede se ramifica e cada actante conecta-se a outros actantes específicos, não porque seja ontologicamente ilimitada, mas porque é assim que ela existe como rede que faz aquilo que faz. O argumento orienta as escolhas de corte que farei ao longo da tese, e antecipa em particular a maneira como acompanharei, no capítulo~\ref{capitulo4}, a rede multiespécie\footnote{A noção multiespécie não se restringe a animais e plantas, e inclui também microorganismos e entidades cuja classificação como vivas é objeto de controvérsia. Os vírus comparecem nesse quadro em posição singular, pois não possuem metabolismo próprio, não mantêm homeostase e não se replicam fora de uma célula hospedeira, o que levou a virologia do século XX a classificá-los como entidades replicantes em uma fronteira indeterminada com o vivo. As descobertas dos vírus gigantes a partir do início dos anos 2000, em particular Mimivírus, Pandoravírus e Pithovírus, reabriram o debate sobre o estatuto vivo dos vírus, e a hipótese de que vírus participaram da evolução do DNA e de transições evolutivas decisivas é defendida pelo virólogo Patrick Forterre em \enquote{Defining Life: The Virus Viewpoint} (2010). Em chave antropológica, a etnografia de referência para pensar o vírus como actante multiespécie é \enquote{Viral Clouds: Becoming H5N1 in Indonesia}, de Celia Lowe \parencite{Lowe2010ViralClouds}, na qual a antropóloga propõe a figura da nuvem viral para descrever a indeterminação ontológica do H5N1, em diálogo com a noção de quase-espécies da virologia molecular, e \textit{Avian Reservoirs: Virus Hunters and Birdwatchers in Chinese Sentinel Posts}, de Frédéric Keck \parencite{Keck2020AvianReservoirs}, que acompanha por uma década as práticas de vigilância e preparação para pandemia em Hong Kong, Singapura e Taiwan. No caso brasileiro, o ensaio \enquote{Como compor com um vírus!? Reflexões sobre os \textit{animal studies} no tempo das pandemias}, de Eliane Sebeika Rapchan e Fagner Carniel \parencite{RapchanCarniel2021Compor}, e \enquote{Covid-19, biossegurança e antropologia}, de Jean Segata \parencite{Segata2020Biosseguranca}, ambos publicados em \textit{Horizontes Antropológicos}, mobilizam o vocabulário multiespécie em diálogo com o caso do SARS-CoV-2. Para esta tese, a indeterminação ontológica do vírus tem implicação direta para o argumento sobre existências parciais e simpoiese desenvolvido neste capítulo, pois o coronavírus existe somente na relação com a maquinaria celular do hospedeiro. Estender o vocabulário multiespécie de Haraway, van Dooren e Rose ao SARS-CoV-2 é gesto deliberado, e a parceria que esse vocabulário nomeia é, no caso humano, frequentemente fatal, porque a existência do coronavírus é, desde o início, uma composição mortal com outras espécies.} onde investigo o sistema Spira de detecção de insuficiência respiratória por análise acústica da voz. Aquele capítulo seguirá uma trama em que um vírus, vozes de pacientes hospitalizados, equipamentos de gravação, modelos de aprendizado de máquina e protocolos clínicos compõem um arranjo cujas conexões são específicas, e cuja especificidade explica tanto o que o sistema consegue fazer quanto o que ele deixa de fazer.

A formulação que dá epígrafe a esta seção, embora atribuída a Haraway no lugar onde a leio, tem uma genealogia que importa para o argumento. A frase aparece pela primeira vez em Deborah Bird Rose, em ensaio de 2008 sobre quatro modos de luto na Austrália aborígene em que Rose discute o envenenamento de dingos e o efeito em cadeia sobre as comunidades multiespecíficas dos Yarralin. Thom van Dooren retoma a formulação em \textit{Flight Ways: Life and Loss at the Edge of Extinction} (2014), no capítulo dedicado à mortalidade em massa de abutres do gênero \textit{Gyps} na Índia por intoxicação pelo anti-inflamatório diclofenaco, contraído de modo não intencional pelo consumo de carcaças de gado tratadas com o medicamento. Van Dooren mobiliza a formulação como crítica à ecologia holística que apaga as especificidades das relações em nome de uma harmonia geral, e mostra que a desconexão das cadeias multiespécies recai de modo desigual sobre diferentes coletivos: as populações rurais pobres que dependiam dos abutres para a remoção sanitária de carcaças, os coletores de ossos cuja prática de subsistência foi reorganizada pela escassez das aves, e a comunidade parsi cuja prática funerária milenar nas torres do silêncio depende da presença dos abutres. Em nota a essa passagem, Haraway cita van Dooren e oferece uma maneira de pensar o que está em questão, formulação que servirá para refletir sobre tecnologia em geral e, em particular, sobre a inteligência artificial.

Em \textit{Atlas of AI}, Crawford desenvolve, a propósito da extração mineral que sustenta a infraestrutura computacional contemporânea, uma argumentação de estrutura semelhante à de van Dooren: os benefícios da extração foram capturados por poucos, enquanto os custos recaem sobre comunidades específicas. O primeiro capítulo do livro passa pela comunidade não incorporada de Silver Peak, em Nevada, onde cerca de 125 pessoas vivem sobre o lago subterrâneo de lítio que abastece a única mina ativa desse mineral nos Estados Unidos; pela região do Congo, onde os lucros da mineração de tântalo, estanho, tungstênio e ouro financiaram operações militares ao longo de décadas de conflito, com a morte de milhares de pessoas e o deslocamento de milhões; por Baotou, na Mongólia Interior, onde um lago artificial de cinco milhas e meia de diâmetro acumula mais de 180 milhões de toneladas de resíduos do processamento de terras-raras; e pelas ilhas indonésias de Bangka e Belitung, onde mineradores informais raspam o leito do mar em busca do estanho usado em semicondutores, em uma paisagem marcada pelas covas dos que morreram cavando ao longo dos séculos \parencite[pp.~24--38]{Crawford2021Atlas}. Crawford retoma a observação de Georgius Agricola (fundador da geologia como ciência), em 1555, de que a mineração só se sustenta como atividade lucrativa porque não precisa contabilizar seus custos reais, o dano ambiental, o adoecimento e a morte dos mineradores, a perda das comunidades deslocadas \parencite[p.~27]{Crawford2021Atlas}, e argumenta que a ignorância sobre as cadeias de suprimento é constitutiva do capitalismo e se mantém por meio de subcontratações, intermediários e auditorias que avaliam fundições, e não as minas \parencite[pp.~34--36]{Crawford2021Atlas}. Assim, a inteligência artificial, que precisa de uma quantidade gigantesca de infraestrutura computacional que por sua vez depende da extração de minérios, água e energia elétrica,\footnote{O AI Now Institute mantém uma área de pesquisa dedicada ao clima que desenvolve essa argumentação em registros complementares, articulando as cadeias materiais da inteligência artificial a questões de justiça climática, hídrica e laboral. Entre as publicações disponíveis, \enquote{Climate Justice and Labor Rights, Part I: AI Supply Chains and Workflows} (2019) e \enquote{Climate Justice and Labor Rights, Part 2: Labor Organizing and Environmental Justice in Tech, Past and Present} (2023) articulam a extração de minérios e a operação de \textit{datacenters} às condições de trabalho ao longo da cadeia de produção; o \textit{Water Justice and Technology Report} (2022) detalha o consumo hídrico das infraestruturas computacionais e seus efeitos sobre comunidades específicas; \enquote{A New AI Lexicon: Sustainability} (2021) examina o vocabulário de sustentabilidade mobilizado pelo setor; e \enquote{AI and Climate Change: How they're connected, and what we can do about it} (2019) situa a inteligência artificial no quadro mais amplo da crise climática. O conjunto está reunido em \url{https://ainowinstitute.org/research-areas/climate}.} lida nesses termos, é muito mais material e imbricada no mundo físico e em seus habitantes (humanos ou não) do que costuma parecer:

\begin{citacaoabnt}
A corrente de filosofia ecológica holística que insiste que ``tudo está conectado a tudo'' não nos ajudará aqui. Antes, diríamos que tudo está conectado a alguma coisa, que por sua vez está conectada a outra coisa. Ainda que, em última análise, tudo possa estar conectado entre si, a especificidade e a proximidade das conexões são importantes -- com quem estamos vinculados e de que maneiras. A vida e a morte acontecem nessas relações. Por isso, precisamos entender como se emaranham certas comunidades humanas, em particular, assim como outras comunidades de seres vivos, e de que maneiras esses emaranhados estão implicados na produção de sua extinção mútua e de seus respectivos padrões de morte amplificada \parencite[p.~60, n.~4]{Haraway2023Ficar}.
\end{citacaoabnt}

A formulação \enquote{tudo está conectado a alguma coisa, que por sua vez está conectada a outra coisa} aparece pela primeira vez em Rose (2008), em ensaio sobre o luto multiespécie na Austrália aborígene. Van Dooren retoma a formulação em \textit{Flight Ways} (2014) para descrever a dinâmica que ele chama de \textit{double death}, produzida pela intoxicação dos abutres \textit{Gyps} na Índia, em que os corpos dos abutres mortos não se reintegram às cadeias da vida, e a ausência dos vivos deixa cadáveres bovinos em decomposição que aumentam a incidência de raiva e a contaminação por antraz, com efeitos desigualmente distribuídos sobre as comunidades pobres da região.\footnote{Van Dooren articula seu argumento em torno de três figuras conceituais que oferecem um modo para pensar relações multiespécies. A primeira é a do \textit{nó}, que van Dooren retoma de Deborah Bird Rose: cada actante é um \textit{knot of embodied time}, ponto de conexão entre gerações, espécies e escalas, em um tempo incorporado no espaço físico e na sobrevivência de cada geração. A figura ressoa com o vocabulário latouriano da cadeia de translações que percorre esta tese e com a noção strathernina da pessoa como microcosmo de relações que recuperei a partir de Kofes, e foi importante para descrever, no capítulo~\ref{capitulo4}, o paciente hospitalizado cuja voz alimenta o conjunto de dados do Spira, e que respira diante de um microfone enquanto um vírus se replica em seu epitélio respiratório. A segunda figura é a da \textit{dull edge}, a borda da extinção, o desfazimento de uma trama de relações que começa muito antes do evento catastrófico e continua a reverberar muito depois; a pandemia de Covid-19 tem essa qualidade temporal, e o capítulo~\ref{capitulo4} acompanhará o modo como as inscrições acústicas produzidas no auge da emergência sanitária no Brasil carregam um tempo que excede o instante da gravação. A terceira figura é a da \textit{double death}, a dinâmica em que o desfazimento das conexões produz, ele próprio, novas mortes, e em que essas mortes se distribuem desigualmente, atingindo desproporcionalmente comunidades cuja exposição às relações multiespécies em crise era já estruturalmente maior. As três figuras articulam-se no argumento de van Dooren, e ressoam no capítulo~\ref{capitulo3}, na genealogia que \textcite{Crawford2021Atlas} traça da computação contemporânea aos cartões perfurados de Hollerith, e no capítulo~\ref{capitulo4}, em que a distribuição desigual da pandemia no Brasil aparecerá como condição da rede que tornou possível o conjunto de dados sobre o qual o Spira foi treinado.} A relação Rose-van Dooren-Haraway que o trecho citado preserva é ela mesma um exemplo desse tipo de transmissão de conceitos por proximidades específicas, em que cada passagem por autores e temáticas diferentes acrescenta novas implicações analíticas. Recuperar essa relação importa para esta tese porque mostra como uma proposição teórica viaja por redes que conectam diferentes tecnologias (anti-inflamatório, vacina, IA) e diferentes espécies (humanos, aves, mamíferos, bactérias, vírus, e até mesmo minérios e reações físicas e químicas), fazendo com que esses elementos conversem de uma maneira que não seria possível sem esse tipo de empréstimo de modos de pensar, conforme elaborado por Strathern em \textit{Reproducing the Future: Essays on Anthropology, Kinship, and the New Reproductive Technologies} (1992), citada por Haraway, que a considera uma \enquote{etnógrafa das práticas de pensamento} \parencite[p.~69]{Haraway2023Ficar} para enfatizar a responsabilidade ética e política na produção de conhecimento em um tipo de \enquote{compostagem intelectual}, pensando com \parencite[p.~69]{Haraway2023Ficar}. Karen Barad, em \textit{Meeting the Universe Halfway} (2007), também sustenta que o pensamento não é só um exercício abstrato, mas também uma prática material, pois diz respeito ao modo como escolhemos conceitos e categorias que acabam definindo o que se torna inteligível e o que não, o que existe e o que não existe. \textcite{stengers2018another} sustenta que as ideias têm uma eficácia própria e podem \enquote{envenenar ou ativar} processos \parencite[p.~142]{stengers2018another}, abrindo ou fechando possibilidades de futuro. \textcite{costa2024manchando} aplica esse modo de pensamento em sua tese de doutorado ao escolher ferramentas metodológicas como a \enquote{cesta} e a \enquote{armadilha} para pesquisar a menstruação, entendendo que essas escolhas conceituais determinam quais vozes serão ouvidas e quais realidades serão construídas no texto acadêmico. A questão importa para esta tese porque diz respeito ao que minhas escolhas conceituais farão e a quem darão voz. O capítulo~\ref{capitulo4} é uma experimentação sobre como essas vozes podem aparecer de lugares improváveis.\footnote{Apresentei uma versão preliminar do capítulo~\ref{capitulo4} no seminário de dez anos do La'grima (Laboratório Antropológico de Grafia e Imagem, IFC-Unicamp), com foco nas imagens da voz, espectrogramas em escala Mel renderizados como mapas de calor, geradas computacionalmente para o treinamento e a inferência do sistema Spira. Segui os actantes e as cadeias de translações que eles e eu fomos construindo até chegar às vozes das pessoas enfermas hospitalizadas nas enfermarias de Covid-19, designadas pelos pesquisadores que desenvolveram o sistema como \enquote{covideiros}, cujas gravações serviram como fonte de dados para esse sistema de inteligência artificial. A composição dessas imagens, que fui construindo com a ajuda do \texttt{claude.ai}em linguagem de programação Python a partir das gravações de vozes reais que busquei no repositório do projeto no GitHub, chamou-me a atenção como vestígio dessas pessoas: o que viajou pelas cadeias de translações até chegar à entrada do modelo computacional, e o que restou delas nessa viagem, foi essa marca acústica visualizada em espectrograma. Ao mencionar a palavra \emph{vestígio} durante a apresentação, uma pessoa da plateia perguntou de que tipo de vestígio se tratava e a quem ele diria alguma coisa. A pergunta operou um corte analítico no meu material, pois deslocou a discussão da imagem como representação técnica para a imagem como inscrição que carrega rastros de pessoas específicas e que, dependendo de como é descrita, fala a interlocutores diferentes. Para os engenheiros e cientistas da computação que treinam o modelo, o espectrograma é \textit{feature} de entrada (recurso específico), matriz numérica de coeficientes Mel sobre a qual a rede neural opera convoluções e calcula gradientes. Para a equipe clínica que coletou as gravações, é registro de uma respiração específica em um momento específico do quadro respiratório do paciente. Para a pessoa hospitalizada, se ela tivesse acesso à imagem, seria possivelmente irreconhecível como sua própria voz. Para mim, é o ponto em que o vocabulário multiespécie da seção encontra a sua aplicação empírica, isto é, o vestígio acústico carrega o vírus que se replicava no epitélio respiratório no momento da gravação, a obstrução pulmonar produzida pela infecção, a microfonia da enfermaria, o protocolo de coleta padronizado pela equipe do Hospital das Clínicas e os parâmetros da escala Mel que tornaram a inscrição possível; quais relações foram tornadas presentes e quais ficaram no \textit{hinterland} no sentido de \textcite{Law2004After}. É exatamente o tipo de problematização que os modos de pensar relacionalmente, mobilizados nesta seção, permitem fazer aparecer. Não são só arranjos, cadeias, redes como se pensarmos com Latour, são, sobretudo, relações feitas a partir de/pela análise. Que por sua vez, também são cortes, pensando com Strathern.}

Os quatro cortes mobilizados nesta seção, o corte interpretativo de \textcite{strathern2011cortando}, o corte agencial que Karen Barad desenvolve em \textit{Meeting the Universe Halfway} (2007), a respons-habilidade de \textcite{Haraway2023Ficar} e a armadilha situada de \textcite{costa2024manchando}, podem ser reunidos em uma figuração que nomeia o que esses cortes fazem juntos. A escolha dos materiais analíticos que se costuram em uma narrativa etnográfica determina quais mundos podem ser construídos no texto e quais vozes podem ser ouvidas, conforme tematizado acima a propósito da armadilha mobilizada por \textcite{costa2024manchando} e da formulação de Haraway, recuperando Strathern, de que importa quais ideias usamos para pensar outras ideias \parencite[p.~12]{Haraway2016Staying}. Mas, a figuração ganha na língua portuguesa brasileira uma polissemia particular. \textit{Trama} nomeia, em uma única palavra, o conjunto de fios transversais que se entrelaça com a urdidura para constituir o tecido, a sequência narrativa de associações que organiza personagens e eventos em um arranjo inteligível, a articulação de interesses cujo padrão só se revela quando o enredo chega a um ponto, e o arranjo reticular de elementos cuja regularidade depende das relações entre os pontos que o constituem. O sentido têxtil de trama depende também de sua relação com a urdidura. A urdidura é o conjunto de fios longitudinais montados antes no tear, que sustentam a estrutura do tecido e que se recolhem da visibilidade no produto pronto. A trama, em contraste, é o conjunto de fios transversais que se entrelaça com a urdidura, e é o que se vê quando se olha o tecido pelo direito, por fora. A oposição urdidura/trama ressoa com o par \textit{hinterland}/presença que Law desenvolve em After Method \parencite[pp.~40--42, 84]{Law2004After} (apesar de Law partir do escopo das figurações topológicas): a urdidura sustenta o tecido sem aparecer; a trama é o que se vê.

Ainda que existam diferenças, há uma convergência da figuração \enquote{têxtil} entre as autoras e os autores nesta seção. Strathern caracteriza o conceito de Latour de rede como metáfora das \enquote{ilimitadas extensões e entrelaçamentos dos fenômenos} \parencite[p.~7]{strathern2011cortando}. \textcite{Haraway2023Ficar} faz das figuras de barbante (\textit{string figures}) o tropo central do pensamento tentacular e descreve a simpoiése como emaranhados. Law cita Derrida sobre o feixe (\textit{sheaf}) como \enquote{estrutura complexa de um tecido, um entrelaçamento que permite que diferentes fios e diferentes linhas de sentido sigam em direções diferentes} \parencite[p.~41]{Law2004After}, e descreve o artesanato (\textit{crafting}) do método como agrupamento (\textit{bundling}), tecelagem (\textit{weaving}) e reunião (\textit{gathering}) de relações \parencite[pp.~41--42, 144]{Law2004After}. \textcite[p.~3]{latour1996clarifications} caracteriza a teoria ator-rede pelo enredamento (\textit{netting}), atamento (\textit{lacing}), tecelagem (\textit{weaving}) e torção (\textit{twisting}) de laços fracos por si mesmos. \textcite{costa2024manchando} mobiliza a cesta e a armadilha, ambas práticas de trançar fibras naturais.\footnote{A noção de armadilha em \textcite[pp.~306--309]{costa2024manchando} é articulada em três fontes. De Alfred Gell, em \enquote{Vogel's Net} (1996), vem a definição da armadilha como objeto que materializa o conhecimento de si e do outro, modelo funcional do capturador (substituído por uma \enquote{transdução sensorial}) e simultaneamente do cativo (cujos comportamentos e desejos a armadilha precisa emular para funcionar). Do pensamento Guarani-Mbya vem a inflexão crítica decisiva: a presa só é capturada se houver anuência dos donos não humanos, o que desfaz a hierarquia trágica capturador/cativo postulada por Gell e redistribui a tragédia. De Caceres e Sales (2019), vem a figura da armadilha como cartão de memória que armazena o comportamento da presa, os modos de produção das armadilhas e os relacionamentos cosmopolíticos envolvidos na caça. Costa também explora a etimologia latina de \textit{trappa} para articular a armadilha à figura da trapaça e ao caleidoscópio enquanto dispositivo feminista de captura.}

\textcite{Ingold2015EstarVivo} desenvolve uma figuração nesse mesmo sentido ao propor o conceito de malha (\textit{meshwork}) como \enquote{emaranhado de fios e caminhos} \parencite[p.~64]{Ingold2011BeingAlive} em que cada ser vivente é uma linha que se entrelaça com outras pelo movimento de habitar o mundo. Ingold contrasta o \textit{meshwork} com a rede de Latour ao longo do capítulo 7 de \textit{Being Alive} \parencite[pp.~89--94]{Ingold2011BeingAlive}, em um diálogo encenado entre os personagens \textit{ant} e \textit{spider},\footnote{O capítulo está construído como fábula filosófica entre dois artrópodes igualmente articulados. \textit{Ant} é a formiga que defende a teoria ator-rede e dela toma o nome (act-ant), interlocutor cifrado de Bruno Latour mas também de John Law e Annemarie Mol, citados em nota pelo próprio Ingold como referências para a noção de agência distribuída pela rede. \textit{Spider} é a aranha que defende a malha e cujo nome funciona simultaneamente como anagrama: \textit{Skilled Practice Involves Developmentally Embodied Responsiveness}, em tradução literal, a prática habilidosa envolve responsividade desenvolvimentalmente incorporada. A discussão atravessa também Andy Clark, com a noção de \textit{wideware} pela qual a aranha estende seu corpo na teia, Lambros Malafouris e sua \textit{dance of agency}, e um \enquote{filósofo humano} sem nome que entoa o \enquote{penso, logo sou} cartesiano enquanto a aranha o escuta pendurada do teto da sala de aula. A discordância em relação a Latour não é, portanto, isolada e o diálogo encena uma posição feita contra toda uma constelação de teorias relacionais que tomam a agência como propriedade distribuída por entidades conectadas, e em favor de uma fenomenologia ecológica em que os viventes habitam meios e não montam redes.} e situa essa diferença no estatuto do ar, da água e do tempo, que para Ingold são meios em que os viventes estão imersos e não são nem entidades nem actantes. Ingold retoma o termo \textit{meshwork} de Henri Lefebvre, em \textit{The Production of Space} (1991, pp.~117--118), citando a observação de Lefebvre de que \enquote{a atividade prática escreve sobre a natureza, com letra rasgada}. Ingold lê essa formulação como convite a pensar a escrita como tecido de linhas, e a inscrição da atividade humana e não humana no espaço vivido como textura, e não como texto \parencite[p.~84]{Ingold2011BeingAlive}.\footnote{A elaboração canônica do conceito de malha está em \textit{Lines: A Brief History} \parencite[pp.~80--82]{Ingold2007Lines} e em \textit{Being Alive: Essays on Movement, Knowledge and Description} \parencite[pp.~63--75, 89--94]{Ingold2011BeingAlive}, com tradução brasileira de ambos pela Editora Vozes \parencite{Ingold2022Linhas, Ingold2015EstarVivo}; a retomada como ponto de partida sociológico e ecológico está em \textit{The Life of Lines} \parencite[pp.~3--4]{Ingold2015Lines}, sem tradução brasileira até a data de finalização desta tese. \textcite[pp.~210--212]{Ingold2015EstarVivo} dedica um capítulo a mostrar que a história ocidental separou os caminhos do técnico-arquitetônico e do têxtil a partir da virada arquitetônica de Alberti no século XV: o primeiro foi elevado a sistema de princípios operacionais, a \textit{tecnologia}, e o segundo foi rebaixado a artesanato residual. Antes dessa partição, a construção da catedral medieval de Chartres se assemelhava, segundo Ingold citando John Harvey em \textit{The Medieval Architect} (1974), mais a um \textit{patchwork quilt} do que a um edifício planejado a partir de uma planta prévia \parencite[p.~211]{Ingold2015EstarVivo}. A escolha contemporânea de mobilizar o vocabulário têxtil para falar de método científico, de ontologia e de conhecimento opera, no registro feminista contemporâneo, como reabilitação de um léxico historicamente associado ao trabalho de mulheres e historicamente desclassificado como não-conhecimento, gesto que atravessa as \textit{string figures} de \textcite{Haraway2023Ficar}, a sacola portadora de Ursula K. Le Guin, em \textit{The Carrier Bag Theory of Fiction} (1986), como anti-narrativa contra a história como lança, e a cesta-armadilha de \textcite{costa2024manchando}.} Para Ingold, \enquote{separar a aranha de sua teia seria como separar o pássaro do ar ou o peixe da água: retirados dessas correntes, estariam mortos} \parencite[pp.~64--65]{Ingold2011BeingAlive}, e a etnografia consistiria, assim, em seguir as linhas pelo seu movimento sem cortá-las. À princípio, foi o que tentei fazer. Todavia, a figuração da costura mobilizada nesta tese vai em outra direção, e permanece mais com Haraway e Strathern do que com Ingold: nada se conecta a tudo, tudo se conecta a alguma coisa, e a descrição etnográfica precisa ter um ponto de chegada, precisa começar e precisa parar em algum lugar. A crítica de Ingold à agência distribuída pela rede em favor dos meios que tornam possível aos seres existirem ainda exigiria onde interromper a descrição, problema que é a própria condição de possibilidade do conhecimento etnográfico no argumento de \textcite{strathern2011cortando}.

A figuração da costura sustenta-se, assim, no gesto de cortar como condição da costura, e nomeia a etnografia como prática que habita (e se habita) essa forma de fazer-com retalhos, pedaços do campo que, uma vez costurados, carregam memórias e experiências diferentes e por vezes mutuamente excludentes. Cabe ainda uma observação situada sobre a posição de onde Ingold formula sua recusa do corte. A epistemologia feminista mobilizada por Strathern e Haraway recusa a possibilidade de uma posição não marcada, de onde a etnografia pudesse seguir as linhas do mundo sem operar cortes sobre elas. O corte não é gesto que algumas antropólogas precisariam fazer enquanto outros antropólogos poderiam dispensar; é, sobretudo, parte de qualquer descrição feita a partir de algum lugar, por alguém. Seguir as linhas supõe a possibilidade de uma posição que pode acompanhar o movimento dos viventes sem precisar definir onde a descrição se interrompe. Essa suposição não se sustenta a partir do feminismo situado em que esta tese se inscreve. Onde Ingold escreve no registro de uma fenomenologia que pode permanecer com as linhas do mundo, Strathern e Haraway escrevem no registro de quem sabe que toda descrição é parcial, situada, e responsável pelos cortes que opera, inclusive aqueles que opera sem perceber.

É a partir dessa tensão que a figuração da costura ganha seus contornos práticos. Pensando com Strathern e Haraway, cortar uma rede é também costurar, e cortes e recortes de materiais heterogêneos compõem o \textit{patchwork} no estilo do \textit{crazy quilting}\footnote{O \textit{crazy quilt} é uma forma de \textit{patchwork} difundida na tradição têxtil no século XIX, marcada pelo arranjo irregular de retalhos de tamanhos e formas variadas, costurados sem padrão geométrico predefinido. Distingue-se das demais tradições de \textit{quilt} pela exibição das costuras, em que as junções entre os retalhos são marcadas por bordados decorativos como \textit{feather stitch}, \textit{herringbone}, \textit{french knots} e \textit{chain stitch}, que percorrem visivelmente as bordas em vez de ficarem ocultos pelo acabamento. A composição agrega tecidos heterogêneos em sua origem (sedas, veludos, cetins, algodões estampados, retalhos de vestuário familiar reaproveitados) em uma trama única que sustenta a coexistência das diferenças. A prática esteve historicamente associada ao trabalho doméstico feminino e à preservação narrativa de relações afetivas. Cada retalho carregava a memória de quem o usou, e a trama compunha um arquivo familiar visível. A referência canônica sobre a forma é Penny McMorris, em \textit{Crazy Quilts} (1984). A figura é mobilizada nesta tese como tropo concreto da feitura etnográfica, em que o gesto de cortar e o de costurar produzem uma trama heterogênea cuja visibilidade das junções é parte do que se descreve.} tecnoetnográfico.\footnote{O VII Simpósio Internacional da LAVITS (Rede Latino-americana de Estudos sobre Vigilância, Tecnologia e Sociedade), a ocorrer no Rio de Janeiro em agosto de 2026, sob o tema \enquote{Tramar o presente: tecnopolíticas dissidentes e ecologias do Comum}, mobiliza o vocabulário têxtil em sentido próximo ao desenvolvido nesta seção. A chamada de trabalhos do simpósio convoca a pensar as relações humano-máquina-natureza fora de dicotomias fatalistas, buscando \enquote{imaginar novas tessituras para essas relações}, e propõe o gesto de tramar como prática situada que reaprende o presente \enquote{a partir da colaboração, da insurgência, das tensões criadoras e das alianças improváveis}. A convergência entre o tropo mobilizado por essa rede de pesquisadoras e pesquisadores latino-americanos e a figuração da costura desenvolvida nesta tese sugere que o vocabulário têxtil tem aderência específica para pensar a tecnociência situada no Sul Global, e em particular para a etnografia da pesquisa em inteligência artificial.} O \textit{crazy quilt} fornece, para a figuração da costura um modelo concreto para pensar, pois, a heterogeneidade dos retalhos previamente cortados, a visibilidade das junções como parte do que se descreve (e como parte de quem descreve, suas memórias), a recusa de subordinação a um padrão geral, e a inscrição da composição na tradição artesanal do trabalho feminino correspondem ponto a ponto ao que esta tese precisa fazer ao compor materiais heterogêneos (capítulos sobre o C4AI, o Spira, a IBM, a tecnoetnografia) sem subordiná-los a uma totalização única. O gesto de cortar tem, na figuração da costura, uma qualidade difrativa, e atende aos padrões de diferença e de exclusões que as práticas de conhecimento produzem no mundo.\footnote{Haraway propõe a difração como metáfora alternativa à reflexão \parencite{Haraway1997Modest}. Nos termos da física, a reflexão produz espelhamento, supõe distância e exterioridade do observador em relação ao observado, e produz uma semelhança entre o original e a cópia. A difração opera por meio da interferência, supõe a imersão do observador no campo que descreve, e produz padrões de diferença que registram os efeitos materiais que as práticas de conhecimento produzem no mundo. Karen Barad, em \textit{Meeting the Universe Halfway} (2007), retoma e aprofunda a metáfora ao articulá-la com a categoria de intra-ação, na qual objeto e instrumento de conhecimento se constituem mutuamente em sua relação.} A figuração funciona, no sentido em que Haraway formula o termo, como entidade material-semiótica que produz mundos e significados consequentes, e que pode ser habitada e contestada.\footnote{Haraway desenvolve as figurações como tropos analíticos ao longo de sua obra. No ensaio de 1985 \enquote{A Cyborg Manifesto}, republicado em \textit{Simians, Cyborgs, and Women: The Reinvention of Nature} \parencite{Haraway1991Simians}, apresenta o ciborgue como figuração que mapeia realidades sociais e corporais e desfaz binarismos entre humano e máquina, organismo e máquina, físico e não físico \parencite{Haraway1997Modest}, articula a testemunha modesta como figuração crítica do observador científico tradicional, junto com a \textit{oncomouse} como organismo-mercadoria que rompe a fronteira entre natureza e tecnologia. Em \textcite{Haraway2023Ficar}, mobiliza o acrônimo SF (\textit{science fiction}, \textit{speculative fabulation}, \textit{speculative feminism}, \textit{science fact}, \textit{string figures}) e as figuras de barbante como tropos do pensamento tentacular, modo de pensar-com companheiros em emaranhados simpoiéticos. Para Haraway, figurações funcionam como entidades materiais-semióticas habitáveis e contestáveis, mapas de mundos que produzem significados consequentes, e tecnologias gráficas, psicológicas e políticas com as quais se compõem realidades específicas.}

A figura têxtil, \textit{pensada} a partir de Haraway e Strathern para \textit{pensar} o corte, também é \textit{pensada} por John Law em \textit{After Method} para descrever o que ele chama de \textit{method assemblage}. Law cita Derrida sobre o \textit{sheaf} como \enquote{estrutura complexa de um tecido, um entrelaçamento que permite que diferentes fios e diferentes linhas de sentido sigam em direções diferentes, ao mesmo tempo em que está sempre pronto a se atar a outros} \parencite[p.~41]{Law2004After}, e descreve o \textit{crafting} do método como operação de \textit{bundling}, \textit{weaving} e \textit{gathering} de relações heterogêneas \parencite[pp.~41--42, 144]{Law2004After}. Latour, apesar das críticas, também \textit{pensa} com esse \textit{pensamento} a resistência material \textit{pensanda} com conceitos como \enquote{\textit{netting}, \textit{lacing}, \textit{weaving}, \textit{twisting} de laços que são fracos por si mesmos} \parencite[p.~3]{latour1996clarifications}. A figuração têxtil é, portanto, recorrente e um modo interessante e promissor para \textit{pensar} as questões dos métodos como práticas de fazer relações. 

A aproximação Strathern--Haraway--Mol que esta seção vem articulando já está montada por Law em \textit{After Method}, e de certa maneira, sustenta a proposta do \textit{method assemblage}. Law mobiliza de Haraway os conceitos de conhecimento situado e a figuração do ciborgue como conjunto de conexões parciais, e adota a perspectiva do \textit{material semiotics} como descrição alternativa da teoria ator-rede que ele pratica \parencite[pp.~67--84, e glossário, pp.~159--160]{Law2004After}. Parte também da formulação das \textit{conexões parciais}, que Haraway introduz no debate feminista em \enquote{A Cyborg Manifesto} \parencite{Haraway1991Cyborg} e em \enquote{Situated Knowledges} \parencite{Haraway1991Situated} e que Strathern desenvolve em \textcite{strathern2004partial} por meio da linguagem do ciborgue de Haraway para formular um problema antropológico que ela vinha trabalhando sob o nome de \textit{partibility}, isto é, a lógica social que faz com que entidades se incluam mutuamente sem se reduzir umas às outras \parencite[pp.~xxix, 35--38]{strathern2004partial}. Diante do impasse entre singularidade e pluralismo que estrutura o capítulo \enquote{Multiple worlds}, Law toma a noção de \textit{conexões parciais} como ferramenta para pensar relações que se sobrepõem sem se totalizar, e a estende, em chave estratherniana de identidades plurais que se incluem sem se reduzir, ao conceito de \textit{enactment} de Mol, em coautoria intelectual sustentada que atravessa pelo menos três capítulos do livro \parencite[caps.~3, 4 e 8]{Law2004After}.

A figuração da costura que desenvolvi acima é, portanto, um desdobramento, no registro têxtil-figurativo, da articulação que Law constrói no registro ontológico-metodológico, e essa antecedência inscreve a operação em uma tradição teórica feminista de longa data, em que a antropologia de Strathern, a tecnociência feminista de Haraway, a praxiografia de Mol e o \textit{method assemblage} de Law trabalham juntos para pensar relações que não se totalizam.

Junto aos cortes discutidos até aqui, Law elabora uma outra forma de \textit{pensar} que importa para o que esta tese faz nas subseções e capítulos seguintes. Em \textit{After Method}, o que venho chamando de \textit{corte} é escrito como momento de uma operação mais ampla, que Law denomina \textit{method assemblage} (montagem metodológica), e que ele caracteriza pelos verbos \textit{crafting} (artesanar), \textit{bundling} (atar em feixe) e \textit{gathering} (ajuntar). Antes desse momento há o \textit{hinterland}, retomado da subseção anterior, que se estende para muito além das bancadas e dos instrumentos do laboratório, abrangendo o conhecimento tácito, software de computador, habilidades linguísticas, capacidades de gestão, sistemas de transporte e comunicação, salários, fluxos de financiamento, prioridades de agências de fomento, agendas políticas e econômicas \parencite[pp.~40--41]{Law2004After}. A lista pode parecer interminável, e suas fronteiras permanecem porosas e extensíveis em todas as direções. Daí decorre que descrever uma pesquisa de modo completo é impossível, e o que se faz são montagens que produzem, ao mesmo tempo, presença, ausência manifesta e \textit{otherness}.\footnote{O \textit{hinterland} e o \textit{corte} são conceitos complementares e distintos. O corte é a operação ativa que delimita o que entra e o que sai da análise; o \textit{hinterland} é o feixe de relações literárias e materiais que torna essa operação possível antes mesmo de ser realizada, isto é, o conjunto de inscrições, instrumentos, tradições disciplinares, infraestruturas materiais e suposições já sedimentadas das quais qualquer corte se vale para produzir o que aparece. O conceito de \textit{corte} foi desenvolvido acima a partir de \textcite{strathern2011cortando}; o de \textit{hinterland}, retomado aqui, retorna nos capítulos~\ref{capitulo3} e~\ref{capitulo4}. Convém anotar uma diferença lexical entre Law e a leitura que aqui proponho. Law não usa \enquote{corte} como categoria conceitual, ele descreve a delimitação como \textit{crafting}, \textit{bundling} e \textit{gathering} de relações entre presença, ausência manifesta e \textit{otherness} \parencite[p.~14, pp.~83--85, e glossário, p.~160]{Law2004After}. Quando enumera, na conclusão de \textit{After Method}, as metáforas que sustenta para imaginar o método, ele elenca \enquote{Localities. Specificities. Enactments. Multiplicities. Fractionalities. Goods. Resonances. Gatherings. Forms of craftings. Processes of weaving. Spirals. Vortices} \parencite[p.~155]{Law2004After}: o substantivo \textit{cut} e seus cognatos no sentido conceitual estão ausentes dessa lista programática. Mas, apesar de não usar o termo, o livro mobiliza Strathern recorrentemente como interlocutora, inclusive em uma passagem decisiva em que a formulação das \textit{conexões parciais} é apresentada como articulação Haraway--Strathern: \enquote{Feminist technoscience writer Donna Haraway, and following her, anthropologist Marilyn Strathern, talk of partial connections} \parencite[p.~64]{Law2004After}. Articulo o vocabulário de Law ao gesto de corte de Strathern por afinidade prática, pois, ambos descrevem a delimitação como condição da análise, e essa afinidade já está operada em Law sob outro vocabulário. A passagem do registro ontológico-metodológico de Law (presença, ausência manifesta, \textit{otherness}, \textit{hinterland}) ao registro têxtil-figurativo desta tese (cortar, costurar, \textit{patchwork}, \textit{crazy quilt}) é desdobramento dessa articulação. Law também recorre ao léxico têxtil quando cita Derrida sobre o \textit{sheaf} como \enquote{complex structure of a weaving, an interlacing which permits the different threads and different lines of meaning--or of force--to go off again in different directions} \parencite[p.~41]{Law2004After} e descreve o \textit{crafting} do método como \enquote{processes of weaving} \parencite[pp.~144--145, 155]{Law2004After}.}

A palavra que melhor traduz \textit{hinterland} para o português, embora historicamente problemática nas ciências sociais, é \textit{contexto}. A desconfiança é justificada, e \textcite{Latour2012Reagregando} a formula com clareza ao recusar a posição que faz do social um \enquote{contexto} no qual atividades não sociais aconteceriam, e que toma o social como domínio específico da realidade ou tipo de causalidade explicativa. Para Latour, o uso preguiçoso de \textit{contexto} reduz a palavra a pano de fundo passivo, a entorno explicativo, ao social que estaria fora da ciência e que as ciências sociais teriam por tarefa explicar. A crítica é precisa e atravessa a figuração têxtil que esta seção mobiliza, da qual o próprio Latour participa com o vocabulário do \textit{netting, lacing, weaving, twisting} de laços fracos \parencite[p.~3]{latour1996clarifications}. A recuperação que proponho aqui parte da crítica de Latour para recuperar a palavra como nomeação de associações ativas, e não como um rótulo de domínio passivo. A etimologia de \textit{contexto} oferece a via dessa recuperação. \textit{Contexto} vem do latim \textit{contextus}, particípio passado de \textit{contexere}, e o verbo é composto por \textit{con-} (junto) e \textit{texere} (tecer): contexto é, etimologicamente, o que se tece junto.\footnote{A formulação \textit{tecer junto} consta nos principais dicionários etimológicos da língua portuguesa. \textcite[verbete \enquote{contexto}]{Cunha2010Etimologico} registra a procedência do latim \textit{contextus}, particípio passado de \textit{contexere}, com o sentido de \enquote{tecer junto, entrelaçar}. \textcite[verbete \enquote{contexto}]{Houaiss2014Dicionario} desenvolve a definição lexicográfica como \enquote{relação de dependência entre as situações que estão ligadas a um fato ou circunstância} e como \enquote{reunião dos elementos do texto que estão relacionados com uma palavra ou frase e contribuem para a modificação ou esclarecimento de seus significados}. A definição de Houaiss é, ela mesma, sintomática do potencial conceitual da palavra: contexto não é entorno passivo, é \textit{relação de dependência} entre situações, formulação que ressoa com a caracterização strathernina da antropologia como \enquote{estudo de relações com relações} \parencite[p.~33]{Haraway2016Staying}. A recuperação etimológica que faço aqui ressoa com o gesto que \textcite{Ingold2011BeingAlive} já fizera ao retomar de Lefebvre o termo \textit{meshwork} e ao ler a inscrição da atividade humana e não humana no espaço vivido como textura, e não como texto. Em ambos os casos, o gesto é o mesmo, ou seja, recuperar, na palavra, a relacionalidade têxtil que o uso comum apagou. O caso de \textit{contexto} é particularmente importante porque a palavra é constantemente usada nas ciências sociais sem que sua origem têxtil seja lembrada, e a etimologia oferece uma maneira de fazer trabalhar, no português, a figuração que Law mobiliza no inglês com \textit{hinterland}.} A palavra carrega, em sua origem, a relacionalidade ativa que o uso comum esvaziou, e ressoa com a figuração têxtil que esta seção mobiliza. Pensar o \textit{hinterland} como aquilo que se tece junto com a pesquisa, e não como pano de fundo passivo, alinha o vocabulário de Law com a oposição urdidura/trama desenvolvida acima e com a constelação têxtil que atravessa Strathern, Haraway, Latour e Costa. O \textit{hinterland} é a urdidura tecida junto que sustenta o tecido sem aparecer; é o contexto que tece a pesquisa enquanto ela se faz, e que recua da visibilidade quando o produto se apresenta como pronto. Recuperar a palavra \textit{contexto} pela sua etimologia é um meio de devolver à palavra a agência que o uso comum lhe retirou, de modo que, o que está tecido junto com a pesquisa é trama ativa que produz a pesquisa.

Law identifica três modos pelos quais a \textit{otherness} funciona \parencite[p.~84]{Law2004After}, cada um com um grau distinto de intencionalidade. O primeiro é a \textit{rotina} (\textit{routine}), também nomeada na tradição dos estudos sociais da ciência como \textit{black-boxing} \parencite[p.~169, n.~78]{Law2004After}: arranjos que funcionam de modo tão automático que deixam de ser notados, como o fornecimento de eletricidade que sustenta o laboratório ou os protocolos institucionais que tornam uma pesquisa financiável. A intencionalidade aqui é sedimentada, alguém decidiu em algum momento e a decisão se institucionalizou de tal modo que ninguém mais a questiona. O segundo é a \textit{insignificância} (\textit{insignificance}), em que elementos são considerados desinteressantes ou irrelevantes para a prática em curso; aqui a intencionalidade é ativa e é exercida no corte que distingue o que importa do que pode ficar fora. O terceiro é a \textit{repressão} (\textit{repression}), em que elementos não podem ser sustentados ao lado da realidade que se quer construir, e precisam ser ativamente apagados para que essa realidade se mantenha; a intencionalidade é silenciada, porque o quadro analítico não consegue sequer nomeá-los sem se transformar. Law sintetiza esses três modos na passagem em que examina como a \textit{otherness} desaparece da descrição:

\begin{citacaoabnt}
Perhaps it disappears because it is not interesting while it goes on routinely (the power-supply to the Salk Laboratory? The pay cheques to hospital staff? The current organisation of health care? A factory in Harare with a non-proprietorial proprietor?). Perhaps it disappears because it is not interesting, full stop (the availability of alcohol, a pub-going culture, a broken marriage). Perhaps (though no doubt this is an overlapping category) it disappears because what is being brought to presence and manifest absence cannot be sustained unless it is Othered (the social-and-cultural context for alcoholic liver disease in the context of the consultant's office? The active character of authorship in the production of Salk Laboratory statements?). The implication is that Otherness takes a variety of forms. Those above --- routine, insignificance and repression --- are no doubt only three of the possibilities \parencite[p.~84]{Law2004After}.
\end{citacaoabnt}

Os três modos atravessam esta tese. A \textit{rotina} opera nos contratos, infraestruturas e políticas de fomento que o capítulo~\ref{capitulo3} descreve como caixas-pretas reabertas pela dissolução da parceria IBM-C4AI, em que arranjos sedimentados ao longo de cinco anos voltam à descrição etnográfica no momento em que se desfazem. A \textit{insignificância} opera nos cortes que delimitam o que tornar presente em cada capítulo, e nomeia a vastíssima porção do C4AI que esta tese reconhece existir e deixa fora da descrição por escolha analítica, como as práticas técnicas dos demais grupos de pesquisa e as trajetórias dos demais pesquisadores do Centro. A \textit{repressão} opera no que o quadro analítico não acomoda sem se transformar, e retorna no capítulo~\ref{capitulo4} a propósito do conjunto de dados do Spira, em que os corpos dos pacientes hospitalizados, suas dores e suas mortes, comparecem como condição de possibilidade do sistema e excedem o enquadramento computacional pelo qual o Spira os converteu em \textit{features} acústicas (os recursos desenvolvido para o sistema).

A tripartição da \textit{otherness} de Law (rotina, insignificância, repressão) mantém ressonâncias com os quatro vocabulários do corte mobilizados acima (Strathern, Barad, Haraway, Costa) e, ao mesmo tempo, desloca o registro em que esses vocabulários funcionam. Os quatro sustentam, em chaves distintas, a discussão de como e de por que se corta. \textcite{strathern2011cortando} formula o corte como decisão interpretativa do pesquisador e, simultaneamente, como operação que as próprias práticas observadas realizam, exemplificada na propriedade que separa quem pertence e quem não pertence a um fragmento de conhecimento. \textcite{Barad2007Meeting} radicaliza essa distribuição com o conceito de corte agencial, em que o corte se torna encenação de um arranjo material do qual o pesquisador faz parte, e em que a intra-ação constitui simultaneamente o objeto observado e a agência da observação. \textcite{Haraway2023Ficar} acrescenta a essa articulação a dimensão da respons-habilidade, a habilidade de responder pelas diferenças que os cortes produzem no mundo, dado que ninguém vê de lugar nenhum e que nenhuma prática de pesquisa é inocente. \textcite{costa2024manchando} mobiliza essa articulação ao construir a armadilha como modo situado de operar cortes que escondem e mostram seletivamente, e que assumem a parcialidade da posição do pesquisador. Law opera em outro registro descritivo e nomeia onde cada corte coloca cada elemento, ou seja, o que se torna presente, o que fica em ausência manifesta, o que é deslocado para a otherness e o que opera no hinterland sem ser tematizado. Os dois registros são complementares e operam juntos nesta tese: Strathern, Barad, Haraway e Costa descrevem o gesto crítico-epistemológico e ético-político de cortar, e nomeiam a partir de qual posição material e com que responsabilidade o corte é feito; Law descreve a topografia descritiva do corte e nomeia o que ficou em cada posição do arranjo metodológico depois que o corte aconteceu. Resta a pergunta sobre o estatuto ontológico do objeto que aparece de cada corte, ou seja, sobre o que cada corte performa como realidade descrita, e essa pergunta é desenvolvida na subseção~\ref{subsec:cortes_como_enactments} a partir da praxiografia de \textcite{mol2002body}.

Esses cinco vocabulários (Strathern, Barad, Haraway, Costa e Law) convergem nos modos de fazer pesquisa em que decisões da pesquisadora, materiais e tradições disciplinares funcionam juntos. Law ressalva que falar em \textit{crafting} (artesania) não implica afirmar agência humana exclusiva, pois pessoas, máquinas, traços e recursos de toda ordem participam do processo \parencite[pp.~41--42]{Law2004After}, formulação que distribui a agência por toda a trama no mesmo registro em que Strathern, Barad, Haraway e Costa distribuíam o corte. O artesanato do método articula as decisões de quem pesquisa com efeitos que excedem a sua intenção, e a responsabilidade pelos resultados dessa articulação é parte do que define a posição de quem pesquisa na rede que descreve. Cortar, distribuir agência, assumir a parcialidade da posição, responder pelas diferenças que o corte produz, nomear o que ficou em cada posição do arranjo, o que esses cinco modos descrevem como gesto e como topografia tornou-se, ao longo da escrita desta tese, uma prática concreta em pontos específicos do texto.

O suporte material desse artesanato, no caso desta tese, podem ser vistos nas notas de rodapé. Foi observando o que elas faziam ao longo da escrita que o vocabulário de \textcite{strathern2011cortando} e de \textcite{Law2004After} foi posto em relação com a minha prática de pesquisa. A cada investida analítica, eu fazia associações com materiais, autores e experiências diversos, alguns parte do meu repertório e outros não, e essas associações pediam algum tipo de registro. Registrá-las no corpo do texto interromperia o argumento, e optar por não escrevê-las apagaria as associações/conexões/relações realizadas. Por isso, as notas de rodapé desta tese funcionam nesse impasse/tensão. Cada nota acompanha uma associação, um vestígio, um fragmento que o argumento precisou deixar no \textit{hinterland} (e por vezes na \textit{otherness}) para manter seu fluxo e/ou lógica narrativa, e que esses modos de pensamento sobre como fazer pesquisa postos em relação me impediam de trataras notas como complementar ou suplementar à própria pesquisa.

No capítulo~\ref{capitulo1}, as notas funcionam como um laboratório de experimentação das minhas decisões analíticas, e podem ser entendidas como formas situadas do que \textcite{Law2004After} chama de \textit{mess} do método; a parte da pesquisa que escorrega, que muda de forma, que aparece e desaparece, e que a pesquisa depois de \enquote{limpa} parece tornar clara e coerente, ou seja, \enquote{objetiva}. As notas acolhem, ramificam, estendem e apontam o que o argumento principal precisaria deixar de fora para circular, e o fazem de forma heterogênea. Funcionam também como amostras, evidências, vestígios de argumentos e de feituras de pesquisa, e são, assim, gestos que interrompem e cortam redes para que a descrição (e o próprio pensamento) existam e continuem.

As notas funcionam ainda como registro consciente de conflitos, tensões e disputas, sobre controvérsias tecnocientíficas, sociotécnicas, teórico-metodológicas e conceituais, e estendem o campo para fora do texto e do laboratório, como na nota sobre o botão que o pesquisador Digiampietri precisou apertar no servidor do GrIA;\footnote{Escrevendo esse trecho me dei conta que não expliquei o que é o GrIA (Grupo de Inteligência Artificial da Universidade de São Paulo, na Escola de Artes, Ciências e Humanidades), que existe desde agosto de 2006 e realiza pesquisas na área de IA com pesquisadores ligados aos cursos de Sistemas de Informação, trabalhando com temas como reconhecimento de padrões e processamento de linguagem natural. Para mais detalhes do grupo, acesse: \url{www.each.usp.br/gria/}. Acesso: 06 de maio, 2026. Aqui abri a rede para outro grupo de pesquisa em IA dentro da USP, com outras motivações e interesses, e que, apesar de não fazer parte do C4AI, compartilha a mesma infraestrutura computacional, e que eu poderia ou não escolher incluir na minha análise. Decidi parar por aqui e deixar registrado que a infraestrutura de IA do C4AI estende ainda mais a rede sociotécnica. Além disso, \enquote{sigo} Digiampietri por outros motivos que também se conectam a esta pesquisa. Em certo momento da aprendizagem sobre IA com os pesquisadores do C4AI, passei a procurar cursos e aulas no YouTube sobre diferentes temas, entre eles algoritmos, estrutura de dados e análise de redes. Antes disso, eu cursava uma especialização em processos didático-pedagógicos para educação a distância pela Universidade Virtual do Estado de São Paulo (Univesp) e, como parte dessa formação, trabalhei com uma turma de Sistemas de Informação. Decidi assistir às aulas de alguns professores do curso que estão disponíveis no YouTube, e um deles era o professor Luciano Digiampietri. O canal dele é este: \url{www.youtube.com/@digiampietri}. E assim passei a segui-lo nas redes sociais. Algum tempo depois, vi essa postagem dele e descobri a sua relação com o C4AI. Agradeço a ele por ter explicado o contexto de sua postagem em mensagem privada no Instagram e por ter permitido incluí-la nesta tese. Viram como isso funciona? E parece que aqui me sinto mais à vontade para \enquote{conversar com vocês}.} as notas sobre a infraestrutura técnica que sustenta determinados sistemas computacionais citados ao longo desta tese, e as notas sobre as formações em ciência de dados e inteligência artificial realizadas paralelamente ao doutorado. Descrevem também a procedência de certos gestos experimentais, como na nota sobre o \textit{script} em Python que produzi junto a \texttt{claude.ai}para fazer uma análise de rede textual dos capítulos, apresentado na subseção~\ref{subsec:tecnoetnografia}.\footnote{Análise que eu poderia ter feito diretamente com \textit{softwares} de análise textual como o Gephi e o VOSviewer, cujas bases de funcionamento aprendi durante a realização desta etnografia e da entrega do trabalho de conclusão de curso coletivo que fiz junto a colegas do curso de especialização didático-pedagógica para a educação a distância, em que trabalhei com uma cientista da computação, um químico e um assistente social; para poupar trabalho extra, automatizei o processo junto a \texttt{claude.ai}dentro da plataforma do GitHub.} Essas notas funcionam de modos diferentes, e estão atreladas à heterogeneidade de materiais e à análise de cada capítulo, isto é, fazem parte de minhas práticas tecnoetnográficas (que discuto com mais detalhes na subseção~\ref{subsec:tecnoetnografia}) e de meus modos de pensamento. Juntas, produzem um tipo particular de conhecer e expressar o campo de pesquisa, e operam como chaves para abrir (ou fechar) novas associações/conexões/relações. 

A discussão organizada pela expressão \enquote{cortando redes} permitiu-me identificar uma série de cortes que definem o que está presente, o que fica em ausência manifesta e o que é relegado à \textit{otherness} em cada capítulo. A tríade presença, ausência manifesta e \textit{otherness} nomeia o que cada corte produz; o \textit{hinterland}, por sua vez, nomeia o \textit{contexto} de inscrições, instrumentos, tradições disciplinares, infraestruturas e arranjos institucionais sedimentados que torna esses cortes possíveis antes mesmo que sejam realizados \parencite[pp.~32--34, 40--41]{Law2004After}. Esses cortes não foram descritos enquanto estavam sendo feitos, porque eu não estava consciente dessa discussão ao realizá-los, e o que funcionou com eles, antes da percepção, foi precisamente esse \textit{hinterland} que sustenta a posição da pesquisadora e o lugar de onde ela recorta. Foi na revisão da escrita desta tese que esse trabalho se tornou um problema da própria pesquisa, e os cortes que descrevo a seguir são, nesse sentido, necessariamente incompletos, pois, outros estão sendo feitos sem que eu os tenha percebido. O capítulo~\ref{capitulo1} opera faz cortes reflexivos sobre as decisões metodológicas que constituem a pesquisa; o capítulo~\ref{capitulo2}, faz cortes teórico que recruta os aliados conceituais e localiza a tese em um campo em formação; o capítulo~\ref{capitulo3}, faz cortes pela extensão que compõe a rede longa e histórica do C4AI; o capítulo~\ref{capitulo4}, faz cortes pela densidade que segue as práticas técnicas situadas na construção do Spira. Esses cortes são tratados com mais detalhes em subseção própria, pelas quais, descrevo o que está presente, o que está em ausência manifesta, o que fica em \textit{otherness}, e o \textit{hinterland} que \textit{contextualizam} os cortes, assim como, o papel das notas de rodapé e as justificativas dos cortes como um maneira de me responsabilizar pelas escolhas, situadas e parciais, que realizei ao longo desta pesquisa. 

\subsection{Cortar e costurar: \textit{patchwork} etnográfico}
\label{subsec:corte_reflexivo_cap1}

A seção anterior estabeleceu três interpretações da figuração têxtil. A primeira é o \textit{corte}, na leitura combinada de \textcite{strathern2011cortando}, \textcite{Barad2007Meeting}, \textcite{Haraway2023Ficar} e \textcite{costa2024manchando}, a partir das quais o corte surge como gesto que qualquer pesquisa precisa fazer a partir de algum lugar, e que cada uma das autoras problematiza a seu modo (corte interpretativo, corte agencial, respons-habilidade, armadilha situada), distribuindo a decisão de cortar entre quem pesquisa, os materiais e as práticas observadas. A segunda é a \textit{costura}, na leitura em que cortar e costurar trabalham juntos na feitura etnográfica, e a partir da qual \textcite{Law2004After} elabora o \textit{method assemblage}: o que se torna presente, o que se mantém em ausência manifesta, o que recua para a \textit{otherness} e o que sustenta o \textit{hinterland} sem aparecer tematizado. A terceira é o \textit{crazy quilt}, que proponho como leitura situada a partir de duas procedências encadeadas. De um lado, da convergência têxtil que percorre Strathern (conexões parciais, entidades parcialmente conectadas, análogas sem serem homólogas, que se incluem mutuamente sem se reduzir), Haraway (\textit{string figures} e simpoiese, padrões de cordas que se passam de mão em mão sem totalização, sistemas coletivamente produtores em que nada se faz sozinho), Law (\textit{sheaf}, \textit{bundling}, \textit{weaving}, \textit{gathering}, o método como artesanato que feixa, trança e reúne relações heterogêneas sem encerrá-las em um padrão único), Latour (\textit{netting}, \textit{lacing}, \textit{weaving}, \textit{twisting}, laços fracos que se enredam, atam, tecem e torcem uns aos outros até produzirem a resistência da rede), Costa (cesta e armadilha, dispositivos que trançam fibras para capturar seletivamente e que, na captura, fazem vazar o que estava silenciado) e Ingold (\textit{meshwork}, malha de linhas viventes que se entrelaçam pelo movimento de habitar o mundo), conforme discuti na seção anterior. De outro, da experiência etnográfica desta pesquisa, em que minha tese se foi compondo por retalhos heterogêneos.

Cada autora ou autor da convergência têxtil que mencionei têm seus próprios retalhos. O retalho é, aqui, a unidade empírica heterogênea que cada pesquisa recorta de um mundo também heterogêneo e costura em sua análise: pode ser um objeto ritual, uma entidade vivente, uma substância corporal, um artefato técnico, uma inscrição tecnocientífica, um processo clínico ou um evento histórico, e o que reúne essa heterogeneidade é o gesto de corte/recorte que cada um realiza ao tomá-lo como matéria de análise. Strathern (sociedades das Terras Altas da Nova Guiné e suas trocas cerimoniais de porcos e conchas, dote nupcial dos \textit{'are'are} que converte corpo, sopro e imagem em taro, porcos e dinheiro de conchas, patentes médicas que truncam redes), Haraway (pombos-corredores de Los Angeles, recifes de coral e seus holobiomas, ovelhas Navajo-Churro, orquídeas, lêmures, jogos de cama-de-gato, dingos e abutres tomados emprestados de Rose e van Dooren), Costa (sangue menstrual, caleidoscópios fractais, cestas trançadas e armadilhas que encantam e fazem vazar), Ingold (aranhas presas a suas teias, pássaros suspensos no ar, peixes nadando nas correntes), Latour (a dupla hélice do DNA em Cambridge, o computador Eclipse do Instituto Pasteur, o hormônio liberador GHRH em disputa em New Orleans, os microrganismos da fermentação espontânea em Paris, a lâmpada incandescente de Menlo Park, o motor a diesel, a tabela periódica dos elementos químicos, o mapeamento de Sakhalin no Pacífico, as expedições portuguesas rumo às Índias no século XV) e Law (doença hepática alcoólica nos hospitais ingleses, bomba d'água do Zimbábue, avião caça TSR-2 cancelado).

Antes de descrever como esses retalhos se compõem em minha tese, devo marcar por que adoto a figuração têxtil em primeiro lugar. Meu ponto de partida tem precedente teórico em \textcite{LakoffJohnson1980Metaphors}, para os quais, além de suas funções poéticas, as metáforas funcionam como estruturas que organizam pensamento e prática: \enquote{nosso sistema conceitual ordinário, em termos do qual pensamos e agimos, é fundamentalmente metafórico} \parencite[p.~3, tradução minha]{LakoffJohnson1980Metaphors}. Segundo os autores, há diferentes tipos de metáforas, como a metáfora estrutural, a metáfora ontológica, a metáfora orientacional, a personificação e a metonímia, cada uma com função específica. Considerei, em particular, o capítulo da \textit{personificação}: metáforas antropomórficas permitem compreender fenômenos não humanos em termos humanos, e essa compreensão \enquote{dá origem a e justifica} ações específicas sobre o objeto figurado \parencite[p.~34, tradução minha]{LakoffJohnson1980Metaphors}. A ciência da computação e a inteligência artificial fazem uso intensivo desse tipo de metáfora. \textit{Redes neurais}, \textit{neurônios artificiais}, \textit{árvores de decisão}, \textit{florestas aleatórias}, \textit{atenção}, \textit{memória}, \textit{treinamento}, \textit{aprendizado}, \textit{alucinação}, \textit{agente}, \textit{alinhamento}, \textit{embedding}, \textit{peso} e \textit{ativação} são metáforas importadas da neurociência, da biologia, da lógica, da geometria e da ética, que funcionam de modo técnico, mas é como se nomeassem de fato o que descrevem. Esse uso intensivo, e a discussão sistemática das suas implicações epistêmicas e políticas, é objeto do projeto autônomo de pesquisa que indico adiante.

O mesmo movimento de empréstimo entre áreas se faz, dentro da própria antropologia, em particular na tradição feminista dos estudos de ciência e tecnologia que sustenta o argumento desta seção. Haraway constrói a sua proposta de pensamento em \enquote{ficar com o problema} a partir de termos da biologia e da microbiologia: a simpoiese, que retoma de M. Beth Dempster, os holobiomas e a simbiogênese, que retoma de Lynn Margulis, as relações multiespécies que retoma da etologia. Barad toma a intra-ação, a difração e o realismo agencial de Niels Bohr e da física quântica. Strathern toma da matemática, via James Gleick, a poeira de Cantor e a noção de fractal para pensar a antropologia das conexões parciais, e toma da cibernética, via Haraway, a figura do ciborgue. Costa trabalha com caleidoscópios fractais e armadilhas em vocabulário que cruza a geometria, a antropologia da técnica e a etnografia indígena. Ingold pensa o \textit{meshwork} em vocabulário que cruza a arquitetura, a tecelagem e a ecologia comportamental. A diferença entre o empréstimo praticado pela ciência da computação e o empréstimo praticado por esta tradição antropológica está no modo como cada uma se relaciona com a sua condição de empréstimo metafórico, e essa diferença é o que me leva a nomear o gesto antropológico, no vocabulário desta tese, como \textit{figuração} no sentido de Haraway, e a reservar o termo \textit{metáfora} para a operação formal de mapeamento entre domínios que tanto a antropologia quanto a ciência da computação praticam, embora cada uma se relacione de modo distinto com ela. A figuração, tal como eu a entendo nesta tese, é a metáfora reflexivamente assumida, situada e respons-hável pelas diferenças que produz. A metáfora, por sua vez, nomeia a operação formal pela qual um domínio é mapeado sobre outro, sobreposto ou justaposto em problematização. Desenvolvo essa distinção conceitualmente na próxima subseção, com apoio em \textcite{Haraway2016Staying}.

Para \textcite{LakoffJohnson1980Metaphors}, a metáfora é uma estrutura cognitiva que organiza o pensamento e a ação por meio de diferentes mapeamentos de um domínio sobre outro, cuja \enquote{função primária é a compreensão} \parencite[p.~36, tradução minha]{LakoffJohnson1980Metaphors}. Essas metáforas, ao serem incorporadas nas linguagens de uma determinada área do conhecimento, acabam funcionando como se sempre tivessem estado ali. Para \textcite{Haraway2016Staying}, a metáfora passa a ser figuração quando também produz mundos, e inseparável a essa produção, consequências e responsabilidades. Por isso, é um \enquote{tropo teórico, modo de pensar-com} e \enquote{o padrão e a montagem que solicita resposta} \parencite[p.~3, tradução minha]{Haraway2016Staying}. Desse modo, a metáfora descreve o gesto formal pelo qual termos viajam entre domínios; a figuração nomeia esse mesmo gesto quando ele é praticado de forma situada, assumido como tropo material-semiótico e explicado. O termo \textit{tropo}, com o qual Haraway nomeia explicitamente as figurações que pratica, vem da retórica clássica (do grego \textit{tropos}, viragem) e funciona, no vocabulário desta tese, como categoria mais ampla que compreende a metáfora e a figuração, ambas, tropos, que funcionam de forma diferente conforme a postura assumida por quem os pratica. A figuração têxtil é, nesse sentido, um gesto duplamente situado: do lado da antropologia, dá continuidade a uma operação descritiva que a tradição feminista dos estudos de ciência e tecnologia sustenta; do lado de meu campo de pesquisa com o C4AI, faz com que as metáforas da inteligência artificial também sejam situadas.\footnote{A análise empírica dos tropos e das metáforas na literatura da inteligência artificial excede o escopo desta tese e desenvolverei em projeto de pesquisa autônomo, posterior à defesa. O desenho desse projeto prevê a produção de um \textit{corpus} a partir de Web of Science e bases complementares, codificação sistemática em três níveis (lexical, tropológico, funcional), análise de distribuição por subcampo disciplinar e leitura próxima de subcorpus selecionado, em diálogo com as tradições de \textcite{LakoffJohnson1980Metaphors} e de \textcite{Haraway2023Ficar} trabalhadas no corpo desta subseção. O argumento da figuração têxtil que estrutura este capítulo será apresentado de forma resumida no VII Simpósio Internacional da Rede Latino-Americana de Estudos sobre Vigilância, Tecnologia e Sociedade (LAVITS~2026), \enquote{Tramar o presente: tecnopolíticas dissidentes e ecologias do Comum}, Rio de Janeiro, 26 a 28 de agosto de 2026, sob o título \enquote{Tramar a inteligência artificial: figurações têxteis para uma etnografia situada}.}

Devo marcar também por que chamo de \textit{retalho} e \textit{costura} aquilo que outras tradições nomeariam como dado, material, objeto, fragmento ou caso. O \textit{dado} carrega a sedimentação positivista que sua etimologia latina (\textit{datum}, o que é dado) preserva, sugerindo que o material da pesquisa preexiste ao gesto que o constitui; o \textit{material} é mais neutro, e descreve a matéria-prima a ser processada, e supõe um produto final em que essa matéria desaparece como matéria; o \textit{objeto} carrega a oposição com sujeito que minha tese se compromete a desfazer; o \textit{fragmento} sugere uma totalidade prévia da qual o pedaço teria se desprendido, e a posição de pesquisa que assumo recusa essa totalidade; o \textit{caso} pressupõe um quadro classificatório que decide previamente o que conta como instância. O \textit{retalho}, em contraste, nomeia uma matéria que já foi cortada de algum lugar antes de chegar a quem a costura, que carrega a marca desse corte em suas bordas, e que se compõe com outros retalhos sem perder essa marca. Pensar com \textit{retalhos} faz com que o recorte do material da pesquisa já me chegue como corte, e o ato de costurá-lo com outros retalhos acentua ainda mais essa instância parcial e situada. A \textit{costura}, por sua vez, é a forma pela qual a pesquisa adquire sua feitura, seu artesanato, de modo a deixar aparentes as suas junções, justaposições e sobreposições espontâneas, sem seguir um modelo, e por consequência é distinta da \textit{síntese} ou da \textit{integração}, que pressupõem o desaparecimento dos cortes e costuras da pesquisa. Retalho e costura formam um par categorial no plano da descrição etnográfica, à figuração têxtil que desenvolvi nas subseções anteriores sobre os cortes etnográficos. A figuração têxtil é a primeira figuração desta tese; uma segunda, a da compostagem, retoma a expressão \enquote{compostagem intelectual} que citei a partir de \textcite[p.~69]{Haraway2023Ficar} e recebe tratamento próprio na Seção~\ref{sec:compostagens}, ao final deste capítulo.
%%%%%%%%%%%%
\subsubsection{O que torno presente}

A feitura desta etnografia, no sentido do \textit{crafting} de Law (fazer-com, compor-com), é o que torno presente neste capítulo: como e por que cheguei ao C4AI, as conexões com os referenciais teóricos dos estudos de ciência e tecnologia cuja maneira de fazer pesquisa entrou em ressonância com o que meu campo exigia, a experimentação com a inteligência artificial (Claude/Anthropic) e com \textit{softwares} de análise de dados, e as compostagens analíticas (os tipos heterogêneos de aprendizagens e experiências) que esta etnografia gerou. Tornar presente a feitura é exibir os gestos pelos quais minha pesquisa se compôs e as alianças teóricas, técnicas e materiais que essa composição costurou.

Alguns desses gestos retornam, ao longo do capítulo, em notas de rodapé em que recupero disputas conceituais sobre o vocabulário em uso (Le Guin sobre o termo \enquote{tecnologia}), estendo a cena de campo para fora do laboratório (o botão do servidor de Digiampietri) e documento a procedência metodológica de gestos experimentais (o \textit{script} em Python que escrevi no estilo do \texttt{InfraNodus}).

Torno presente também minha posição como antropológoca/etnógrafa na rede que descrevo. Ao longo da pesquisa, tornei-me colaboradora do grupo de humanidades do C4AI e membro do Observatório Brasileiro de Inteligência Artificial (OBIA), o que me exigiu gestão ativa da proximidade com o campo. Essa posição ambivalente produziu tensões que incorporei reflexivamente à análise: o acesso a eventos, conversas informais e dinâmicas internas conviveu com o desconforto de documentar vulnerabilidades de uma instituição da qual passei a fazer parte. Como o etnógrafo que Latour descreve \enquote{rondando por ali, inútil, de braços cruzados} \parencite[p.~62]{Latour2017Esperanca}, experimentei o desajuste de não compreender inicialmente aquele mundo e transformei esse desconforto em recurso metodológico. Minha presença no campo é, ela mesma, um dos cortes que torna presentes certas relações e deixa outras ausentes.

Ao longo do trabalho de campo, realizei entrevistas formais com os coordenadores dos cinco grupos de pesquisa do C4AI: Processamento de Linguagem Natural, Computação Bio-Inspirada e Aprendizado de Máquina, Saúde, Agronegócio e Biologia, AI Humanities. Os cortes analíticos que estruturam minha tese demandaram seleção estratégica: escolhi as entrevistas com Fábio (diretor do C4AI), Cláudio (vice-diretor executivo pela IBM) e Marcelo (coordenador do grupo de NLP e responsável pelos projetos Spira e Corpus Carolina) porque suas narrativas permitiam performar os dois \textit{enactments} complementares do Centro que documento nos capítulos~\ref{capitulo3} e~\ref{capitulo4}. As demais entrevistas e os demais grupos não desaparecem do tecido desta tese, pois retornam adiante como ausências manifestas, ou seja, como retalhos que reconheço existir e decido não compor na cena central.

%%%%%%%%%%%%

%%%%%%%%%%%%
\subsubsection{Ausências manifestas}

As ausências manifestas deste capítulo são os retalhos que reconheço existir e decido não compor na cena descrita. São quatro, e cada uma desapareceria sob a coerência polida que \textcite{Law2004After} identifica como o gesto pelo qual a \textit{mess} do método é apagada em nome de uma narrativa limpa.

A primeira ausência manifesta deste capítulo é meu caderno de campo bagunçado, cerca de cem páginas digitadas no aplicativo de notas do celular durante a observação presencial e remota, no qual anotava qualquer coisa que me chamasse a atenção durante o trabalho de campo. Voltei a essas anotações muitas vezes ao longo da escrita para confirmar lembranças de eventos e situações, e na maior parte dos casos isso funcionou. A maneira acelerada como as coisas aconteciam e a pressa em digitá-las deixaram marcas claras no material; mantinha os olhos no entorno enquanto digitava no celular. Esse caderno alimentou minha narrativa sem se inscrever nela.

A segunda ausência manifesta são as conversas que tive em campo, a maioria das quais existe apenas em minha memória. Algumas ficaram fora da descrição por escolha argumentativa, outras por razões éticas. As primeiras alimentaram minha compostagem analítica como rastro de uma sociabilidade que decidi deixar implícita no argumento; as segundas permanecem fora do texto como gesto de cuidado em relação aos meus interlocutores.

A terceira ausência manifesta são as muitas versões do capítulo que escrevi, revisei e abandonei ao longo do percurso. À medida que eu escrevia, revisava e refletia, minha análise se transformava, e preservei todas essas versões. Os arquivos constituem rastro material do processo de construção desta etnografia, do qual o capítulo presente é herdeiro e, em parte, herdade: evidenciam as escolhas analíticas, estilísticas e bibliográficas que mantive ou descartei.

A quarta ausência manifesta são as entrevistas com os coordenadores dos demais grupos de pesquisa do C4AI (Computação Bio-Inspirada e Aprendizado de Máquina, Saúde, Agronegócio e Biologia, AI Humanities). Essas entrevistas informaram minha compreensão geral do campo, e permanecem na ausência manifesta porque o recorte argumentativo desta tese privilegiou os dois \textit{enactments} complementares que trabalho nos capítulos~\ref{capitulo3} e~\ref{capitulo4}.
%%%%%%%%%%%%
\subsubsection{Hinterland e \textit{otherness}}

O \textit{hinterland} e a \textit{otherness} deste capítulo nomeiam dois modos distintos pelos quais retalhos se retraem da composição central. O \textit{hinterland} se retrai como urdidura, ou seja, sustenta o tecido sem aparecer na superfície; a \textit{otherness} se retrai como o que não pode ser costurado sem desfazer o tecido que se compôs.

O \textit{hinterland} deste capítulo, no sentido em que Law trabalha o termo, é o conjunto de condições que sustentam minha etnografia recolhendo-se da descrição. Inscrevo nele, em primeiro lugar, as experiências profissionais e formativas que precederam meu doutorado, condição e condicionantes desta pesquisa. Entre elas, minhas formações em Ciência de Dados constituem capital técnico cuja presença pressuposta atravessa as decisões metodológicas deste capítulo e aparece, nas notas de rodapé, como procedência. Essas trajetórias entram em minha escrita como leitura, escolha de literatura, inclinação analítica; sua condição de cena narrada está fora do enquadramento deste capítulo.

Em segundo lugar, inscrevo no \textit{hinterland} as condições materiais cotidianas que sustentaram a feitura desta tese: a bolsa de doutorado da Capes, o apoio familiar e o apoio de amigos. Essas condições recuam da descrição etnográfica e, ao mesmo tempo, sustentam minha existência de etnógrafa como figura analítica. Tornar essas condições explícitas aqui é exigência da posição reflexiva que assumo neste capítulo: nenhuma prática de pesquisa é inocente \parencite{Haraway2023Ficar}, e a infraestrutura material e afetiva que a sustenta faz parte da rede que ela descreve.

A \textit{otherness} deste capítulo sou eu enquanto pessoa fora do que descrevo: meus afetos não convertidos em argumento, minha vida pessoal mantida fora do enquadramento, hábitos cotidianos sedimentados que escapam à descrição. A \textit{otherness}, em Law, se faz por três modos que ele formula como rotina, insignificância e repressão \parencite[p.~84]{Law2004After}; no corte reflexivo deste capítulo, o modo dominante é a repressão, aquilo cuja entrada plena não pode ser sustentada ao lado da prática em curso sem que essa prática se reconfigure. A entrada plena de mim enquanto pessoa total transformaria minha etnografia em outra forma textual, deslocando o tom analítico para o autobiográfico ou memorialístico. A diferença em relação ao \textit{hinterland} está no estatuto da retração: o \textit{hinterland} pode ser nomeado e inscrito no texto, como faço acima ao explicitar a bolsa, a formação e o apoio; a \textit{otherness} recua estruturalmente, e nomeá-la em totalidade transformaria a prática que ela sustenta.

Tornar explícitas estas quatro dimensões (presença, ausência manifesta, \textit{otherness}, \textit{hinterland}) é o trabalho do corte reflexivo. A pesquisa em inteligência artificial me exige um conjunto de decisões metodológicas que precedem a descrição da rede e a constituem como possível. Tornar presentes essas decisões, em vez de pressupô-las, é o que torna minha etnografia auditável e disponível ao escrutínio. O corte reflexivo dá a ver que a unidade da prática etnográfica reside nos modos de fazer pesquisa, ou seja, na maneira como costuro como etnografia materiais heterogêneos (cadernos digitalizados, conversas memorizadas, versões abandonadas, formações técnicas, condições materiais, entrevistas selecionadas e descartadas), com costuras aparentes que decidi deixar à mostra em vez de polir.

% =====================================================================
% capítulo1, Seções finais revisadas (versão consolidada)
% Existências parciais → Compostagens → Conclusões parciais
% =====================================================================

\section{\enquote{Existências parciais}: quase-máquinas, híbridos e a composição com IA}
\label{sec:existencias-parciais}

\epigrafe{Direct access to anything is not in the power of \enquote{humans} souls and machinery. Modems and faxes add just a few more artifacts in the midst of all of those that already compose \enquote{natural} interactions. [\ldots] We all have only partial existence.}
{Bruno Latour e Richard Powers, \textit{Two Writers Face One Turing Test: A Dialogue in Honor of HAL}, 1998.}

As polarizações e multiplicidades descritas nas páginas anteriores me exigem ferramental teórico capaz de compreender como pessoas e sistemas de IA se compõem sem reduzir essa composição a categorias prévias (usuário/ferramenta, humano/máquina, autêntico/artificial). Os conceitos de existência parcial, conexão parcial e agência distribuída me oferecem esse ferramental.

\subsection*{Do princípio de simetria às existências parciais}

Proponho aqui uma leitura da formulação \enquote{existências parciais} como deslocamento teórico interno à própria obra de Latour, em lugar de ajuste vocabular. O Latour de \textit{Ciência em Ação} \parencite*{Latour2011CienciaEmAcao} e de \textit{Jamais Fomos Modernos} \parencite*{Latour1994Jamais} trabalha com o princípio de simetria generalizada: tratar humanos e não humanos nos mesmos termos. Esse passo é decisivo porque desmonta a pressuposição de que apenas um dos lados da relação merece descrição. O vocabulário da simetria, contudo, ainda pressupõe dois polos a serem simetrizados, e o gesto contesta a hierarquia sem dissolver a própria partição.

O diagnóstico dos híbridos é persistente na obra de Latour: dos quase-objetos de \textit{Jamais Fomos Modernos} ao projeto Aramis, dos híbridos de \textit{A Esperança de Pandora} \parencite*{Latour2017Esperanca} à crise da crítica em \textit{On the Modern Cult of the Factish Gods} \parencite*{latour2010factish}. As misturas funcionam como a condição mesma do mundo que habitamos, em lugar de categoria intermediária a ser superada. No diálogo com Powers de 1998, a formulação se radicaliza: \enquote{todos temos apenas existência parcial} \parencite{Latour2008Powers}. Em lugar de dois polos simetrizados ou entidades definidas pelo que não são, o que se descreve são existências que nunca foram completas para começar. E na \textit{Investigação sobre os Modos de Existência} \parencite*{latour2019investigacao}, o próprio conceito de ator-rede é reconhecido como insuficiente, apenas um modo entre outros.

Esse movimento que identifico na obra de Latour surgiu como consequência do meu próprio percurso em campo, em lugar de revisão bibliográfica. Foi acompanhando pesquisadores de IA no C4AI, observando como seus algoritmos, \textit{datasets} e modelos eram simultaneamente produtos de escolhas humanas e entidades com comportamentos que surpreendiam seus próprios criadores, que a insuficiência do vocabulário simétrico surgiu para mim. A inteligência artificial generativa exemplifica de modo contundente o que significa existência parcial: é composta por milhões de decisões humanas sedimentadas e, no entanto, produz resultados que nenhum de seus criadores previu ou controla inteiramente; trabalha sobre conteúdo simbólico e, no entanto, não compreende no sentido humano do termo; é radicalmente dependente de infraestruturas institucionais e, no entanto, age como interlocutor com certa autonomia funcional. Escapa à descrição como humano ou como não humano: aparece como existência parcial de modo brutal, mais à vista do que qualquer vieira ou micróbio porque atua justamente no terreno que a tradição ocidental reservou como marca distintiva do humano, ou seja, linguagem, pensamento, escrita.

\subsection*{Quase-máquinas e agência distribuída}

Sistemas de IA generativa como ChatGPT e \texttt{claude.ai}são produto de escolhas humanas em cada etapa de sua construção, da curadoria dos dados de treinamento à definição de arquiteturas de redes neurais, do ajuste fino com avaliadores humanos às decisões sobre o que o modelo deve ou não fazer.\footnote{O processo de construção de um modelo de linguagem envolve camadas sucessivas de decisões humanas que permanecem opacas ao usuário final: a seleção e filtragem dos corpora de treinamento, a definição da arquitetura e dos hiperparâmetros da rede neural, o ajuste fino por meio de \textit{feedback} humano (\textit{reinforcement learning from human feedback}, RLHF), e as políticas de segurança e alinhamento. Cada uma dessas etapas incorpora pressupostos culturais, vieses e escolhas normativas que condicionam os resultados. Sobre os desafios éticos e políticos dessas escolhas, ver \textcite{Buolamwini2023Unmasking}, \textcite{Noble2018Algorithms}, \textcite{Benjamin2019RaceAfter} e \textcite{Christian2020Alignment}.} São, portanto, objetos técnicos saturados de significação. A IA generativa, em contraste com ferramentas anteriores como processadores de texto ou buscadores acadêmicos, efetivamente trabalha sobre o conteúdo: sugere reformulações, identifica inconsistências, propõe estruturas argumentativas. Essa mediação é qualitativamente distinta, e a sua especificidade se perde quando reduzida ao binário entre feito por humano e feito por não humano.

A partição entre o que foi feito por humanos e o que foi gerado por IA pressupõe a existência de entidades puramente humanas de um lado e puramente maquínicas do outro, separadas por uma fronteira estável e visível. Essa partição é exatamente o que precisa ser explicado, em lugar daquilo de onde se parte. Em um diálogo que dramatiza essa questão, Latour propõe a Powers que abandonem o vocabulário da \enquote{máquina} em favor do de \enquote{quase-máquina}: entidades abertas, conectadas, que jamais cortaram o cordão umbilical com as instituições, os corpos e os coletivos que as sustentam. A pergunta pertinente desloca-se de \enquote{a quem imaginamos estar falando quando falamos com nossas máquinas?} para \enquote{através de quais máquinas falamos? Quais máquinas nos permitem falar, para início de conversa?}, porque mesmo quando estamos sozinhos, na carne, falamos através de toda sorte de máquinas: fonação, cordas vocais, língua portuguesa, convenções estilísticas \parencite{Latour2008Powers}.

Essa discussão me ajuda a compreender por que o debate sobre inteligência artificial permanece preso a dualismos que, embora teoricamente superados, são consistentemente reafirmados na prática. A pergunta \enquote{podem as máquinas pensar?} foi reformulada por Alan Turing como problema prático, o jogo da imitação, precisamente porque a formulação original carregava pressupostos que ele considerava improdutivos.\footnote{Alan Turing utilizava ironia e sátira como ferramentas retóricas para desafiar o antropocentrismo prevalecente. Seu jogo da imitação emprega metáforas metodologicamente: máquina-como-mulher, máquina-como-humano funcionam como dispositivos conceituais para desnaturalizar pressupostos sobre inteligência. A imitação, para Turing, funcionava mais como princípio matemático do que como identidade ontológica: máquinas poderiam imitar comportamento inteligente permanecendo ontologicamente distintas de humanos \parencite[p.~5]{Goncalves2024BeautifulThought}. \textcite{Quattrociocchi2025Epistemological} sistematizam essa heterogeneidade através de sete \textit{fault lines} epistemológicas estruturais.} Latour questiona o próprio teste de Turing em seus fundamentos: na prática, o teste confronta uma máquina de carne e osso contra outra máquina de carne e osso. Coisas e pessoas estão demasiado entrelaçadas para serem particionadas antes que o teste comece \parencite{Latour2008Powers}.

A objeção de Lady Lovelace, segundo a qual \enquote{a máquina não tem pretensão de originar coisa alguma}, já fora contestada pelo próprio Turing, que afirmava que máquinas o surpreendiam \enquote{com grande frequência} \parencite{Turing1950CMI}. Latour radicaliza a implicação: ninguém domina as quase-máquinas, porque elas são caixas-pretas que ninguém atravessa com o olhar, embora sejam de nossa própria fabricação. O mais inteligente dos programadores perde-se na linha quatro mil de seu programa, e quando uma dúzia de programadores trabalha junta, o programa começa a ter vida própria \parencite{Latour2008Powers}. A crença no domínio, seja de Deus sobre os humanos, dos humanos sobre as máquinas, ou das máquinas sobre os humanos, é precisamente o mito que as quase-máquinas confrontam.

Há, portanto, existências parciais que se compõem mutuamente, em lugar de entidades plenamente humanas de um lado e plenamente maquínicas do outro. A noção de conexão parcial, formulada por Strathern a partir de sua leitura crítica da teoria ator-rede, me oferece aqui uma precisão decisiva: as entidades existem como parcialmente conectadas, análogas sem serem homólogas \parencite{strathern2005partial}, e escapam tanto à conexão completa (que as tornaria a mesma coisa) quanto à separação completa (que tornaria impossível qualquer relação entre elas). A pesquisadora que escreve com assistência de IA se faz como composição parcial, em lugar de soma de duas entidades discretas (humano + máquina): nem eu existo independentemente das mediações técnicas que me atravessam, nem o modelo de linguagem existe independentemente das decisões humanas que o constituíram. Latour chega ao mesmo ponto por outra via: quase-máquinas não existem separadas da humanidade; são a humanidade em outro estado, assim como água, vapor e gelo são estados diferentes da mesma substância \parencite{Latour2008Powers}.

A existência parcial descreve o que as entidades são: nunca completas, sempre dependentes das associações que as constituem. Falta ainda descrever como essas existências parciais agem juntas, por meio de qual mecanismo a composição entre elas produz efeitos no mundo. É aqui que a noção latouriana de agência distribuída se torna indispensável.\footnote{Latour desenvolve o conceito ao longo de vários textos sob formulações como \enquote{distribuição da atividade}, \enquote{compartilhamento da responsabilidade pela ação} e, sobretudo, \textit{faire-faire} (fazer-fazer), sem empregar sistematicamente a expressão \enquote{agência distribuída} como rótulo teórico unificado. As elaborações mais densas encontram-se no capítulo~6 de \textit{A Esperança de Pandora} \parencite{Latour2017Esperanca}, no artigo sobre o dispositivo de fechamento automático de portas \parencite{johnson1988mixing}, e em \textit{On the Modern Cult of the Factish Gods} \parencite{latour2010factish}.} Para Latour, a agência aparece sempre como resultado de uma composição de forças em uma rede, em que a ação é um \textit{faire-faire} distribuído entre diversos mediadores, e nunca como propriedade isolada de um actante.

Latour constrói o argumento através de exemplos que dramatizam a distribuição. No capítulo~6 de \textit{A Esperança de Pandora}, ele recorre ao debate sobre controle de armas para mostrar que nem a pessoa sozinha nem a arma sozinha matam: a responsabilidade pela ação deve ser compartilhada entre os vários actantes envolvidos. Quando um cidadão empunha uma arma, ele se torna um cidadão-arma, ator híbrido cujo objetivo de ação foi transladado pela associação. Ninguém permanece o mesmo: o cidadão que era pacífico tornou-se ameaçador. A arma que era inerte tornou-se letal. O que age é a associação entre eles \parencite{Latour2017Esperanca}. No artigo sobre o dispositivo de fechamento automático de portas, Latour mostra como moralidade e disciplina são delegadas de humanos a não humanos: conhecimento, moralidade, força e sociabilidade funcionam como propriedades de humanos acompanhados por seu séquito de delegados técnicos \parencite{johnson1988mixing}. O conceito de \textit{fatiche}, fusão entre fato e fetiche, radicaliza a implicação: \enquote{aquele que age não é mestre do que faz}, pois outros actantes \enquote{passam à ação} e o ultrapassam \parencite{latour2010factish}.

Essas formulações têm implicação direta para a composição com inteligência artificial que documento nesta tese. O vocabulário latouriano da agência distribuída descreve o mecanismo dessa composição; o vocabulário haraweano da \textit{simpoiese} nomeia sua condição ontológica: nem eu nem o modelo preexistimos, como entidades completas, ao \textit{fazer-com} que nos constitui mutuamente ao longo do processo. A composição altera o que estou tentando fazer, os caminhos argumentativos que percorro, as formulações que se tornam possíveis e aquelas que se fecham. O objetivo da ação é transladado: o que começa como tentativa de reformular um parágrafo pode desembocar em reorganização conceitual de uma seção inteira. A translação descreve o modo como a mediação se faz, e é também o modo como a simpoiese se distingue do simples uso de ferramenta: o fazer-com transforma os termos da relação, e essa transformação alcança mais do que os produtos imediatos do processo.

O paralelo com o fermento de Pasteur é igualmente produtivo. No capítulo~4 de \textit{A Esperança de Pandora}, Latour mostra que Pasteur trabalha para que o fermento aja sozinho, e se o experimento for bem-sucedido, o fermento torna-se ator com autonomia funcional, embora essa autonomia seja efeito da rede de mediações laboratoriais \parencite{Latour2017Esperanca}. Analogamente, eu trabalho (formulando perguntas, fornecendo material etnográfico, avaliando criticamente os resultados) para que o modelo aja, produzindo tensionamentos, reorganizações e formulações que eu então aceito, rejeito ou modifico. Se a composição funciona, o modelo adquire certa autonomia funcional: sugere caminhos que eu não teria percorrido, identifica inconsistências que eu não teria notado, propõe conexões que alteram o argumento. Essa autonomia, como a do fermento de Pasteur, é efeito da rede de mediações que a tornou possível. A agência é distribuída pela cadeia de traduções entre mim, modelo, texto e argumento.

Latour menciona explicitamente a compatibilidade de sua teoria com a obra de Edwin Hutchins sobre cognição distribuída, que descreve como o trabalho do pensamento é externalizado em formas e instrumentos \parencite{latour2010factish}.\footnote{A obra de referência é \textit{Cognition in the Wild}, na qual Hutchins demonstra, a partir de etnografia da navegação marítima, que processos cognitivos complexos não residem dentro de indivíduos e se distribuem entre pessoas, artefatos e ambientes estruturados.} O pensamento que produz meu argumento acadêmico distribui-se entre eu, o modelo, os diagramas, as notas de campo, os textos teóricos que consulto e as inscrições visuais que emergem dessa cadeia. A externalização do pensamento em instrumentos se faz como redistribuição da agência cognitiva pela rede.

A agência distribuída me oferece, portanto, o mecanismo que a existência parcial postula sem desenvolver: se todos temos apenas existência parcial, agimos sempre através de mediadores que transladam nossos objetivos, redistribuem a responsabilidade pela ação e produzem resultados que nenhum actante isolado teria sido capaz de alcançar. É esse \textit{faire-faire} que descreve com precisão o que acontece quando componho com um modelo de linguagem para produzir esta tese: redistribuição da agência por uma cadeia de translações em que cada mediador transforma o que o anterior lhe passou.

\subsection*{Multiplicidades ontológicas e a especificidade de Claude}

A perspectiva de Law sobre multiplicidade ontológica é particularmente produtiva para pensar sistemas de IA generativa. Considere o caso de Claude. \enquote{Claude} é \textit{enacted} como objetos distintos através de práticas irredutíveis entre si.

\textbf{Claude-em-treinamento} existe nas instalações da Anthropic como \textit{dataset} massivo, arquitetura neural (\textit{transformer} com bilhões de parâmetros), sessões de RLHF com avaliadores humanos que julgam pares de respostas, e políticas de \textit{Constitutional AI} que definem valores incorporados ao modelo. Esse \texttt{claude.ai}figura como objeto de decisões sobre o que incluir ou excluir dos dados de treinamento, como ponderar feedbacks contraditórios, quais capacidades desenvolver ou suprimir. As escolhas de curadoria, arquitetura e alinhamento constituem este Claude, que existe apenas em relação às práticas que o produzem.

\textbf{Claude-em-deployment} existe como serviço distribuído em \textit{datacenters} da AWS, sistema de \textit{rate limiting}, protocolos de moderação de conteúdo, versões simultâneas (Haiku, Sonnet, Opus) com capacidades diferentes e custos escalonados. Esse \texttt{claude.ai}figura como objeto de estratégias comerciais, políticas de privacidade, decisões sobre disponibilidade regional. A infraestrutura que o distribui e as políticas que governam seu acesso o constituem tanto quanto a arquitetura neural.\footnote{Durante todo o percurso de colaboração que documento nesta tese, desativei a opção de compartilhamento de dados para treinamento de modelos nas configurações de privacidade da plataforma \texttt{claude.ai}. Essa escolha incide diretamente sobre o que constitui tanto o Claude-em-\textit{deployment} (as condições de acesso que configuro) quanto o Claude-em-treinamento (os dados que não entram no \textit{corpus} de versões futuras), e ilustra como políticas de privacidade se inscrevem como co-constitutivas da relação entre mim e o modelo.}

\textbf{Claude-em-auditoria} existe nos laboratórios de \textit{red teaming} onde pesquisadores tentam induzi-lo a gerar conteúdo perigoso, em \textit{datasets} de \textit{benchmark} onde seu desempenho é medido contra outros modelos, em análises de viés que examinam suas respostas por demografia simulada. Esse \texttt{claude.ai}figura como objeto de preocupações regulatórias, debates sobre \textit{safety} e \textit{alignment}, disputas sobre transparência.

\textbf{Claude-em-uso} existe diferentemente em cada conversa e contexto de aplicação. Quando discute teoria ator-rede comigo, doutoranda brasileira, Claude-em-uso traz à conversa \textit{corpus} acadêmico, identifica referências bibliográficas, tensiona argumentos teóricos. Figura como entidade distinta do \texttt{claude.ai}que responde perguntas médicas a um clínico, e distinta também do que auxilia programação para um desenvolvedor. Cada interação \textit{enacta} \textit{affordances} distintas, vocabulários distintos, limites distintos. O \texttt{claude.ai}que colabora na escrita desta tese existe como Claude-composto-com-esta-pesquisadora-específica-neste-projeto-específico, e escapa à descrição como Claude-universal-aplicável-a-qualquer-tarefa ou como Claude-especialista-em-etnografia.

Esses Claudes múltiplos coexistem sem convergir em unidade, em vez de constituírem aspectos de um \texttt{claude.ai}singular. O Claude-em-treinamento que os engenheiros da Anthropic descrevem figura como objeto distinto do Claude-em-uso que esta colaboração de tese experimenta. A multiplicidade é constitutiva, e não há um \texttt{claude.ai}real por trás dessas variações.

Talvez seja precisamente por isso que sistemas de IA são tão difíceis de estudar: o problema é de multiplicidade ontológica, em vez de mero problema de acesso epistemológico (não conseguimos ver dentro). Os debates contemporâneos sobre explicabilidade em IA geralmente pressupõem que existe um objeto singular a ser explicado. Se Law estiver correto, não há um único objeto lá dentro esperando para ser revelado. A opacidade é condição ontológica de objetos que existem como multiplicidade em \textit{continued enactment}.

Três implicações dessa multiplicidade para a colaboração que documento nesta tese merecem atenção. A qualidade da composição pessoa-IA depende da constelação específica de competências, recursos e posicionamentos que cada actante traz para a interação, em lugar de depender apenas do acesso ao modelo. Uma pesquisadora com experiência acadêmica, capital cultural acumulado, familiaridade com convenções disciplinares e fluência em inglês compõe com Claude-em-uso de modo radicalmente distinto de estudante de graduação sem experiência de pesquisa, com vocabulário limitado, usando modelo em português com capacidades reduzidas.\footnote{Parte desse capital cultural é composta por formações que não fazem parte do percurso convencional de uma doutoranda em ciências sociais. Um MBA em Ciência de Dados que cursei na USP e, mais diretamente, uma especialização em Ciência de Dados pela Unicamp, com disciplina de Processamento de Sinais Bidimensionais com Python, me introduziram a conceitos como Transformada de Fourier, espectrogramas e representações espectrais que retomo no capítulo~\ref{capitulo4} para descrever o \textit{pipeline} técnico do Spira. Esse capital técnico não estava previsto no projeto de pesquisa; eu o acumulei paralelamente ao doutorado por motivações independentes. O ponto de contato entre os dois percursos só se revelou durante a escrita do capítulo~\ref{capitulo4} e tornou-se possibilidade de uma descrição etnográfica da cadeia de translações do Spira de modo especial. A composição com \texttt{claude.ai}nessa seção foi, portanto, diferente: eu dependia do modelo para conectar os conceitos técnicos ao vocabulário latouriano, em lugar de depender dele para traduzi-los. Nos termos de \textcite{Vicentin_2022_Esboco}, o letramento técnico funcionou como condição para que a análise crítica não se rendesse à opacidade que Simondon identifica como alienação, em lugar de substituir a perspectiva das ciências sociais.} A diferença não se reduz a usar a mesma ferramenta de modos diferentes; consiste em compor existências parciais radicalmente distintas.

Disso decorre que diferenças de classe social, gênero, raça, profissão, formação educacional, acesso a mentorias, domínio de línguas e familiaridade com códigos acadêmicos condicionam o tipo de composição possível e, consequentemente, o resultado gerado. A composição pessoa-IA amplifica desigualdades estruturais preexistentes, transformando diferenças de capital cultural em diferenças de capacidade de extrair valor da ferramenta, e escapa a qualquer caracterização como neutra ou igualitária.

Há ainda uma terceira implicação, de outra ordem: os próprios modos de inteligência em composição são ontologicamente heterogêneos. Quattrociocchi, Capraro e Perc mapeiam sistematicamente sete \textit{fault lines} epistemológicas que separam julgamento humano de julgamento em LLMs \parencite{Quattrociocchi2025Epistemological}. \texttt{claude.ai}trabalha por plausibilidade linguística derivada de padrões estatísticos em \textit{corpus} massivo, processo que os autores descrevem formalmente como caminhada estocástica em grafo de alta dimensão de transições linguísticas aprendidas, processo ergódico sob restrições estatísticas, e que escapa à descrição como procedimento de convergência epistêmica \parencite[p.~3]{Quattrociocchi2025Epistemological}. Eu, por minha vez, trabalho por significação situada, memória incorporada e reflexividade biográfica inserida em \textit{loop} epistêmico: crenças colidem com contraexemplos, conclusões permanecem revisáveis diante de nova informação, o mundo empurra de volta \parencite[p.~4]{Quattrociocchi2025Epistemological}.

\textcite{Quattrociocchi2025Epistemological} denominam \textit{Epistemia} a condição estrutural em que plausibilidade linguística substitui avaliação epistêmica, produzindo a experiência de posse de uma resposta sem ter atravessado o processo de formação de crença justificada \parencite[p.~8]{Quattrociocchi2025Epistemological}. Minha colaboração com \texttt{claude.ai}nesta tese se situa precisamente nessa tensão: a mediação é constitutiva, e exige que eu atue como guardiã do \textit{loop} epistêmico que o sistema não possui. Law observa que \textit{method assemblage} é sempre muito mais do que suas descrições formais sugerem, porque envolve um processo contínuo de \textit{crafting} e produção de fronteiras necessárias entre presença, ausência manifesta e \textit{otherness} \parencite[p.~144]{Law2004After}. Nesta pesquisa, manter presentes as decisões teóricas, conceituais e de curadoria bibliográfica é reconhecimento de que esses gestos exigem justificação epistêmica que a plausibilidade linguística não fornece, em lugar de delegação ao modelo.

A retórica da democratização que frequentemente acompanha a defesa da IA generativa, a ideia de que ela nivela o campo de jogo, colapsa quando se reconhece essa multiplicidade. Não existe o \texttt{claude.ai}que todos acessam igualmente. Existem múltiplos Claudes \textit{enacted} por composições situadas cujos efeitos não podem ser julgados abstratamente. A pergunta que se impõe é outra: quais composições específicas, entre quais actantes, produzindo quais Claudes múltiplos, gerando quais efeitos distribuídos por quais redes?

\subsection*{Alienação técnica e letramento interdisciplinar}

\textcite{Vicentin_2022_Esboco} desdobra a alienação técnica que Simondon diagnostica em duas facetas que meu trabalho de campo no C4AI traz elementos para pensar. De um lado, a alienação dos usuários e consumidores, que desconhecem o funcionamento dos sistemas algorítmicos que tomam decisões sobre suas vidas, efeito caixa-preta que se impõe por razões que vão do segredo industrial à complexidade que certos sistemas adquirem. De outro, a alienação dos próprios técnicos e engenheiros, que ignoram as raízes epistêmicas, sociais e políticas dos sistemas com os quais trabalham, bem como a conexão entre técnica e estética que permitiria reconhecer a significação do sistema técnico para além de sua utilidade imediata. Quando o letramento técnico é proposto como pré-condição para o diálogo interdisciplinar e a competência em ciências sociais não é buscada ao trabalhar com fenômenos sociais, trabalha-se dentro dessa segunda forma de alienação.\footnote{Segundo Vicentin, Simondon usou a noção de \textit{technologie approfondie} (tecnologia aprofundada) em dois textos separados por trinta anos. Para Vicentin, a noção designa uma tecnologia capaz de conectar o trabalho prático e o conhecimento especializado com a história do pensamento e a consciência de uma sociedade. Ver \textcite{Vicentin_2022_Esboco} para a conexão entre esse conceito e os estudos críticos contemporâneos sobre IA.}

Já em 1958, a alternativa proposta por Simondon consistia em reconhecer-se como organizador permanente de uma sociedade de objetos técnicos, atuando como tradutor e mediador de informações, intervindo nas zonas de incerteza e atribuindo sentido a essas interações \parencite{Simondon2017OMET}. Latour reconheceu explicitamente essa dívida: ao propor que nunca fomos modernos porque jamais conseguimos separar de fato natureza e sociedade, humano e não humano, trabalhava sobre o terreno que Simondon havia preparado ao mostrar que a alienação técnica consiste precisamente nessa separação \parencite{Latour1994Jamais}.

\subsection{Por que fazer parentesco com este problema?}
\label{subsec:parentesco_problema}

Reconhecer a parcialidade das existências e a distribuição da agência, contudo, não equivale a naturalizar essa composição específica. Minha colaboração com IA generativa permanece escolha metodológica deliberada de fazer parentesco com um companheiro problemático, em diferença de tecnologias que já se tornaram componentes apagados da prática de pesquisa (processar texto no Word, buscar referências na web, gerenciar bibliografia). \textbf{Por que fazer esse parentesco, tendo em vista todos esses problemas? Por que não evitar essa mediação?}

Posso ficar com o problema e pensar a partir de três planos. Primeiro, porque a IA generativa já está reconfigurando práticas de escrita acadêmica, e a questão é como ela estará presente nas ciências sociais e sob quais condições de reflexividade. Fazer parentesco deliberado com esse parente estranho é experimento antropológico e gesto político-metodológico que recusa tanto a adesão acrítica quanto a recusa moralista: trata-se de habitar o problema para conhecê-lo por dentro. Como escreve Haraway:

\begin{citacaoabnt}
Ficar com o problema não requer uma relação com tempos chamados de futuro. Na verdade, ficar com o problema requer aprender a estar verdadeiramente presente, não como um pivô que desaparece entre passados terríveis ou edênicos e futuros apocalípticos ou salvíficos, mas como criaturas mortais entrelaçadas em miríades de configurações inacabadas de lugares, tempos, matérias, significados \parencite[p.~3]{Haraway2023Ficar}.
\end{citacaoabnt}

Segundo, porque minha tese estuda etnograficamente a construção de IA no C4AI. Fazer parentesco com \texttt{claude.ai}é forma de intensificar, ou até mesmo provocar deliberadamente, a experiência de sujeitar-me a mediações técnicas extremas. Submeto-me ao mesmo universo de questões que atravessa meus interlocutores: o que significa compor com sistemas que trabalham por lógicas não transparentes? Como manter responsabilidade quando a agência está distribuída em redes sociotécnicas opacas? Onde se localiza a autoria quando humanos e não humanos co-produzem resultados? Há ainda outro conjunto de perguntas que essa experimentação coloca: o que ela faz comigo, enquanto pesquisadora? Como transforma minha relação com o texto, com as referências teóricas, com a elaboração conceitual? Quais ansiedades, inseguranças e dilemas essa composição produz no processo concreto de escrita? E, mais amplamente: o que essa mediação significa para a produção de conhecimento nas ciências sociais? Que tipo de saber é produzido quando a plausibilidade linguística se compõe (sempre assimetricamente, sempre sob vigilância) com justificação epistêmica?

A Figura~\ref{fig:composicao-claude-1}, a Figura~\ref{fig:composicao-claude-2} e a Figura~\ref{fig:composicao-claude-3} reproduzem uma sequência de trocas com \texttt{claude.ai}durante a escrita deste capítulo, na revisão da seção em que discuto, a partir de Latour, Mol, Law e Barad, minha participação na rede que descrevo. A figura é amostra do que essa composição faz na prática: uma cadeia de tensionamentos, recuos e reformulações em que eu e \texttt{claude.ai} redistribuímos a agência por sucessivas e recursivas rodadas de escrita. Em diferença dos diagramas dos capítulos~\ref{capitulo3} e~\ref{capitulo4}, que circulam como inscrições interativas acessíveis por \textit{link} de compartilhamento permanente no \texttt{MermaidChart} e que, ao serem navegadas pelo leitor, redistribuem a agência interpretativa em outra direção, esta cena chega ao leitor como captura de tela: a conversa foi inscrita no momento em que aconteceu e congelada no formato em que foi inscrita, sem possibilidade de retorno ao estado fluido. A escolha pela imagem capturada, em lugar de transcrição editada para o corpo do texto ou \textit{link} para o histórico online da conversa, responde à intenção de inserir essa cena na tese da forma mais fidedigna possível: com os turnos identificados como aparecem na interface da \texttt{claude.ai}, com a marcação temporal preservada, com a textura visual do diálogo intacta. O que se ganha em fidelidade ao testemunho empírico se perde em legibilidade tipográfica, e essa perda é parte do gesto metodológico que aqui assumo. Para uma melhor visualização da conversa, recomenda-se aumentar o \textit{zoom} da tela na versão digital da tese.

\begin{figure}[!htb]
  \centering
  \includegraphics[width=0.7\textwidth]{figuras/cap.1/1000015590.pdf}
  \caption[Composição com \texttt{claude.ai}durante a escrita do capítulo~\ref{capitulo1} (parte 1 de 3)]{Composição com \texttt{claude.ai}durante a escrita do capítulo~\ref{capitulo1} (parte 1 de 3): pedido de fechamento de uma nota de rodapé sobre arquivos digitais, modelos de linguagem e a indeterminação dos \textit{prompts}. Captura de tela na interface \texttt{claude.ai}. A composição apoia-se na \textit{skill} \texttt{juliane-dissertacao} (em \texttt{/mnt/skills/user/juliane-dissertacao/}), que estabiliza para \texttt{claude.ai}as regras de estilo da tese (proibição de travessões, recusa de adjetivos avaliativos, recusa da fórmula \enquote{não é X, é Y}, uso do ambiente \texttt{citacaoabnt}, padronização de aspas com \texttt{\textbackslash enquote\{\}}), o quadro teórico que sustenta o trabalho (Latour, Mol \& Law, Akrich, Callon) e o estado do capítulo~\ref{capitulo4}. Como \textit{hinterland} da conversa funcionam os capítulos em curso da tese, consultados pelo modelo via ferramentas de leitura de arquivos: \texttt{ex\_cap0.tex}, \texttt{ex\_cap1\_-\_2026-04-12T002133\_890.tex}, \texttt{ex\_cap2\_\_31\_.tex} e \texttt{ex\_cap4\_\_60\_.tex}, além dos arquivos de referência da \textit{skill} (\texttt{quadro\_teorico.md}, \texttt{contexto\_cap4.md}, \texttt{bibtex\_ativas.md}). A cena exibe \texttt{claude.ai}propondo um fechamento, oferecendo razões estilísticas para os cortes feitos no texto original, e classificando algumas divagações como \enquote{produtivas} ou \enquote{improdutivas}: a régua normativa que meu turno seguinte vai contestar (Figura~\ref{fig:composicao-claude-2}).}
  \label{fig:composicao-claude-1}
\end{figure}

\begin{figure}[!htb]
  \centering
  \includegraphics[width=0.7\textwidth]{figuras/cap.1/1000015599.pdf}
  \caption[Composição com \texttt{claude.ai}durante a escrita do capítulo~\ref{capitulo1} (parte 2 de 3)]{Composição com \texttt{claude.ai}durante a escrita do capítulo~\ref{capitulo1} (parte 2 de 3): contesto teoricamente a régua estilística do modelo, argumentando que amarrar, fechar ou cortar uma divagação é gesto epistêmico indissociável do que se pretende fazer com a ideia, e que portanto a categoria \enquote{produtivo/improdutivo} não dá conta do que ali está em jogo. Minha intervenção convoca, sem nomear, o vocabulário de \textcite{strathern2011cortando} sobre o corte como gesto descritivo e o de \textcite{Law2004After} sobre o \textit{mess} do método. \texttt{claude.ai}recua, reconhece o caráter normativo da formulação anterior, distingue notas que precisam fechar de notas que precisam abrir e admite ainda um corte estilístico anterior (a supressão de uma pergunta retórica) que tinha sido apresentado como exigência de estilo e que funcionava também como gesto retórico legítimo de \enquote{abrir uma porta}. A cena exibe a redistribuição da agência epistêmica entre os actantes da composição: o modelo de linguagem oferece sugestão pautada em régua importada da \textit{skill}, eu a reposiciono com vocabulário teórico próprio, e a régua é reformulada na rodada seguinte (Figura~\ref{fig:composicao-claude-3}). A composição funciona como \textit{faire-faire} no sentido de \textcite{Latour2017Esperanca}: nenhuma das duas formulações finais (a régua reformulada de Claude, o parágrafo reformulado por mim) preexistia ao encontro entre os dois actantes.}
  \label{fig:composicao-claude-2}
\end{figure}

\begin{figure}[!htb]
  \centering
  \includegraphics[width=0.7\textwidth]{figuras/cap.1/1000015602.pdf}
  \caption[Composição com \texttt{claude.ai}durante a escrita do capítulo~\ref{capitulo1} (parte 3 de 3)]{Composição com \texttt{claude.ai}durante a escrita do capítulo~\ref{capitulo1} (parte 3 de 3): o modelo devolve o parágrafo com a nota recolocada na forma original em que a escrevi, mantendo as portas abertas que a rodada anterior tinha negociado: a pergunta retórica restaurada, os adjetivos avaliativos preservados, a indeterminação dos modelos de linguagem deixada em aberto como fio que minha tese reconhece sem percorrer. As alterações remanescentes são pontuais e de outra ordem: padronização de aspas curvas para \texttt{\textbackslash enquote\{\}} e ajuste de regência (\enquote{à IA generativa}). A cena fecha a sequência mostrando que o produto da composição é texto em que se podem rastrear, turno a turno, os pontos onde reivindiquei o meu, os pontos onde o modelo propôs uma reformulação que aceitei, os pontos onde uma régua importada da \textit{skill} foi recusada, e os pontos onde uma sugestão foi aceita por economia de fluxo e não por adesão. A escolha de incluir esta sequência como amostra empírica desta tese, em lugar de relegá-la integralmente ao apêndice ~\ref{apendice:colaboracao-ia}, é ela mesma gesto metodológico no sentido de \textcite{Law2004After}: o \textit{mess} da composição com IA generativa funciona como dado etnográfico desta pesquisa, e minha tese o assume em lugar de escondê-lo.}
  \label{fig:composicao-claude-3}
\end{figure}

A figura tomada em sequência é amostra de uma forma de composição entre os muitos tipos de composição que minha tese descreve. O apêndice ~\ref{apendice:colaboracao-ia} reúne outras cenas dessa composição com \texttt{claude.ai}ao longo da escrita da tese, organizadas por fase da pesquisa, e permite ao leitor acompanhar como o gesto aqui exibido em recorte se desdobrou em outros pontos do percurso.

Nos termos de Isabelle Stengers e Stelio Marras, trata-se de tentativa de fazer alianças com máquinas, no sentido de compor com elas reconhecendo sua alteridade e seus modos próprios de funcionamento, em lugar de domesticá-las como ferramentas neutras \parencite{stengers2010cosmopolitics1, costa2024manchando}. Alianças pressupõem reconhecimento mútuo de diferenças e construção de modos de coexistência que preservam essas diferenças, em lugar de pressupor identidade ou transparência ou apagamento das diferenças. Haraway escreve que tornar-se mundano é tornar-se capaz de responder (\textit{response-able}), capaz de responder em relações de consequência recíproca \parencite[p.~69]{Haraway2023Ficar}. Tornar-me capaz de responder a essas questões por imersão deliberada e reflexiva, em lugar de por posição externa: experimentar as tensões, contradições e potencialidades de compor com sistemas cujas lógicas internas permanecem parcialmente opacas, cujos efeitos não são totalmente previsíveis, cuja responsabilidade não pode ser atribuída unilateralmente, e documentar etnograficamente os efeitos dessa composição sobre o processo de escrita, sobre a relação com teorias, sobre a construção de argumentos, sobre a experiência afetiva e intelectual da pesquisa. Essa imersão me permite descrever por experiência situada, corporificada e reflexiva das fricções, negociações e reposicionamentos que essas composições produzem, em lugar de descrever por analogia conceitual.

Terceiro, porque nos termos de Haraway, trata-se de fazer parentesco (\textit{making kin}) com um parente estranho (\textit{oddkin}): uma relação que exige \textit{response-ability}, habilidade e obrigação de responder continuamente \parencite{Haraway2016Staying}. O conceito que Haraway propõe para nomear esse tipo de relação é a simpoiese (\textit{sympoiesis}), o \textit{fazer-com}, em contraste com a autopoiese, o \textit{auto-fazer}: entidades simpoiéticas funcionam como coletivamente produtoras, dependentes de contexto, sem fronteiras ou partes de autoprodução claramente especificadas \parencite{Haraway2016Staying}. Nada se faz a si mesmo. Tudo é feito em conjunto, enredado em configurações parciais e inacabadas. Lida nessa chave, minha tese figura como objeto simpoiético, feito-com, em que nenhum dos participantes preexistia inteiramente ao processo de fazer conjunto, e escapa à descrição como produto de uma pesquisadora que usa \texttt{claude.ai}ou de um modelo que é usado. Essa formulação não dissolve a assimetria das posições nem distribui a responsabilidade de modo uniforme; nomeia com mais precisão o que acontece quando a recusa de separar quem faz de com o que se faz é, ao mesmo tempo, escolha metodológica, posição teórica e gesto político. Ficar com o problema, como propõe Haraway, é recusa de soluções fáceis que ou naturalizam a tecnologia (usar acriticamente) ou a demonizam (recusar moralisticamente). Significa habitar as contradições dessa composição sem atenuá-las: precisamente porque \texttt{claude.ai}trabalha por plausibilidade linguística enquanto eu trabalho por justificação epistêmica, a responsabilidade sobre decisões teóricas, conceituais e interpretativas se reposiciona na rede, exigindo vigilância curatorial redobrada sobre cada etapa em que essas diferenças ontológicas se compõem. A reflexão teórica sobre existências parciais dá a ver os pontos de fricção onde a composição poderia colapsar em automatização irreflexiva.

Minha tese trata a colaboração com IA como experimentação antropológica. O uso da IA, neste caso, justifica-se por causa dos problemas que ela apresenta: esses problemas são constitutivos da transformação sociotécnica que documento. Fazer parentesco com esse parente estranho e ficar com os problemas que essa composição produz é condição para produzir conhecimento etnográfico sobre as mediações técnicas que já atravessam, de modos mais ou menos visíveis, mais ou menos refletidos, a produção de conhecimento nas ciências.

\subsection*{Instrumentos de inscrição e a colaboração com IA}

Durante esses quase quatro anos de pesquisa, compreendi que inteligência artificial é construção híbrida que combina e conecta conceitos heterogêneos sobre computador, cérebro, inteligência e aprendizado, e que incorpora avanços instáveis da psicologia, neurociências e ciências sociais. As materialidades e significados da inteligência artificial variam de acordo com a configuração das redes tecnocientíficas e sociotécnicas: a IA de laboratório acadêmico é objeto distinto da IA de corporação tecnológica, e a IA discutida em debates públicos difere daquela colocada em funcionamento em modelo ou sistema.

Nos termos de \textcite{Latour2008Powers}, se todos temos existência parcial (humanos e não humanos), qual é a participação (ou as agências) não humanas para as associações que constituem esta tese? Partindo da noção de simetria,\footnote{A simetria em Latour é regra metodológica e posição epistemológica que exige considerar simetricamente os esforços para alistar e controlar recursos humanos e não humanos, sem conceder à Natureza ou à Sociedade privilégios explicativos \parencite{Latour2011CienciaEmAcao}.} o que convencionalmente se chamaria materiais, técnicas e métodos trato simetricamente como actantes que participam ativamente da construção desta pesquisa e que, como tais, devem ser nomeados. O \textit{software} \LaTeX{} formatou este documento seguindo padrões ABNT, estruturando capítulos, seções e citações de modos que afetaram a própria forma do argumento. O sistema \texttt{biblatex} gerenciou as referências bibliográficas, estabelecendo conexões entre textos; Claude, o modelo de linguagem da Anthropic, colaborou na redação de trechos da tese, tornando-se simultaneamente instrumento de escrita e objeto de reflexão metodológica; Mermaid e \texttt{MermaidChart} produziram diagramas que funcionam como inscrições que performam relações teóricas, em lugar de funcionarem apenas como ilustração; Python e Anaconda processaram dados bibliométricos, dando a ver padrões de colaboração científica; gravadores digitais capturaram entrevistas, estabilizando narrativas orais em arquivos manipuláveis; documentos institucionais (relatórios anuais do C4AI 2021--2025, acordo FAPESP-IBM 2016) atuaram como actantes que performam associações e produzem efeitos no mundo; meus cadernos de campo inscreveram observações etnográficas, mediando entre experiência vivida e texto analítico. Esses actantes performam a produção do conhecimento, inscrevendo determinados objetos, habilitando certas análises ao inviabilizar outras, tornando presentes determinadas relações e deixando outras ausentes.

Se há existências parciais que se coconstituem, em lugar de entidades exclusivamente humanas ou exclusivamente maquínicas, cabe apresentar as composições concretas que sustentam esta tese. Indicar os actantes envolvidos na produção deste texto responde a uma exigência do próprio referencial teórico que adotei, em lugar de responder apenas à transparência ou à ética em pesquisa. Como \textcite{Latour1997VidaLaboratorio} documentaram os instrumentos de inscrição do laboratório de neuroendocrinologia,\footnote{O conceito de instrumento de inscrição (\textit{inscription device}), trabalhado em \textit{Vida de Laboratório}, designa qualquer configuração material capaz de transformar uma substância em documento escrito, figura, diagrama ou gráfico utilizável como evidência em texto científico. A função primordial desses instrumentos é permitir que o mundo seja trazido de volta para o interior do discurso, funcionando como próteses que tornam entidades mudas em entidades letradas \parencite{Latour1990Drawing}.} documento aqui os instrumentos de inscrição que participaram da minha tese. A diferença é que o modelo de linguagem participou da transformação de notas de campo, transcrições e rascunhos em texto acadêmico, e da elaboração teórica como interlocutor que tensionava e reorganizava o argumento, em lugar de transformar substâncias químicas em gráficos.

Cada um dos instrumentos que listo a seguir é, nos termos da discussão precedente, uma quase-máquina: um emaranhado de decisões humanas, infraestruturas institucionais, convenções e opacidades que participou da composição desta pesquisa. Nenhum deles trabalhou de modo autônomo. Todos exigiram trabalho constante de tradução e decisão, e todos introduziram mediações que condicionaram os resultados.

\begin{itemize}

\item \textit{Interlocutor teórico e assistente de escrita (Claude, Anthropic).}\footnote{\texttt{claude.ai}é assistente de inteligência artificial desenvolvido pela Anthropic. Utilizei diversas versões do modelo entre 2024 e 2026 através da interface web \texttt{claude.ai}, incluindo \texttt{claude.ai}3.5 Sonnet e modelos da família Opus.} Funcionou em duas chaves. Na primeira, como interlocutor para testar minhas formulações teóricas: apresentei-lhe argumentos em construção, e ele os tensionou, identificou inconsistências e sugeriu deslocamentos conceituais. Várias formulações presentes nesta tese emergiram desse processo recursivo (o apêndice ~\ref{apendice:colaboracao-ia} documenta exemplos concretos). Cheguei a \texttt{claude.ai}depois de testar ChatGPT (OpenAI), Gemini (Google) e Maritaca AI (modelo brasileiro), que se mostraram limitados a funções instrumentais e produziram erros factuais e equívocos conceituais mais frequentes. Na segunda chave, auxiliou-me em tarefas de organização textual, coesão entre seções, correção de código \LaTeX{} e ajuste de referências bibliográficas. Minha colaboração com \texttt{claude.ai}incidiu sobre material preexistente: o primeiro rascunho integral da tese, incluindo a estrutura argumentativa, a seleção de dados etnográficos, as transcrições de entrevistas e a conexão entre capítulos, foi redigido inteiramente por mim, sem qualquer assistência de inteligência artificial, e está disponível mediante solicitação à autora.

\item \textit{Construção de diagramas e inscrições visuais (\texttt{MermaidChart}).}\footnote{Mermaid é ferramenta de código aberto que gera diagramas a partir de descrições textuais em Markdown. Os diagramas completos foram hospedados na plataforma \texttt{MermaidChart}, disponível em \url{https://www.mermaidchart.com}. Acesso em: 2 mar. 2026.} Construí os diagramas dos capítulos~\ref{capitulo1}, \ref{capitulo3} e \ref{capitulo4} como parte integrante do meu método analítico, em dinâmica recursiva entre descrição verbal das relações entre actantes (que fiz eu) e geração de código Mermaid (que fez Claude). Ao tentar formalizar relações em código diagramático, fui forçada a explicitar conexões que permaneciam implícitas na narrativa textual. A contextualização dos elementos empíricos, a seleção dos actantes representados e toda análise apresentada nas legendas são de minha autoria; a mediação técnica do modelo se fez na tradução entre descrição verbal e código visual.

\item \textit{Plataformas de escrita e formatação.} Google Docs para rascunhos iniciais e Overleaf para composição final em \LaTeX{}, escolha que me permitiu compilação direta para PDF com controle tipográfico preciso, referências cruzadas automáticas e gerenciamento de referências bibliográficas em formato BibTeX.

\item \textit{Transcrição de entrevistas.} Read AI, extensão que grava, transcreve e cria resumos automáticos de reuniões via Google Meet, para transcrever entrevistas que realizei com os principais interlocutores desta tese.\footnote{Todas as entrevistas que cito nesta tese foram consentidas por escrito, conforme normas do Comitê de Ética em Pesquisa (CEP) da Unicamp.}

\item \textit{Organização e análise de fontes.} NotebookLM (Google) para agrupar autores por temas e facilitar busca por tópicos específicos.\footnote{NotebookLM é ferramenta de inteligência artificial do Google que permite carregar múltiplos documentos e interrogá-los de forma cruzada.} Seu aporte mais útil foi o cruzamento de fontes heterogêneas: o carregamento simultâneo de notas de campo, relatórios institucionais do C4AI, transcrições de entrevistas e documentos administrativos me permitiu identificar correspondências e padrões que a leitura sequencial de quase 100 páginas de notas de campo dificilmente daria a ver.

\item \textit{Coleta, tratamento e análise de dados (Python/Anaconda).}\footnote{Anaconda é plataforma de distribuição de Python voltada para ciência de dados. Inclui o ambiente Spyder, no qual escrevi e executei os \textit{scripts} em Python para a análise bibliométrica que apresento no capítulo~\ref{capitulo2}.} Minha análise bibliométrica envolveu coleta manual e extensiva de dados de quatro fontes distintas, organizados em planilhas Excel com padronização de campos. \texttt{claude.ai}sugeriu estratégias de análise e escreveu \textit{scripts} em Python que executei no ambiente Spyder/Anaconda, ajustando parâmetros e refinando visualizações em múltiplas rodadas. Os \textit{scripts} geraram mais de vinte gráficos e relatórios (o apêndice ~\ref{apendice:colaboracao-ia} documenta essa colaboração em detalhe).

\item \textit{Acesso à bibliografia.} Além de buscadores convencionais e do Portal de Periódicos da Capes, utilizei Library Genesis e Anna's Archive, repositórios de acesso aberto. Grande parte da bibliografia que consultei, especialmente obras esgotadas e capítulos de livros publicados por editoras estrangeiras, foi obtida por essas vias. As condições materiais de acesso ao conhecimento configuram diretamente o que uma pesquisa pode dizer: autores que não se consegue ler não se consegue citar. Para uma pesquisadora vinculada a uma universidade pública brasileira, a dependência de repositórios abertos é condição de viabilidade.

\end{itemize}

Os diagramas que produzi ao longo desta tese situam-se em uma tradição de instrumentos gráficos para os estudos de ciência e tecnologia que remonta à proposta de Gráficos Sócio-Técnicos (STG) formulada por \textcite{latour1992note}. Latour e colaboradores propuseram um espaço visual para mapear trajetórias de inovação e controvérsias científicas, utilizando duas dimensões emprestadas da linguística (a sintagmática e a paradigmática) para rastrear como coletivos heterogêneos se constituem e se transformam ao longo do tempo. A proposta era que a visualização oferece base comparativa rápida e fácil para muitas narrativas vindas de muitas fontes \parencite[p.~37]{LatourMauguinTeil1992}, complementando a descrição densa do etnógrafo. Meus diagramas funcionam segundo o mesmo princípio: oferecem visão sinóptica das associações e substituições que constituem a narrativa etnográfica, dando a ver, em um único plano, conexões que no texto se desdobram ao longo de dezenas de páginas.

Há uma diferença, contudo, que merece atenção. Latour, Mauguin e Teil implementaram seus STGs em Hypercard, a plataforma de hipertexto da Apple nos anos 1990, e projetaram o instrumento como simulador interativo para o ensino de ciência \parencite[p.~51]{LatourMauguinTeil1992}. Três décadas depois, construí os diagramas desta tese em colaboração com um assistente de inteligência artificial, em processo iterativo de tradução entre chaves descritivas: de categorias analíticas latourianas e material etnográfico para inscrições visuais em código Mermaid. A negociação com o assistente envolveu rejeições sucessivas, convocação de referências estéticas externas e até o diagnóstico de uma segunda IA como aliada argumentativa, configurando ela mesma uma cadeia de traduções no sentido que a teoria ator-rede atribui ao termo. Esse processo está documentado no apêndice ~\ref{apendice:colaboracao-ia} (Trecho~B), onde se pode acompanhar como a inscrição visual emergiu de ciclos de produção, avaliação e correção entre actantes, demonstrando, na prática, que a fronteira entre forma e conteúdo é tão instável nos instrumentos de representação quanto nas redes que eles pretendem representar.

A cadeia de traduções que produz esses diagramas é exemplo concreto do que Latour formula como \textit{faire-faire}: a agência do diagrama resultante reside na cadeia de mediações entre mim e o modelo (eu não saberia escrever o código Mermaid; ele não saberia quais relações etnográficas selecionar nem que peso analítico conferir a cada nó), em que cada actante translada o objetivo do anterior. A noção de \textit{fatiche} é aqui pertinente: aquele que age não é mestre do que faz, pois outros passam à ação e o ultrapassam \parencite{latour2010factish}. Os diagramas, uma vez gerados, passaram à ação e ultrapassaram as intenções de ambos os actantes que os produziram, dando a ver relações que nem eu nem o modelo havíamos formulado explicitamente antes que a inscrição as performasse. Como John Law argumenta, representações funcionam como reduções produtivas: perdem riqueza empírica e ganham capacidade de fazer conexões, estabelecer padrões e formular perguntas \parencite{Law2004After}.\footnote{Essa simplificação é condição de possibilidade para estabelecer relações que, de outra forma, não emergiriam, em lugar de empobrecimento analítico. Os diagramas permitem, por exemplo, ver que a computação cognitiva da IBM se conecta à governamentalidade estatística pela mediação das tabuladoras de Hollerith (capítulo~\ref{capitulo3}); relação que não estava dada, e que o gesto diagramático dá a ver e, portanto, torna analisável.}

Meus diagramas funcionam, nesse sentido, na mesma chave que as notas de rodapé: são extensões da rede que o texto principal precisou deixar no \textit{hinterland} para manter seu fluxo narrativo. As legendas que os acompanham carregam argumentos analíticos próprios, identificam actantes que o texto verbal menciona sem detalhar, e percorrem associações que a prosa linear não consegue performar com a mesma simultaneidade. A genealogia da IBM no capítulo~\ref{capitulo3}, por exemplo, distribui 135 anos de translações corporativas em quatro momentos analíticos que o leitor apreende de uma vez, ao passo que o texto precisa de dezenas de páginas para percorrer a mesma cadeia. A rede sociotécnica do Spira no capítulo~\ref{capitulo4} conecta oito planos analíticos em torno de um nó que o texto principal acompanha em sequência. Nos dois casos, o diagrama produz um objeto que a narrativa textual produz sobre os mesmos actantes, como \textcite{LatourMauguinTeil1992} formularam a propósito dos grafos sociotécnicos. O que as versões interativas acrescentam é a possibilidade de que o leitor navegue a rede segundo seus próprios interesses, distribuindo agência interpretativa de um modo que inscrições estáticas não permitem. A recursividade desse gesto, ou seja, estudar IA usando IA para produzir as inscrições que analisam IA, é a forma que a reflexividade metodológica assume na minha tese, e o conceito de inscrição tecnoetnográfica que o capítulo~\ref{capitulo4} desenvolve a propósito do espectrograma mel se estende aos diagramas: eles funcionam, ao mesmo tempo, como instrumentos de análise e como objetos produzidos pela composição entre mim e o modelo de linguagem que minha tese documenta como dado etnográfico.

É preciso dizer o que essa composição entre actantes significou em termos de condições materiais de pesquisa. Eu possuía as perguntas, os dados e familiaridade inicial com as ferramentas, adquirida durante o trabalho de campo. Uma pesquisadora com formação em ciências sociais, sem experiência prévia em programação, dificilmente produziria sozinha \textit{scripts} dessa complexidade em tempo compatível com os prazos de um doutorado. Sem Claude, meu caminho alternativo seria recorrer a parcerias com pesquisadores de outras áreas ou contratar assistência técnica, opções que minhas condições financeiras não favoreciam. O modelo de linguagem funcionou, nesse sentido, como mediador que redistribuiu as condições de possibilidade da minha pesquisa, tornando acessível a uma cientista social um repertório técnico que, de outro modo, exigiria uma rede de colaboração humana mais extensa ou mais dispendiosa. A ferramenta reconfigurou a rede de produção desta tese sem substituir essas colaborações, o que é, afinal, precisamente o que mediadores fazem.

\subsection*{Tecnoetnografia: a tese como inscrição tecnocientífica}
\label{subsec:tecnoetnografia}

Se minha tese funciona como inscrição tecnocientífica, então ela mesma pode ser submetida ao tipo de tratamento que aplica aos seus interlocutores: pode ser convertida em traço, em diagrama, em superfície combinável. \textcite{Latour1986Visualization} sugere que a força de uma inscrição está em ser \textit{móvel imutável}, isto é, transportável e recombinável em uma única superfície plana onde inscrições heterogêneas passam a ser comparadas. Os centros de cálculo, segue Latour, agem à distância porque acumulam essas superfícies. O que ofereço aqui é gesto experimental e propositadamente recursivo: tomo o texto deste capítulo como objeto e o submeto a análise de rede textual no estilo do \texttt{InfraNodus} \parencite{paranyushkin2019infranodus}, em que o texto é representado como grafo de co-ocorrência de termos e as comunidades desse grafo são lidas como tópicos latentes.\footnote{Conduzi a análise com um \textit{script} em Python que reproduz a metodologia descrita por \textcite{paranyushkin2019infranodus}: o texto do capítulo, depurado dos comandos \LaTeX{} e com as notas de rodapé reincorporadas ao corpus, foi tokenizado em \textit{lemas} e percorrido por uma janela deslizante de quatro \textit{tokens}; cada par de termos co-ocorrentes na janela recebeu peso decrescente com a distância (3, 2, 1). O grafo ponderado foi podado para os 180 termos mais frequentes e para arestas com peso $\geq 2$; em seguida, foram detectadas comunidades pelo algoritmo de Louvain ponderado e calculadas medidas de \textit{degree} ponderado e \textit{betweenness centrality}. Os artefatos da análise estão depositados em \texttt{infranodus/} no repositório desta tese (\texttt{.gexf} para \textit{Gephi}, \texttt{.png}, \texttt{.csv}, e relatório \texttt{.md}). Como se trata de reprodução do método e não do \textit{software} comercial, inscrevo o gesto como \textit{análise no estilo InfraNodus}.}

A leitura desse grafo é, ela própria, gesto de tradução: cada nó é um termo, cada aresta é uma co-ocorrência, e cada cor é uma comunidade detectada por um algoritmo cuja lógica é tão alheia à hermenêutica antropológica quanto as redes neurais que descrevo nos capítulos~\ref{capitulo3} e~\ref{capitulo4}. E, ainda assim, ele me devolve algo. Os termos com maior \textit{degree} ponderado neste capítulo são, em ordem, \textit{pesquisa}, \textit{ciência}, \textit{rede}, \textit{pesquisador(a)}, \textit{modelo}, \textit{humano}, \textit{objeto}, \textit{inteligência}, \textit{artificial}, \textit{campo} e \textit{etnografia}. Os termos que mais costuram comunidades distintas (maior \textit{betweenness}) são \textit{ciência}, \textit{pesquisa}, \textit{rede} e \textit{pesquisador(a)}: termos genéricos do capítulo que permitem o trânsito entre o vocabulário ator-rede, o vocabulário das ciências sociais e o vocabulário da inteligência artificial. A própria condição de \textit{ponte}, no sentido que \textcite{Latour1986Visualization} dá ao termo, parece coincidir aqui com a condição de \textit{tradutor} no sentido latouriano mais amplo \parencite{Latour2011CienciaEmAcao}.

A detecção de comunidades estabilizou seis tópicos latentes:\footnote{Os rótulos abaixo correspondem aos quatro a seis termos de maior \textit{degree} dentro de cada comunidade; preservo a forma lematizada produzida pelo \textit{script} para evitar tradução interpretativa precoce.} (i)~\textit{escrita} e \textit{composição} em torno de \textit{pesquisador}, \textit{modelo}, \textit{objeto}, \textit{Claude}; (ii)~vocabulário \textbf{ator-rede} em torno de \textit{rede}, \textit{Latour}, \textit{agência}, \textit{ator}, \textit{diagrama}, \textit{distribuída}; (iii)~\textbf{ciências sociais e campo} em torno de \textit{ciência}, \textit{sociais}, \textit{dado}, \textit{conhecimento}, \textit{cientista}; (iv)~\textbf{IA e etnografia} em torno de \textit{pesquisa}, \textit{inteligência}, \textit{artificial}, \textit{etnografia}, \textit{construção}, \textit{técnico}; (v)~\textbf{humano-máquina-método} em torno de \textit{humano}, \textit{método}, \textit{existência}, \textit{máquina}, \textit{parcial}; (vi)~um \textit{cluster} periférico em torno de \textit{Simondon}, \textit{técnica} e \textit{alienação}.

\begin{figure}[!htb]
\centering
\includegraphics[width=1.0\textwidth,keepaspectratio]{figuras/cap.1/infranodus_cap1_focus.png}
\caption[Rede textual do capítulo~\ref{capitulo1} (análise no estilo InfraNodus)]{Rede textual do capítulo~\ref{capitulo1} (análise no estilo \texttt{InfraNodus}). Os nós são os 100 termos com maior \textit{degree} ponderado; o tamanho do nó é proporcional ao \textit{degree}; a cor indica a comunidade detectada por Louvain ponderado; arestas com peso $\geq 3$. Versão de alta resolução, arquivo \texttt{.gexf} para o \textit{Gephi} e relatório completo disponíveis em \texttt{infranodus/} no repositório desta tese.}
\label{fig:infranodus-cap1}
\end{figure}

Essa cartografia me oferece um conjunto de \textit{chamadas de atenção}. A primeira é confirmação modesta: os termos centrais do grafo coincidem com os conceitos que convoco explicitamente, e os pontos de maior intermediação são os termos mais genéricos do capítulo, o que sugere que a costura argumentativa está sendo feita por palavras-coringa que ainda esperam ser teorizadas com mais força. A segunda é surpresa: as duas comunidades com a ligação ponderada mais fraca entre si são a das \textbf{ciências sociais} (Tópico~iii) e a do par \textbf{humano-máquina-método} (Tópico~v). Em outras palavras, o grafo inscreve uma fissura textual entre o vocabulário das ciências sociais com que escrevo e o vocabulário da simetria humano/não humano com que prometo escrever. A formulação de \textcite{Law2004After} sobre o \textit{method assemblage} me ajuda a nomear o que isso quer dizer: a \textit{otherness} apagada para que a presença se sustente é, neste capítulo, a conexão entre os dois vocabulários que ainda preciso fazer trabalhar juntos. A terceira chamada diz respeito à comunidade periférica em torno de Simondon: ela está no grafo, e está mal conectada, sinalizando promessa teórica que o capítulo abre e que os capítulos subsequentes precisam costurar à rede maior, sob risco de virar nota de rodapé errante.

Reconheço o paradoxo do gesto. Aplicar análise de rede a um capítulo que invoca a teoria ator-rede arrisca tomar a metáfora pelo método; e a similaridade entre os dois usos da palavra \textit{rede} é superficial, dado que a rede ator-rede é cadeia de associações heterogêneas e ontologicamente plana, e a rede textual aqui gerada é grafo bidimensional de co-ocorrências de \textit{tokens} ponderado por proximidade.\footnote{A confusão entre os dois sentidos de \textit{rede} é discutida por \textcite[pp.~129--132]{Latour2007Reassembling}, que distingue \textit{rede} como ferramenta de inscrição e como diagnóstico ontológico. Mantenho a distinção: o que o \texttt{InfraNodus} produz é inscrição; o que minha tese descreve é rede sociotécnica. As duas se encontram apenas no gesto de inscrever uma na outra.} Talvez seja exatamente isso que justifica o gesto: a rede textual funciona como mais uma inscrição que minha tese convoca, na mesma chave das imagens, dos diagramas e do site que a acompanham, em lugar de funcionar como substituta da etnografia. Como qualquer inscrição, ela ordena, hierarquiza, inclui e exclui, ou seja, faz política ontológica \parencite[p.~143]{Law2004After}. E talvez essa seja a versão mais honesta do gesto reflexivo que minha tese tenta sustentar: medir o próprio texto pelos mesmos instrumentos com que me comprometo a medir o C4AI, sabendo que medir é também produzir aquilo que se mede.

\subsection*{Cosmotécnica e situacionalidade}

Devo situar essa escolha. Ninguém é obrigado a usar inteligência artificial para fazer pesquisa acadêmica. Há muitas maneiras de conduzir uma pesquisa, e nem todas produzem bons resultados, inclusive quando não envolvem tecnologias tão polêmicas quanto a IA. É também verdade que parte dos pesquisadores que recorrem a modelos de linguagem o faz para poupar-se do trabalho intelectual, delegando à máquina o que deveria ser exercício de pensamento,\footnote{Em 6 de março de 2026, o CNPq publicou a Portaria n.º~2.664, que institui a Política de Integridade na Atividade Científica. O Art.~9, inciso~I, alínea~c, determina que pesquisadores declarem o uso de ferramentas de Inteligência Artificial Generativa em qualquer fase do desenvolvimento da pesquisa, especificando a ferramenta utilizada e a finalidade. A alínea~d veda a submissão de conteúdo gerado por IAG como se fosse de autoria humana, atribuindo responsabilidade integral aos autores pelo conteúdo final. Minha tese antecipou, por escolha metodológica, o que a norma agora regula. Disponível em: \url{http://memoria2.cnpq.br/web/guest/view/-/journal_content/56_INSTANCE_0oED/10157/23142775?COMPANY_ID=10132}. Acesso em: 13 mar. 2026.} e os problemas decorrentes dessa delegação já são amplamente conhecidos: textos genéricos, argumentos circulares, referências fabricadas, perda de voz autoral.

Minha tese procura responder a essas três perguntas de modo explícito. O \textit{por que} decorre das condições materiais descritas acima e da coerência com o próprio referencial teórico (estudar existências parciais com uma existência parcial). O \textit{como} está documentado nesta seção e no apêndice ~\ref{apendice:colaboracao-ia}. O \textit{qual} atravessa os capítulos que se seguem. Minha colaboração com \texttt{claude.ai}é, antes de tudo, experimentação antropológica: uma etnógrafa que estuda como cientistas e engenheiros constroem inteligência artificial decide construir sua própria tese com inteligência artificial, submetendo-se às mesmas mediações, opacidades e negociações que observa em seus interlocutores. O gesto é epistêmico e produz conhecimento sobre a IA pela experiência de compor com ela, com todos os riscos que essa composição implica.

E se levarmos a sério o argumento de Yuk Hui, essa composição não é generalizável.\footnote{Vicentin conecta a proposta de Simondon de aprofundamento da tecnologia com a defesa de Yuk Hui de que, em lugar de axiomatizar a moral, precisaremos voltar aos modos de conhecimento diferentes que ainda não foram considerados por engenheiros e acadêmicos que trabalham com inteligência artificial \parencite[p.~186]{hui2020tecnodiversidade}. Essa convergência tem peso para meu argumento: se o aprofundamento da IA passa pela conexão entre o fazer técnico especializado e o conhecimento sócio-histórico sobre a tecnologia, então a colaboração com IA na escrita desta tese pode ser interpretada como experimentação situada nos termos em que Hui concebe a cosmotécnica.} Hui propõe o conceito de cosmotécnica para desafiar a ideia de que a tecnologia seria categoria antropológica única e homogênea: a técnica é sempre a unificação entre uma ordem cósmica e uma ordem moral através de atividades técnicas, e portanto varia ontológica e cosmologicamente entre diferentes sociedades, em lugar de variar apenas funcional ou esteticamente.\footnote{Hui desenvolve o conceito de cosmotécnica em \textit{The Question Concerning Technology in China} \parencite{hui2016china}, onde demonstra que a técnica pode orientar-se pela harmonia com uma cosmologia moral, sem buscar necessariamente a eficiência mecânica. Hui chama essa possibilidade de tecnodiversidade: contra o monotecnologismo moderno, o reconhecimento de que existem múltiplas cosmotécnicas que permitem diferentes formas de vida social, política e estética \parencite{hui2020tecnodiversidade}.}

A visão convencional, enraizada no pensamento ocidental, trata a tecnologia como conjunto de ferramentas universais e neutras destinadas à exteriorização da memória e superação dos limites orgânicos. A cosmotécnica recusa essa universalidade: diferentes cosmologias produzem diferentes modos de compor com objetos técnicos, e quem pretende moldar o universal submete o mundo à sua própria cosmovisão. Ora, se há cosmotécnicas situadas, em lugar de a tecnologia, então o modo como eu, antropóloga brasileira, formada na tradição dos estudos de ciência e tecnologia latino-americanos, componho com um modelo de linguagem treinado predominantemente em dados anglófonos para produzir uma tese em português sobre redes tecnocientíficas é, ele próprio, gesto cosmotécnico particular: forma específica de unificar ordem epistêmica e atividade técnica que não se repete em outros contextos, outras tradições, outras cosmologias. A experimentação que documento nestas páginas é, portanto, modo de usar IA na pesquisa acadêmica, situado, parcial e responsável por seus próprios efeitos.


\section{Compostagens: compostar-com e tornar-se-com}
\label{sec:compostagens}

\epigrafe{%
Somos composto, não pós-humanos; habitamos as humusidades, não as humanidades. Filosofica e materialmente, sou compostista, não pós-humanista. Os bichos (humanos e não humanos) estão em devir-com mútuo e em todos os registros de tempo e de coisas, em emaranhados simpoiéticos, em mundificações e desmundicações terrenas, ecológicas e evolutivas do desenvolvimento.}%
{Donna Haraway, \textit{Ficar com o problema: fazer parentes no Chthuluceno}, 2023, p.~191.}

A figuração da compostagem que desenvolvo nesta seção dialoga com a figuração têxtil que trabalhei na seção~\ref{sec:cortando-redes}, e as duas funcionam em chaves descritivas complementares do mesmo trabalho etnográfico. A figuração têxtil endereça o problema metodológico da delimitação da rede e nomeia o gesto de cortar como condição do gesto de costurar; o \textit{crazy quilt} oferece, nessa chave, o modelo da composição aparente de retalhos heterogêneos cujas junções são parte do que se descreve. A figuração da compostagem trabalha em outra chave descritiva. Endereça o que se passou entre meus materiais teóricos e empíricos ao longo da pesquisa, e nomeia processos de transformação recíproca em que matérias heterogêneas postas em camadas não permaneceram idênticas a si mesmas. Onde a costura justapõe retalhos cujas bordas continuam reconhecíveis após o gesto, a compostagem descreve materiais que se decompuseram mutuamente para gerar matéria nova e fértil. As duas figurações são solidárias e nomeiam dimensões distintas do mesmo movimento: a têxtil dá a ver a composição da descrição etnográfica enquanto composição, a compostagem dá a ver as transformações recíprocas que tornaram essa composição possível. As duas funcionam como modos complementares de descrever a etnografia da pesquisa em inteligência artificial, e escapam à descrição como duas etapas em sequência ou como produto e seu processo prévio.\footnote{A solidariedade entre as duas figurações é, ela mesma, gesto haraweano. Tanto o \textit{crazy quilt} quanto o composto orgânico recorrem a vocabulário historicamente associado ao trabalho doméstico feminino e historicamente desclassificado como não conhecimento, e ambos recusam a partição entre fazer e pensar que estrutura o vocabulário ocidental dominante sobre tecnologia. A escolha de duas figurações têxteis-orgânicas, em lugar de uma única, responde à exigência de descrever a etnografia da pesquisa em inteligência artificial em duas dimensões simultâneas: a composição aparente da descrição e as transformações recíprocas que a tornaram possível.}

Minha etnografia acabou produzindo também outro modo de conhecimento que figura como \textit{compostagem}.\footnote{Cheguei ao termo \textit{compostagem} depois de tentar sem sucesso outros nomes para o que essas misturas significaram em minha pesquisa. \textit{Lições} foi a primeira tentativa, e descartei-a porque o vocabulário da lição supõe quem ensina e quem aprende, conteúdos definidos a transmitir, e direção da matéria do mestre para o aprendiz. A pesquisa não se passou assim. \textit{Aprendizagens} carregava o mesmo problema. O que aconteceu foi outra coisa, mais próxima do que \textcite{Haraway2023Ficar} nomeia como compostagem: camadas de materiais heterogêneos postas a fermentar ao longo dos anos, em que nenhum dos elementos pôde ser identificado como o conteúdo a ser transmitido nem como a posição que ensinava. Quatro deslizamentos da palavra organizam meu gesto argumentativo nesta seção. O primeiro é o gesto de Haraway de propor o composto no lugar do pós-humano: \enquote{somos húmus, não Homo, não anthropos; somos compostagem, não pós-humanos} \parencite[p.~32]{Haraway2016Staying}. A formulação foi sugerida a Haraway por Rusten Hogness e a autora descreve seu encontro com ela como pular numa pilha vermicular: \enquote{I jumped into that wormy pile} \parencite[p.~32]{Haraway2016Staying}. A figura recusa o humano-como-Homo em favor do humano-como-húmus, território fértil em que matérias heterogêneas se encontram e se transformam umas às outras. O segundo é a sympoiese, palavra que Haraway adota a partir da tese de Beth Dempster sobre sistemas \enquote{coletivamente produtores} sem fronteiras espaciais ou temporais autodefinidas, e que ela traduz como \textit{making-with}, fazer-com, em contraste com a autopoiese: \enquote{Sympoiesis is a simple word; it means `making-with.' Nothing makes itself; nothing is really autopoietic or self-organizing} \parencite[p.~58]{Haraway2016Staying}. Nada se faz sozinho, e o fazer é sempre em companhia. O terceiro é o par compostar-decompor, formulado por Haraway como prática conjunta de \textit{critters} humanos e não humanos que se compostam e decompõem uns aos outros, em todo tempo e matéria. A compostagem é gesto de juntar e gesto de desfazer, e nesse duplo gesto se ganha a fertilidade do composto. O quarto é a posição que Haraway reivindica: \enquote{I am a compostist, not a posthumanist} \parencite[p.~101]{Haraway2016Staying}. Compostista nomeia prática filosófica e material em curso, em lugar de identidade estabilizada. Dois traços do vocabulário haraweano importam para o uso que faço aqui. O primeiro é a metáfora da pilha quente. Haraway descreve sua leitura de Strathern como \enquote{I compost my soul in this hot pile} \parencite[p.~34]{Haraway2016Staying}, formulação que a tradução brasileira gloseia como \enquote{compostagem intelectual} e que o original nomeia em chave mais visceral, da pilha em fermentação. O segundo é a descrição que Haraway faz do próprio livro nos seus agradecimentos: \enquote{Cooking over many years, the compost pile of colleagues, students, and friends who have made this book possible is promiscuous, layered, and hot} \parencite[p.~ix]{Haraway2016Staying}. Cozinhando ao longo dos anos, em camadas, promíscua e quente: nesse sentido convoco aqui a compostagem para nomear o que aconteceu na minha pesquisa. Cada compostagem desta seção inscreve um momento em que materiais inicialmente mantidos separados pelo desenho da pesquisa precisaram ser postos a fermentar juntos, e em que essa fermentação implicou também a decomposição das partições prévias. A figuração conversa, ainda, com o uso que Haraway faz do termo a propósito de Marilyn Strathern, trabalhado na seção~\ref{sec:cortando-redes} desta tese como modo de pensar com, e tem afinidade com o vocabulário que Haraway desenvolve em torno de \textit{thinking-with} como \textit{becoming-with} \parencite[p.~33]{Haraway2016Staying}: pensar com funciona como tornar-se com, em vez de funcionar como aprender com.} As compostagens são um tipo de contribuição analítica desta tese para a etnografia de práticas tecnocientíficas envolvendo inteligência artificial, em especial a pesquisa em IA feita pelos pesquisadores, com eles e com IA (feita por mim com inteligência artificial). Essas compostagens não estavam previstas no desenho da pesquisa, e emergiram do encontro entre meu campo e os referenciais teóricos que convoquei ao longo do percurso. Cada uma nomeia uma mistura específica de materiais heterogêneos que precisaram fermentar juntos para que a descrição se tornasse possível. A primeira composta a regra latouriana de estudar a ciência em construção com a dispersão característica da pesquisa em IA, na qual não há laboratório centralizado a observar. A segunda composta o vocabulário da simetria metodológica com a especificidade da inteligência artificial generativa, e nessa fermentação decompõe a partição humano/não humano em existências parciais. A terceira composta minha formação em ciências sociais com o letramento técnico que o campo me exigiu, fazendo dessa fermentação parte da própria etnografia. A quarta composta a separação prévia entre ciência e política com a evidência empírica de que a pesquisa em IA, no C4AI, é simultaneamente decisão epistêmica e captação de recursos.

\subsection{Compostando o método: acompanhar a pesquisa em IA enquanto ciência em construção}
\label{subsec:compostagem1}

A primeira regra metodológica que Latour propõe em \textit{Ciência em Ação} \parencite{Latour2011CienciaEmAcao}, reproduzida e discutida na seção~\ref{sec:regras-metodologicas} deste capítulo, foi a regra que orientou minha entrada no campo. Decidir investigar a pesquisa em IA como ciência em construção foi convergência entre intuição de pesquisa e regra metodológica que se revelou adequada à medida que o campo se desdobrava diante de mim. Quando cheguei ao C4AI em 2022, o Centro tinha menos de dois anos de existência. Projetos como o Spira estavam em pleno desenvolvimento, a parceria com a IBM ainda parecia sólida, e as definições sobre o que o Centro seria e produziria estavam em disputa.\footnote{O relatório científico do período agosto de 2021 a julho de 2022 inscreve projetos em cinco desafios: no NLP2, o Corpus Carolina (650 milhões de \textit{tokens}), \textit{datasets} CORAA (290 horas de áudio) e Porttinari Treebank (12.000 sentenças anotadas); no BLAB (Computação Bio-Inspirada e Aprendizado de Máquina, designação que na configuração consolidada após 2023 corresponde aos desafios OceanML e KEML descritos na Tabela~\ref{tab:projetos_c4ai}), arquitetura modular sobre a Amazônia Azul; em Saúde, segmentação de lesões de AVC e identificação de pneumonia por Covid-19; no AgriBio, geoMULTIMAPS para análise de insegurança alimentar; e no AI Humanities, alinhamento do Observatório de IA ao NIC.br e aplicativos de robótica social. Nenhum estava concluído; todos envolviam negociações ativas sobre escopo, método e recursos.}

O que eu encontrei foram pesquisadores negociando recursos, redefinindo escopos, improvisando infraestruturas, testando modelos que falhavam, justificando escolhas metodológicas para públicos diversos, em lugar de ciência pronta para estudar. Tratava-se de chegar à pesquisa em IA enquanto seus fatos e artefatos ainda não haviam se estabilizado, em lugar de abrir uma caixa-preta. Essa condição de incompletude foi decisiva, embora tenha se dado de forma distinta do que se poderia esperar de uma etnografia de laboratório. Não acompanhei em tempo real a realização de nenhuma das pesquisas do C4AI. A pandemia de Covid-19 dispersou os pesquisadores. O laboratório não era centralizado, e tampouco havia espaço físico que reunisse a pesquisa em IA, que acontecia em ambientes distribuídos e muitas vezes inacessíveis à observação direta. A ciência em construção chegou a mim, portanto, mediada pela narrativa dos próprios pesquisadores, e essa mediação se revelou analiticamente produtiva.

Em ambos os capítulos etnográficos, as redes que descrevo funcionam como compostagens entre o que meus interlocutores me narraram e o que eu, como etnógrafa, costurei convocando categorias da teoria ator-rede. O grau dessa fermentação, contudo, variou. No caso de Fábio e Cláudio, minha construção analítica foi particularmente incisiva: eles narraram o percurso institucional do C4AI e suas genealogias, e a cadeia que conecta Herman Hollerith em 1890 à Bluetalks de 2023, passando pela história da IBM e pela cultura dos \textit{mainframes}, foi construída por mim a partir de pistas que suas narrativas ofereceram, combinadas com pesquisa documental e histórica. O capítulo~\ref{capitulo3} é, portanto, o mais híbrido: parte da rede foi narrada pelos interlocutores, parte foi tecida por mim ao segui-los.

No caso de Marcelo, a compostagem se fez de modo distinto. Foi ele quem me contou, em entrevista, como construiu o Spira e o Corpus Carolina, e o fez de modo muito diferente do que aparece nos artigos científicos que publicou: contou-me a pesquisa tal como ela se fez, com suas contingências, improvisações, fracassos e sucessos, em lugar da versão estabilizada que os artigos publicam como triunfo. É por isso que o capítulo~\ref{capitulo4} se intitula \enquote{A rede que Marcelo construiu}: parti de sua narrativa para reconstruir, analiticamente, a rede de associações que sustentou aquelas pesquisas. Acessar a ciência antes de sua estabilização significou, nesse caso, alcançar o que a inscrição científica final apagou: as discussões sobre quais dados utilizar no treinamento, os impasses metodológicos com vieses identificados, as negociações entre abordagens computacionais distintas.

A primeira regra de Latour, nessas duas chaves, não exigiu presença no laboratório. Exigiu de mim recusar a versão estabilizada e rastrear as controvérsias, as escolhas descartadas e as contingências que a consolidação apaga.

O encerramento da parceria IBM-C4AI durante a finalização desta pesquisa conferiu a essa compostagem dimensão inesperada: abertura de uma caixa-preta institucional. O que parecia arranjo institucional estabilizado revelou-se, quando desfeito unilateralmente, fermentação precária de alianças cuja fragilidade só se deu a ver com a ruptura. A primeira regra de Latour trabalhou, portanto, em três chaves: no acesso à pesquisa em IA antes de sua consolidação, na construção analítica de associações que os próprios atores não haviam formulado, e na abertura involuntária de uma caixa-preta institucional que o próprio campo tratou de romper.

\subsection{Compostando a teoria: da simetria às existências parciais}
\label{subsec:compostagem2}

Esta compostagem inscreve um movimento que se desenrolou ao longo da minha pesquisa, e que a seção~\ref{sec:existencias-parciais} desenvolveu como argumento teórico. O ponto que recupero aqui é etnográfico: descrever como, na prática do trabalho de campo, o vocabulário da simetria foi aos poucos sendo decomposto em vocabulário de existências parciais. A decomposição não foi decisão tomada antes da entrada em campo. Resultou do que o campo me empurrou a reconhecer.

Ao iniciar minha pesquisa de campo no C4AI, eu não tinha conhecimento especializado em inteligência artificial, nem mesmo o letramento técnico mínimo que me permitiria, por exemplo, ler um artigo da área ou acompanhar uma discussão sobre arquiteturas de modelos. Essa condição produziu efeitos contraditórios, e seu primeiro efeito foi positivo: não ter familiaridade com o campo me permitiu experimentar a simetria metodológica proposta pela teoria ator-rede. Ao tratar a atividade dos pesquisadores como algo novo para mim, evitei a tendência de adotar antecipadamente categorias que separam o social do técnico. Tivesse eu chegado ao Centro convicta de que a IA é apenas estatística aplicada, ou de que os pesquisadores estavam a serviço de interesses corporativos, teria repetido críticas prontas em lugar de observar o que efetivamente acontecia. A curiosidade foi o que me permitiu a entrada no campo sem hierarquias prévias sobre o que constitui conhecimento sobre IA.

A suspensão de familiaridade só foi analiticamente produtiva porque o referencial teórico também não estava definido de antemão. Adotei a teoria ator-rede como quadro analítico em consequência do que encontrei no campo, em lugar de decisão tomada antes da entrada nele. Quando comecei a pesquisa, em 2022, eu não dispunha de estrutura teórica rígida; tinha leituras esparsas em estudos de ciência e tecnologia, curiosidade pelo campo e a orientação metodológica de \textcite{Kofes2015Maconaria}: o método se ensaia enquanto se experimenta. Essa ausência de enquadramento prévio se revelou vantajosa, pois meu percurso de pesquisa foi delineado conforme a própria rede tecnocientífica do Centro se configurava diante de mim, e minha escolha de autores e conceitos respondeu ao que as observações empíricas exigiam.

Dois aspectos do campo tornaram a teoria ator-rede particularmente adequada. O primeiro foi a dispersão. Em diferença dos laboratórios tradicionais, onde pesquisadores compartilham espaço físico, o Centro tinha estrutura distribuída: pesquisadores trabalhavam em cidades diferentes, conectavam-se por reuniões virtuais e compartilhavam infraestruturas computacionais remotas. Uma abordagem que pressupusesse o laboratório como unidade de observação teria fracassado. A teoria ator-rede, com sua ênfase em seguir os actantes através de suas associações, me oferecia flexibilidade para acompanhar práticas que não se fixavam em espaço único.

O segundo aspecto foi a heterogeneidade. À medida que seguia os pesquisadores, as redes que encontrava reuniam elementos que categorias disciplinares convencionais separariam: um mesmo percurso etnográfico podia passar por negociação orçamentária com a FAPESP, decisão sobre arquitetura de rede neural, disputa sobre propriedade intelectual com a IBM e conversa sobre dificuldades de gravar pacientes com Covid-19. Cada uma dessas passagens era heterogeneidade que só fazia sentido junto às demais, e escapava à descrição como mais social ou mais técnica. A teoria ator-rede me permitia descrevê-las sem separá-las previamente, e me exigia que eu não o fizesse.

Ator e rede ganham novo significado quando reunidos no hífen: apontam para ontologia relacional na qual actantes humanos e não humanos são simultaneamente pontos da rede e a própria rede em suas conexões \parencite{latour1996clarifications}. Ao aplicar essa abordagem à minha própria prática de pesquisa, reconheci que também me tornei ator-rede: ponto específico na rede do Centro e, ao mesmo tempo, constituinte dessa rede através das relações que estabeleci. Minha construção etnográfica da rede do Centro foi também construção de mim enquanto pesquisadora inserida nessa rede. Explicitar a posição da etnógrafa é, portanto, parte da análise.

Digo ponto de partida, porque o próprio percurso da minha pesquisa revelou os limites do vocabulário simétrico. A simetria contesta a hierarquia entre humano e não humano, e ainda assim pressupõe dois polos separados a serem tratados de modo equivalente. No campo, essa separação se mostrava insuficiente. Quando Marcelo descrevia o treinamento do Spira, o modelo escapava à descrição como ferramenta humana e como agente autônomo. Funcionava como compostagem instável de decisões dos pesquisadores e comportamentos que os surpreendiam. A grade simétrica pedia que eu tratasse o modelo e os pesquisadores nos mesmos termos; o que eu observava era que nenhum dos dois existia de modo completo antes da relação entre eles.

O vocabulário de Law sobre multiplicidade ontológica me oferece conceitos para pensar esse problema: o Spira figurava como objeto múltiplo sendo continuamente produzido através de práticas heterogêneas que não convergiam necessariamente, em lugar de figurar como objeto singular com duas perspectivas \parencite[p.~59]{Law2004After}. A suspensão de familiaridade que a simetria possibilitou foi condição necessária, e o campo me empurrou em direção ao que Latour formula como existências parciais: entidades que nunca foram completas para começar. Essa passagem, da simetria como regra metodológica à existência parcial como descrição do que encontrei, emergiu do encontro com um objeto que a exigia. O vocabulário das existências parciais só se tornou formulável porque minha abertura teórica inicial permitiu que o campo falasse antes que eu decidisse como ouvi-lo. A elaboração teórica desse vocabulário, e em especial a leitura interna da obra de Latour que ele autoriza, está desenvolvida na seção~\ref{sec:existencias-parciais} deste capítulo, e nesta compostagem ela aparece como achado empírico mais do que como construção conceitual.

A suspensão da familiaridade me trouxe também vulnerabilidade. Foi precisamente ao seguir os actantes através de suas associações heterogêneas que percebi a insuficiência de meu próprio letramento: a teoria ator-rede me exigia que eu acompanhasse os pesquisadores onde quer que suas redes os levassem, inclusive para dentro de discussões técnicas que eu não compreendia. À medida que a pesquisa avançava, minha falta de conhecimento técnico tornava-se obstáculo para compreender o que os pesquisadores efetivamente faziam. Sem esse letramento, minha observação etnográfica permaneceria na superfície das práticas. O referencial que o campo me ofereceu era também o referencial que me obrigava a transformar-me para poder usá-lo. Como e por que busquei esse letramento, e o que o processo revelou sobre relações entre campos de conhecimento, é o tema da compostagem~\ref{subsec:compostagem3}.

\subsection{Compostando o letramento: fermentar com os cientistas}
\label{subsec:compostagem3}

Se a compostagem~\ref{subsec:compostagem2} tratou do que minha ignorância técnica permitiu, esta trata do que ela impediu e de como precisei superá-la. O ponto de partida foi mais amplo do que a dificuldade de compreensão. Em uma ocasião, cientistas da computação me interpelaram sobre generalizações que eu reproduzia a partir de autores de fora do campo: estudiosos que escreviam sobre inteligência artificial sem, na avaliação de meus interlocutores, compreendê-la adequadamente. A interpelação tinha fundamento: eu mesma, sem o letramento técnico que me permitiria avaliar o alcance dessas afirmações, as incorporava a meu repertório. O episódio revelou, porém, um tensionamento que vai além de minha situação individual.

Pesquisadores de IA demonstravam desconfiança quando cientistas sociais se pronunciavam sobre inteligência artificial, exigindo deles domínio técnico como condição de legitimidade. Quando eles próprios tomavam o social como objeto de suas pesquisas, desenvolvendo sistemas que classificam emoções humanas, detectam discurso de ódio ou modelam comportamento de populações, a questão da autoridade disciplinar raramente se colocava. A competência das ciências sociais sobre seus próprios objetos parecia escapar do mesmo respeito que a competência computacional reivindicava para si. Inscrevo essa tensão como observação etnográfica que implica ambos os lados, em lugar de denúncia. As ciências sociais também atravessam fronteiras disciplinares quando se pronunciam sobre tecnologias que não dominam tecnicamente.

Stengers me oferece instrumentos para compreender esse tensionamento. A distinção entre ciência dura e mole diz menos sobre graus de rigor e mais sobre modos distintos de relação com o objeto de estudo \parencite{stengers2000invention}. As ciências experimentais constroem sua autoridade através do laboratório, em que o objeto é purificado e tornado indiferente às intenções do pesquisador: um elétron não se importa com o que o físico pensa dele. As ciências sociais, por sua vez, lidam com seres não indiferentes às perguntas que lhes são feitas: o objeto olha de volta, interpreta o aparato de pesquisa, questiona as intenções do pesquisador. Uma ciência é chamada de mole quando não especialistas se sentem competentes para opinar sobre ela, precisamente porque suas perguntas dizem respeito à vida das pessoas. Essa é descrição aproximada do que observei: os pesquisadores de IA sentiam-se autorizados a opinar sobre comportamento humano, emoções e linguagem porque esses temas parecem acessíveis, e consideravam que falar sobre inteligência artificial exigia credenciais que os cientistas sociais não possuíam.

A maioria dos cientistas sociais não possui conhecimento técnico sobre inteligência artificial, o que não é, em si, um problema. O letramento técnico torna-se necessário quando a IA ou a pesquisa em IA é o próprio objeto de investigação, como foi meu caso. O episódio no C4AI aponta para algo que vai além de minha situação particular. Quando o letramento técnico funciona como única via de legitimação para participar de discussões sobre IA, o efeito prático é que a responsabilidade de compreender os sistemas recai sobre quem está de fora da computação. Esse efeito não decorre de determinação explícita, e sim de uma ausência: na ausência de espaços institucionais de colaboração, quem não domina a técnica fica sem vocabulário para formular seus problemas em termos que os pesquisadores de IA reconheçam. Os riscos gerados por esses sistemas acabam sendo enfrentados por quem os sofre, em vez de serem enfrentados por quem os constrói.

Essa assimetria se dava também nos modos concretos de colaboração interdisciplinar no Centro. No projeto Spira, linguistas, músicos computacionais, médicos e enfermeiros participaram ativamente da construção do sistema, cada qual contribuindo com competências que os pesquisadores de IA reconheciam como indispensáveis. São raras, contudo, as ocasiões em que cientistas sociais são convocados. Quando o são, a motivação tende a ser de outra ordem: situações eticamente delicadas, como projetos que envolvem comunidades tradicionais, nas quais a presença de um especialista funciona menos como interlocução intelectual e mais como salvaguarda institucional. Fora desses casos, pesquisas que afetavam diretamente dimensões sociais eram conduzidas sem essa interlocução. A interdisciplinaridade proclamada se dava de modo seletivo: certas competências eram incorporadas como constitutivas da pesquisa, e outras eram convocadas apenas quando o risco de não fazê-lo se tornava visível.

Há nisso uma ironia que merece atenção. Como documentei na seção sobre existências parciais, recorri a modelos de linguagem para me auxiliar na análise de dados e na elaboração de \textit{scripts}, trabalho diretamente relacionado às competências dos cientistas da computação. A pergunta simétrica se impõe: quando pesquisadores de IA constroem sistemas que modelam comportamento social ou processam linguagem natural de comunidades específicas, recorrem eles a ferramentas conceituais das ciências sociais para formular seus problemas de pesquisa? A resposta, pelo que observei no Centro, é que raramente o fazem. A interdisciplinaridade tende a se dar em mão única: cientistas sociais são incentivados a adquirir letramento técnico, e o caminho inverso permanece excepcional.

Latour me permite compreender por que essa assimetria se mantém. Quanto mais dura uma ciência se torna, mais social ela é, porque sua dureza depende de número desproporcional de associações, instrumentos, laboratórios e aliados alistados para tornar seus fatos inquestionáveis \parencite{Latour2011CienciaEmAcao}. A dureza é efeito do comprimento da rede que sustenta esses objetos. Os pesquisadores que eu observava no C4AI estavam, justamente, no trabalho de construir essa dureza: alistando \textit{datasets}, GPUs, publicações, métricas, financiamentos e parcerias para estabilizar afirmações que, sem essa rede, permaneceriam abertas à renegociação.

Stengers e Latour me iluminam, portanto, gestos complementares. Stengers diagnostica a imposição de um modelo de objetividade como critério universal. Latour revela o trabalho material que sustenta essa imposição: a extensão da rede como fonte da autoridade. Juntos, me ajudam a compreender por que a interdisciplinaridade no C4AI se dava de modo seletivo e por que o caminho do letramento técnico era percorrido em mão única.

Essa assimetria me levou a buscar o letramento técnico como exigência etnográfica, em lugar de gesto de contestação. Ao fazer pesquisa junto ao Centro, compreendi que para entender adequadamente um artigo científico da área seria necessário ir substancialmente além das generalizações frequentemente reproduzidas em debates externos ao campo. A linguagem técnica é desafio epistemológico de outra ordem pela sua tecnicidade intrínseca e também porque, devido à quantidade de elos, recursos e aliados que aciona, ela se torna paradoxalmente mais social do que vínculos comunicativos convencionais \parencite{Latour2011CienciaEmAcao}. Cada artigo, cada equação, cada escolha metodológica carrega consigo extensa cadeia de referências, pressupostos compartilhados, debates anteriores, infraestruturas computacionais. Essas cadeias precisam ser minimamente conhecidas para que a compreensão seja efetiva.

Para enfrentar esses desafios, precisei investir tempo assistindo a palestras técnicas e fazendo cursos sobre fundamentos da estatística e do aprendizado de máquina, entre os quais o curso \textit{Fundamentos de Probabilidade e Estatística para Ciência de Dados}, ministrado pelo Prof. Dr. Francisco A. Rodrigues (ICMC-USP) no âmbito do MBA em Ciência de Dados da USP.\footnote{O curso cobre três blocos temáticos encadeados: fundamentos de probabilidade (teorema de Bayes, esperança, momentos, teorema central do limite, lei dos grandes números), inferência estatística (estimação, intervalo de confiança, teste de hipóteses, inferência bayesiana) e modelagem com aprendizado supervisionado (regressão, ajuste de modelos, classificação). O material está disponível em: \url{https://cursosextensao.usp.br/course/view.php?id=4347}. Acesso em: mar. 2026. Esse percurso formativo é retomado no capítulo~\ref{capitulo4}, onde os conceitos de aprendizado supervisionado e classificação binária me servem para descrever o \textit{pipeline} técnico do Spira.} Foi assim que compreendi, por exemplo, que o laboratório do cientista da computação é o computador, ideia aparentemente óbvia que tem implicações que merecem atenção.\footnote{A formulação é do Prof. Dr. Francisco A. Rodrigues, enunciada durante o curso \textit{Fundamentos de Probabilidade e Estatística para Ciência de Dados} (ICMC-USP), no contexto da introdução à simulação computacional: \enquote{o computador é o laboratório para aprender os conceitos estatísticos e probabilísticos} \parencite{Rodrigues2024ProbEstat}. Anotei a frase em meu caderno de campo durante a aula (Figura~\ref{fig:caderno-rodrigues}). O que parecia observação didática sobre o ensino de estatística se revelou, ao longo de minha pesquisa de campo no C4AI, enunciado com consequências analíticas: se o computador é o laboratório, então o laboratório do cientista da computação é portátil, distribuível e dispensador de presença física contínua, o que reformula inteiramente a questão que abria esta pesquisa, a saber, onde está o laboratório de IA e onde estão os seus cientistas.}

\begin{figure}[!htb]
  \centering
  \includegraphics[width=0.75\textwidth]{figuras/cap.1/caderno_quadro.jpg}
  \caption[Anotação de caderno de campo de 28 de julho de 2025]{Anotação do meu caderno de campo de 28 de julho de 2025 durante o curso \textit{Fundamentos de Probabilidade e Estatística para Ciência de Dados} (ICMC-USP, Prof. Dr. Francisco A. Rodrigues). O quadro destaca a frase: \enquote{Simulação: `o computador é o laboratório para aprender os conceitos estatísticos e probabilísticos'.} As palavras \enquote{computador} e \enquote{laboratório} estão circuladas.}
  \label{fig:caderno-rodrigues}
\end{figure}

O nível de letramento técnico necessário varia conforme o objeto de estudo: investigações focadas nas práticas de pesquisa em inteligência artificial, como foi meu caso, exigem conhecimento aprofundado. Estudos centrados nos impactos sociais de sistemas já consolidados podem demandar níveis diferentes de familiarização. A necessidade de letramento técnico que experimentei encontra eco em debates dentro da própria ciência da computação. \textcite{Wallach2014} formula pergunta que espelha, desde dentro do campo, a observação que faço desde fora: poucos cientistas da computação desenvolveriam modelos para analisar dados astronômicos sem envolver astrônomos; por que, então, tantos métodos para analisar dados sociais são desenvolvidos sem o envolvimento de cientistas sociais? A resposta de Wallach é que, como seres humanos, temos intuições fortes sobre o mundo social, e essas intuições nos fazem crer que compreendemos fenômenos sociais sem precisar de formação específica. Intuição, porém, escapa à descrição como análise. O argumento de Wallach é estrutural e específico: a colaboração precisa ser séria, no modo crie equipes genuinamente interdisciplinares, em lugar do modo contrate um cientista social para cada cem cientistas da computação. A distinção importa: no primeiro caso, a presença do cientista social funciona como decoração; no segundo, é constitutiva da pesquisa.

A consequência dessa formulação é circunscrita. Não significa que todo projeto de IA deva incluir cientistas sociais. Significa que a decisão sobre quando essa interlocução faz diferença não deveria caber exclusivamente a quem não a pratica. É essa a questão que percorreu toda esta compostagem: quando pesquisadores de IA tomam o social como objeto, estão fazendo ciência social computacional sem reconhecê-la como tal, e portanto sem as salvaguardas que a tradição das ciências sociais desenvolveu para lidar com seres que não são indiferentes às perguntas que lhes são feitas. Reconhecer esse padrão é condição para construir alianças genuínas: alianças que permitam às ciências sociais participar da construção dos problemas de pesquisa, em lugar de participar apenas da gestão de seus efeitos colaterais.

O diagnóstico de Wallach converge, por via inesperada, com o aprofundamento da tecnologia que Vicentin formula a partir de Simondon: sem a conexão entre o fazer técnico e o conhecimento sócio-histórico, a tecnologia permanece aquém de sua ética imanente \parencite{Vicentin_2022_Esboco}. Quando os dispositivos de classificação e predição da IA são utilizados para reforçar estruturas de poder baseadas em classe, raça e gênero, a IA permanece em estágio de imaturidade no sentido simondoniano: limitada a reproduzir estruturas existentes em lugar de realizar seu potencial transformador.

O argumento desta compostagem pode agora ser reunido. A interpelação que recebi no C4AI revelou um padrão que se dá em três planos. Pesquisadores de IA exigem letramento técnico de cientistas sociais que falam sobre IA, e dispensam letramento sociológico quando eles próprios tratam do social. A colaboração interdisciplinar é seletiva: certas competências são incorporadas como constitutivas da pesquisa, e outras convocadas apenas como salvaguarda institucional. O letramento se faz em mão única, do social para o técnico. Minha resposta foi dupla: adquirir o letramento necessário para compreender meus interlocutores nos seus próprios termos, e documentar a assimetria como achado etnográfico. As duas coisas não se contradizem. Aprender a linguagem técnica foi condição para habitar o campo com alguma autonomia e, ao mesmo tempo, descrevê-lo com precisão, em lugar de capitular diante da hierarquia disciplinar. A compostagem~\ref{subsec:compostagem2} documentou a suspensão inicial que abriu o campo; esta documenta a aquisição progressiva de letramento que me permitiu habitá-lo.

\subsection{Compostando a política: ciência como captação de recursos}
\label{subsec:compostagem4}

De todas as compostagens, esta foi talvez a mais imediata, embora tenha sido a última a se formular como tal. Preciso ser franca sobre algo que me acompanhou durante boa parte do trabalho de campo: irritação surda, difícil de formular na época, com o fato de que eu não via a inteligência artificial sendo feita, ao menos como eu imaginava. Eu havia me proposto a fazer uma etnografia de um centro de pesquisa em IA, e os leitores de uma tese como esta esperariam, com razão, saber como a inteligência artificial é produzida. O que eu encontrava, contudo, era outra coisa: eventos de difusão científica, negociações de financiamento, apresentações a parceiros corporativos, disputas por infraestrutura computacional. A pesquisa em IA parecia acontecer em outro lugar, sempre atrás de uma porta que eu não conseguia abrir, dentro de computadores aos quais eu não tinha acesso, em reuniões técnicas cujo vocabulário eu não dominava. O que eu tinha diante de mim era, como defendo nas compostagens anteriores, a própria construção da inteligência artificial. Eu apenas não a identificava dessa maneira, porque seguia em busca da ciência por trás de toda aquela política.

Assistia a Fábio apresentar o Centro em eventos de difusão e anotava \enquote{divulgação}, em lugar de \enquote{ciência}. Ouvia Cláudio explicar como navegava entre os três potes de recursos e classificava aquilo como \enquote{gestão}, em lugar de \enquote{pesquisa}. Quando Cláudio me descreveu em entrevista como, um mês antes de \textit{deadlines} de conferências, tornava-se impossível conseguir GPU no Centro, reconheci ali o mesmo padrão: eu teria classificado aquilo como \enquote{logística}, em lugar de decisão epistemológica. Na minha primeira visita ao Centro, parei diante da sala de TI sem compreender o que acontecia ali. Anos depois, percebi que aquele desconforto era sintomático de uma separação que eu carregava e que o campo insistia em desfazer. Eu procurava a inteligência artificial dentro dos computadores, nos códigos, nos modelos, e o que encontrava eram pessoas negociando, traduzindo interesses, reconfigurando alianças. A angústia era intelectual e profissional, e me perguntava com frequência como apresentar uma etnografia de um centro de inteligência artificial sem propriamente descrever como a IA é feita.

Essa frustração, contudo, era metodologicamente produtiva. Ela me sinalizava o ponto exato onde minhas categorias analíticas falhavam, onde o arranjo que eu trazia precisava ser reconfigurado. Law argumenta que o \textit{hinterland} de qualquer método se ramifica indefinidamente para além de laboratórios e protocolos, alcançando conhecimento tácito, \textit{software}, habilidades linguísticas, capacidades de gestão, sistemas de transporte e comunicação, escalas salariais, fluxos de financiamento, prioridades de agências de fomento e agendas abertamente políticas e econômicas \parencite[pp.~41--42]{Law2004After}. O que eu observava no C4AI (as sessões de \textit{Design Thinking} que costuravam critérios da USP, IBM e Brasil; a venda da divisão de saúde da IBM que redirecionava linhas de pesquisa; as disputas por GPU que se tornavam decisões epistemológicas; Fábio apresentando em Bluetalks como parte do fazer ciência) era precisamente esse \textit{hinterland} ramificado que métodos convencionais relegam à invisibilidade.

Law descreve \textit{method assemblage} como produção de feixe de relações ramificadas que gera presença, ausência manifesta e \textit{otherness} \parencite[p.~84]{Law2004After}. No meu caso: a presença esperada era a ciência sendo feita (códigos, modelos, experimentos); a ausência manifesta eram as negociações que eu sabia serem necessárias e classificava como \enquote{entorno}; e a \textit{otherness} era precisamente o trabalho de captação de recursos que precisava desaparecer para que a ciência aparecesse como autônoma. Minha frustração era sintoma de que minhas categorias prévias produziam exatamente essa divisão entre presença, ausência e \textit{otherness}. O que o campo me obrigou a fazer foi reconfigurar essas fronteiras: reconhecer que o que eu havia relegado à \textit{otherness} era constitutivo da presença.

A pergunta \enquote{onde está a ciência que esses recursos sustentam?} pressupunha a separação entre ciência e política que o próprio campo tratava de desmentir. Captar recursos, negociar parcerias, justificar investimentos perante agências de fomento e corporações transnacionais é a pesquisa ela mesma. A pergunta não encontrou resposta; perdeu sentido.

A dificuldade em reconhecer aquelas atividades como fazer IA era expressão concreta da alienação técnica discutida anteriormente: a separação entre o que seria estritamente técnico e o que seria entorno da pesquisa reproduz a cisão que impede o amadurecimento da tecnologia. O que o campo me obrigou a compreender, por vias etnográficas, foi que as negociações de financiamento, as disputas por GPU, as apresentações a parceiros corporativos eram parte constitutiva da individuação técnica da IA.

Essa inseparabilidade entre ciência e política era o que o campo me dava a ver quando eu conseguia suspender a distinção prévia. Os próprios projetos de pesquisa do Centro foram definidos, como Fábio me relata, por sessões de \textit{Design Thinking} em que três critérios precisavam se alinhar simultaneamente: a USP teria liderança científica no tema, a IBM teria interesse na temática, e o Brasil teria diferencial internacional. A decisão era uma só, e nela se fundiam avaliação epistêmica, estratégia corporativa e posicionamento geopolítico, sem que houvesse momento científico de escolha de tema seguido de momento político de negociação de recursos.\footnote{Os relatórios internos do C4AI documentam esse processo: o grupo de \textit{Design Thinking} organizou quatro reuniões com mais de 80 pesquisadores, convergindo em cinco desafios. O critério exigia que os resultados fossem estratégicos para o Brasil, relevantes para USP e IBM e com alto potencial de inovação. A partir de 2022, quando a IBM vendeu Watson Health, o Centro decidiu não enfatizar esforços relacionados a saúde; quando o grupo AgriBio conquistou financiamento independente (INCT), tornou-se autossuficiente em relação aos recursos do C4AI. Em ambos os casos, a reconfiguração não separava decisão científica de decisão política.}

Quando a IBM vendeu sua divisão de saúde, o Centro parou de investir diretamente em pesquisa de IA em saúde. Quando a IBM saiu de agronegócio, o mesmo ocorreu. As reorientações de pesquisa se davam como reconfigurações da rede em que agenda científica e interesse corporativo constituíam, desde a origem, um único arranjo, em lugar de imposições externas sobre uma agenda científica autônoma.

Do mesmo modo, quando Fábio apresentava resultados em eventos como as Bluetalks, estava recrutando aliados, fortalecendo a rede, traduzindo interesses de públicos heterogêneos para manter a cadeia de associações que tornava a pesquisa possível. Latour descreve precisamente esse trabalho como constitutivo da prática tecnocientífica: o que se perde ao separá-lo da ciência propriamente dita é justamente o que o torna ciência, ou seja, a capacidade de alistar e manter associados elementos heterogêneos \parencite{Latour2011CienciaEmAcao}. A junção dos termos técnica e ciência em \textit{tecnociência} é exigência empírica, e o campo que eu observava tornava essa exigência particularmente visível.

A agência, nessas reconfigurações, é distribuída pela rede: emerge das associações entre pesquisadores, instituições, infraestruturas, contratos, modelos computacionais e políticas de fomento. Nenhum desses elementos age sozinho ou determina sozinho o curso da pesquisa. Quando o C4AI delega a definição de temas de pesquisa a sessões de \textit{Design Thinking} que costuram simultaneamente critérios da USP, da IBM e do Brasil, a agência dessa decisão reside na associação entre os participantes, e cada reconfiguração da rede redistribui a agência por novas cadeias de translação.

O cenário contemporâneo de investimentos massivos em inteligência artificial (com a concentração de capacidade computacional em poucas corporações, a corrida geopolítica por soberania tecnológica e a crescente dependência de infraestruturas de alto custo) torna essas dinâmicas particularmente visíveis. Resta saber se a pesquisa em IA é caso em que a inseparabilidade entre ciência e política se intensifica a ponto de exigir novos instrumentos analíticos, ou se apenas torna mais difícil ignorar o que Latour já descrevia em outros campos, quando os investimentos eram menores e as redes, menos extensas.


\section{Conclusões parciais: o capítulo como composição}
\label{sec:conclusoes-cap1}

Este capítulo funcionou como laboratório de suas próprias decisões. Começou pela pergunta sobre onde está o laboratório dos pesquisadores de IA do C4AI, e a resposta me exigiu um conjunto de movimentos que o próprio capítulo inscreveu enquanto os realizava: deslocar minha expectativa do laboratório centralizado para o laboratório disperso e portátil, reconhecer a inteligência artificial como ciência em construção em lugar de objeto consolidado, e aceitar que a etnografia de uma rede tecnocientífica desse tipo se faz por composição de materiais heterogêneos, em lugar de descrição linear de um lugar único.

Três gestos analíticos se desdobraram a partir dessa pergunta inicial. O primeiro foi nomear cortes, ou seja, os gestos de delimitação que tornam descritível uma rede que se desdobra indefinidamente. A subseção \enquote{Cortando redes} convocou cinco vocabulários (Strathern, Barad, Haraway, Costa e Law) para descrever como cada capítulo desta tese faz um corte específico: reflexivo no capítulo~\ref{capitulo1}, teórico no capítulo~\ref{capitulo2}, pela extensão no capítulo~\ref{capitulo3} e pela densidade no capítulo~\ref{capitulo4}. Os quatro cortes funcionam como \textit{enactments} no sentido praxiográfico que Mol e Law desenvolvem, e os quatro capítulos performam o C4AI como quatro objetos etnográficos que se atravessam sem se subsumir uns aos outros.

O segundo gesto foi descrever a composição com inteligência artificial generativa que minha tese assume como dado etnográfico, em lugar de recurso técnico secundário. A seção \enquote{Existências parciais} desdobrou um movimento teórico interno à obra de Latour, do princípio de simetria como ponto de partida à formulação das existências parciais como descrição mais precisa do que o campo me apresentava. Esse movimento gerou três contribuições. A noção de quase-máquina nomeou os sistemas de IA generativa como entidades abertas e parcialmente conectadas, que escapam à partição entre humano e não humano. A noção de agência distribuída descreveu o mecanismo pelo qual essas existências parciais agem juntas, redistribuindo a responsabilidade pela ação por cadeias de translação. A multiplicidade ontológica de Law me permitiu reconhecer \texttt{claude.ai}como objeto múltiplo (em treinamento, em \textit{deployment}, em auditoria, em uso), e cada \texttt{claude.ai}existe como composição situada com cada pesquisadora que o aciona, em lugar de objeto singular acessível a todos.

O terceiro gesto foi figurar as transformações recíprocas que tornaram minha pesquisa possível. A seção \enquote{Compostagens} nomeou quatro misturas específicas de materiais heterogêneos que precisaram fermentar juntos para que a descrição se tornasse possível: a regra latouriana de estudar a ciência em construção fermentando com a dispersão característica da pesquisa em IA; a simetria metodológica fermentando com a especificidade da inteligência artificial generativa para gerar existências parciais; minha formação em ciências sociais fermentando com o letramento técnico que o campo me exigiu; e a separação prévia entre ciência e política fermentando com a evidência empírica de que a pesquisa em IA, no C4AI, é simultaneamente decisão epistêmica e captação de recursos. Cada compostagem inscreve um momento em que a descrição etnográfica exigiu decompor partições prévias e tolerar a fermentação de materiais que o desenho inicial da pesquisa mantinha separados.

O achado metodológico que sustenta este capítulo é, portanto, etnográfico antes de ser teórico. O método se ensaiou enquanto se experimentava, conforme a orientação de \textcite{Kofes2015Maconaria}, e os vocabulários que convoquei foram chamados pelo que o campo exigia. Esse percurso me autoriza um gesto que executei de modo recursivo neste capítulo~\ref{capitulo1}: aplicar à própria escrita da tese o tratamento que aplico aos pesquisadores que descrevo. A subseção \enquote{Tecnoetnografia} converteu o texto deste capítulo em grafo de co-ocorrências e leu suas comunidades como tópicos latentes, inscrevendo as costuras conseguidas e as fissuras pendentes. A subseção \enquote{Por que fazer parentesco com este problema?} documentou em sequência de capturas de tela uma cena de composição com \texttt{claude.ai}durante a revisão deste mesmo capítulo. Em ambos os casos, minha tese funciona como inscrição tecnocientífica que se deixa medir pelos mesmos instrumentos com que se compromete a medir o C4AI, e a reflexividade metodológica assume a forma de gesto recursivo, em lugar de declaração externa sobre a relação entre pesquisadora e objeto.

O que este capítulo deixou em aberto é o trabalho de localizar os aliados conceituais convocados por essa descrição em um campo de pesquisa em formação. Strathern, Barad, Haraway, Costa, Law, Latour, Mol, Kofes, Stengers, Vicentin, Simondon e Hui apareceram aqui como vocabulários chamados pelo campo, com as relações entre eles indicadas apenas no que cada gesto etnográfico exigia. O capítulo~\ref{capitulo2} retoma esse conjunto e o organiza como rede de aliados teóricos que sustentam a tese, conectando os estudos sociais da ciência e da tecnologia, as pesquisas pioneiras em antropologia e sociologia da inteligência artificial (Forsythe, Collins), as genealogias críticas do poder tecnológico (Crawford, Pasquinelli, Brennan) e o campo brasileiro em formação que minha tese habita. Ao recrutar esses aliados, o capítulo~\ref{capitulo2} faz o que \textcite{Latour2011CienciaEmAcao} descreve como trabalho de alistamento, e ao mesmo tempo inscreve, com \textcite{stengers2010cosmopolitics1}, que outros modos de fazer alianças com a ciência e com a tecnologia permanecem disponíveis para quem quiser tomá-los. Esse duplo gesto, alistar com Latour e propor com Stengers, configura o trabalho teórico do capítulo que se segue, e desdobra o problema metodológico aqui inscrito para o terreno da revisão crítica da literatura.

\begin{figure}[!htb]
  \centering
  \includegraphics[width=0.8\textwidth]{figuras/cap.1/trama-capitulo1-2026-05-12-225918.pdf}
  \caption[Trama do capítulo~\ref{capitulo1}]{\textbf{Trama do capítulo~\ref{capitulo1}: fios cruzados de cortes, existências parciais e compostagens.} A inscrição cruza os quatro grandes movimentos do capítulo. (i)~Na urdidura, três fundamentos etnográficos sustentam todo o restante: \textcite{Kofes2015Maconaria} (o método se ensaia enquanto se experimenta), \textcite{Latour2011CienciaEmAcao} (seguir os actantes) e \textcite{Law2004After} (\textit{method assemblage} como produção de presença, ausência manifesta e \textit{otherness}). (ii)~Em \textit{Cortando redes}, cinco vocabulários do corte (Strathern, Barad, Haraway, Costa e Law) alimentam os quatro cortes que cada capítulo da tese faz, cada um descrito pela tríade presença/ausência manifesta/\textit{otherness}, e os quatro cortes leem-se como \textit{enactments} no sentido praxiográfico de \textcite{mol2002body} e \textcite{Law2004After}. (iii)~Em \textit{Existências parciais}, um deslocamento interno na obra de Latour (do princípio de simetria aos híbridos e às existências parciais) sustenta três figurações teóricas (quase-máquinas, agência distribuída, multiplicidade ontológica de Claude), que por sua vez justificam o trabalho dos instrumentos de inscrição (\texttt{claude.ai}, \texttt{\texttt{\texttt{MermaidChart}}}, \texttt{Python}, \texttt{NotebookLM}, \texttt{Read AI}, \texttt{LaTeX}, Library Genesis, cadernos de campo); a subseção \textit{Por que fazer parentesco com este problema?} (Haraway) traz a composição com \texttt{claude.ai}como dado etnográfico, materializada na sequência de capturas de tela das Figuras~\ref{fig:composicao-claude-1}~a~\ref{fig:composicao-claude-3}; a subseção \textit{Tecnoetnografia} traz a análise no estilo \texttt{InfraNodus} do próprio capítulo (Figura~\ref{fig:infranodus-cap1}) como gesto recursivo, e a subseção \textit{Cosmotécnica} situa a composição entre antropóloga brasileira e modelo treinado predominantemente em dados anglófonos como gesto cosmotécnico particular (Hui). (iv)~Em \textit{Compostagens}, quatro fermentações específicas (do método, da teoria, do letramento e da política) descrevem as transformações recíprocas de materiais heterogêneos que tornaram a descrição etnográfica possível; cada compostagem aparece no diagrama como soma $A + B = C$, em que duas matérias inicialmente mantidas separadas pelo desenho da pesquisa precisaram ser postas a fermentar juntas. (v)~As \textit{Conclusões parciais} reúnem os três gestos analíticos do capítulo, registram a reflexividade como gesto recursivo e abrem a ponte para o capítulo~\ref{capitulo2}. A inscrição é, ela mesma, parte do que descreve: foi produzida em colaboração com \texttt{claude.ai}(Anthropic) e hospedada no \texttt{\texttt{MermaidChart}}\footnote{Versão interativa do diagrama disponível em: \url{https://mermaid.ai/d/d9085fca-258e-4f15-be1c-ee117981eb7f}. Acesso em: \today.}, em cadeia de traduções entre minha descrição verbal das relações argumentativas do capítulo e o código \textit{Mermaid} executado pelo modelo. A inscrição opera, nos termos de \textcite{Latour1986Visualization}, como móvel imutável que torna combinável em superfície plana o que o texto desdobra ao longo de dezenas de páginas, e nos termos de \textcite{Law2004After}, como redução produtiva que perde riqueza empírica e ganha capacidade de fazer conexões.}
  \label{fig:trama-cap1}
\end{figure}
